\documentclass[11pt,oneside]{article}

\usepackage[a4paper,margin=27mm,headheight=15pt]{geometry}
\usepackage[T1]{fontenc}
\usepackage[utf8]{inputenc}
\IfFileExists{libertinus.sty}{\usepackage{libertinus}}{\usepackage{lmodern}}
\usepackage{microtype}
\usepackage{amsmath,amssymb,amsthm,mathtools,mathrsfs}
\usepackage{bm}
\usepackage{booktabs,longtable,array,tabularx}
\usepackage{graphicx}
\usepackage{pgfplots}
\usepackage{xcolor}
\usepackage{enumitem}
\usepackage{fancyhdr}
\usepackage{titlesec}
\usepackage{listings}
\usepackage{xurl}
\usepackage{hyperref}
\usepackage[nameinlink,noabbrev]{cleveref}
\usepackage{multicol}
\IfFileExists{needspace.sty}{\usepackage{needspace}}{\newcommand{\Needspace}[1]{}}
\IfFileExists{bookmark.sty}{\usepackage{bookmark}}{}

\definecolor{linkblue}{RGB}{25,75,135}
\definecolor{shadegray}{RGB}{246,247,249}
\definecolor{rulegray}{RGB}{195,201,208}
\pgfplotsset{compat=1.18}
\hypersetup{
  colorlinks=true,
  linkcolor=linkblue,
  citecolor=linkblue,
  urlcolor=linkblue,
  pdftitle={A Mathematical Glossary for the Fabius-Function Development},
  pdfauthor={OpenAI, based on Vladimir Reshetnikov's ProveIt repository},
  pdfsubject={Encyclopedic glossary for Analysis/FabiusFunction},
  pdfkeywords={Fabius function, Rvachev up-function, glossary, dyadic rational, q-binomial, Lambert W, saddle point, Lean}
}

\pagestyle{fancy}
\fancyhf{}
\fancyhead[L]{\small A mathematical glossary for the Fabius-function development}
\fancyhead[R]{\small\nouppercase{\leftmark}}
\fancyfoot[C]{\thepage}
\renewcommand{\headrulewidth}{0.4pt}

\titleformat{\section}{\Large\bfseries}{\thesection}{0.7em}{}
\titleformat{\subsection}{\large\bfseries}{\thesubsection}{0.7em}{}
\titleformat{\subsubsection}{\normalsize\bfseries}{\thesubsubsection}{0.7em}{}
\setcounter{secnumdepth}{3}
\setcounter{tocdepth}{2}
\setlist[itemize]{leftmargin=2em,itemsep=0.25em,topsep=0.4em}
\setlist[enumerate]{leftmargin=2.3em,itemsep=0.25em,topsep=0.4em}
\setlength{\emergencystretch}{2em}

\newtheorem{theorem}{Theorem}[section]
\newtheorem{proposition}[theorem]{Proposition}
\newtheorem{lemma}[theorem]{Lemma}
\newtheorem{corollary}[theorem]{Corollary}
\newtheorem{conjecture}[theorem]{Conjecture}
\theoremstyle{definition}
\newtheorem{definition}[theorem]{Definition}
\newtheorem{algorithm}[theorem]{Algorithm}
\newtheorem{example}[theorem]{Example}
\theoremstyle{remark}
\newtheorem{remark}[theorem]{Remark}
\newtheorem{warning}[theorem]{Warning}

\newcommand{\R}{\mathbb R}
\newcommand{\C}{\mathbb C}
\newcommand{\Q}{\mathbb Q}
\newcommand{\N}{\mathbb N}
\newcommand{\Z}{\mathbb Z}
\newcommand{\Up}{\operatorname{up}}
\newcommand{\sinc}{\operatorname{sinc}}
\newcommand{\wt}{\operatorname{wt}_2}
\newcommand{\e}{\mathrm e}
\newcommand{\ii}{\mathrm i}
\newcommand{\dd}{\,\mathrm d}
\newcommand{\Li}{\operatorname{Li}}
\newcommand{\Var}{\operatorname{Var}}
\newcommand{\E}{\mathbb E}
\newcommand{\Prob}{\mathbb P}
\newcommand{\bigO}{\mathcal O}
\newcommand{\smallo}{o}
\newcommand{\supp}{\operatorname{supp}}
\newcommand{\sgn}{\operatorname{sgn}}
\newcommand{\lcm}{\operatorname{lcm}}
\newcommand{\vtwo}{v_2}
\newcommand{\qbinom}[3]{\genfrac{[}{]}{0pt}{}{#1}{#2}_{#3}}
\newcommand{\repo}{\nolinkurl{Analysis/FabiusFunction}}

\newenvironment{proofidea}{\par\smallskip\noindent\textit{Proof architecture.}\ }{\hfill$\square$\par\smallskip}
\newenvironment{boxedremark}{%
  \par\smallskip\noindent\begin{center}\begin{minipage}{0.94\textwidth}%
  \color{black}\hrule\vspace{0.6em}\small
}{%
  \vspace{0.6em}\hrule\end{minipage}\end{center}\smallskip
}

\lstdefinestyle{pseudo}{
  basicstyle=\ttfamily\small,
  backgroundcolor=\color{shadegray},
  frame=single,
  rulecolor=\color{rulegray},
  columns=fullflexible,
  keepspaces=true,
  showstringspaces=false,
  xleftmargin=0.7em,
  xrightmargin=0.7em,
  aboveskip=0.8em,
  belowskip=0.8em
}

% Glossary-specific additions, retaining the typography and palette of the
% reference article.
\newcommand{\glslink}[2]{\hyperlink{gls:#1}{#2}}
\newcommand{\letterlink}[1]{\hyperlink{letter:#1}{\textbf{#1}}}
\newcommand{\inactiveletter}[1]{\textcolor{rulegray}{\textbf{#1}}}
\newcommand{\alphabetbar}{%
  \begin{center}\small
  \letterlink{A}\enspace\textcolor{rulegray}{\textperiodcentered}\enspace  \letterlink{B}\enspace\textcolor{rulegray}{\textperiodcentered}\enspace  \letterlink{C}\enspace\textcolor{rulegray}{\textperiodcentered}\enspace  \letterlink{D}\enspace\textcolor{rulegray}{\textperiodcentered}\enspace  \letterlink{E}\enspace\textcolor{rulegray}{\textperiodcentered}\enspace  \letterlink{F}\enspace\textcolor{rulegray}{\textperiodcentered}\enspace  \letterlink{G}\enspace\textcolor{rulegray}{\textperiodcentered}\enspace  \letterlink{H}\enspace\textcolor{rulegray}{\textperiodcentered}\enspace  \letterlink{I}\enspace\textcolor{rulegray}{\textperiodcentered}\enspace  \inactiveletter{J}\enspace\textcolor{rulegray}{\textperiodcentered}\enspace  \inactiveletter{K}\enspace\textcolor{rulegray}{\textperiodcentered}\enspace  \letterlink{L}\enspace\textcolor{rulegray}{\textperiodcentered}\enspace  \letterlink{M}\\[-0.15em]
  \letterlink{N}\enspace\textcolor{rulegray}{\textperiodcentered}\enspace  \letterlink{O}\enspace\textcolor{rulegray}{\textperiodcentered}\enspace  \letterlink{P}\enspace\textcolor{rulegray}{\textperiodcentered}\enspace  \letterlink{Q}\enspace\textcolor{rulegray}{\textperiodcentered}\enspace  \letterlink{R}\enspace\textcolor{rulegray}{\textperiodcentered}\enspace  \letterlink{S}\enspace\textcolor{rulegray}{\textperiodcentered}\enspace  \letterlink{T}\enspace\textcolor{rulegray}{\textperiodcentered}\enspace  \letterlink{U}\enspace\textcolor{rulegray}{\textperiodcentered}\enspace  \letterlink{V}\enspace\textcolor{rulegray}{\textperiodcentered}\enspace  \letterlink{W}\enspace\textcolor{rulegray}{\textperiodcentered}\enspace  \inactiveletter{X}\enspace\textcolor{rulegray}{\textperiodcentered}\enspace  \inactiveletter{Y}\enspace\textcolor{rulegray}{\textperiodcentered}\enspace  \letterlink{Z}
  \end{center}
}
\newenvironment{glossaryentry}[3]{%
  \Needspace{8\baselineskip}%
  \par\addvspace{0.9em}%
  \hypertarget{gls:#1}{}%
  \noindent{\large\bfseries\color{linkblue}#2}\par
  \noindent{\small\color{black!64}\textit{Also:} #3}\par
  \smallskip\noindent
}{%
  \par\addvspace{0.45em}\noindent\textcolor{rulegray}{\rule{\linewidth}{0.35pt}}\par
}
\newcommand{\aliasindexentry}[3]{%
  \Needspace{2.7\baselineskip}%
  \noindent #1\hspace{0.25em}\allowbreak
  \textcolor{rulegray}{$\longrightarrow$}\hspace{0.25em}\allowbreak
  \glslink{#2}{#3}\par
}
\setlength{\columnsep}{1.6em}
\setlength{\columnseprule}{0.2pt}
\def\columnseprulecolor{\color{rulegray}}

\begin{document}
\hypersetup{pageanchor=false}
\pagenumbering{gobble}

\begin{titlepage}
\centering
\vspace*{1.55cm}
{\Huge\bfseries A Mathematical Glossary\par}
\vspace{0.3cm}
{\LARGE\color{linkblue} for the Fabius-Function Development\par}
\vspace{0.9cm}
\textcolor{linkblue}{\rule{0.82\textwidth}{1.3pt}}\par
\vspace{0.4cm}
\textcolor{rulegray}{\rule{0.62\textwidth}{0.7pt}}\par
\vspace{1.15cm}
{\Large 215 encyclopedic headwords\par}
\vspace{0.35cm}
{\large Real analysis, probability, Fourier theory, arithmetic,\par
$q$-calculus, asymptotics, approximation, and computability\par}
\vfill
{\large Source corpus\par}
\vspace{0.15cm}
{\ttfamily\large Analysis/FabiusFunction\par}
\vspace{0.7cm}
{\large Main branch inspected 25 August 2026\par}
\vfill
{\small Typography, page geometry, color palette, headings, and running-header style\par
follow \texttt{docs/Fabius\_Function\_and\_Rvachev\_Up/}\par
\texttt{Fabius\_Function\_and\_Rvachev\_Up.tex}.\par}
\end{titlepage}
\clearpage
\hypersetup{pageanchor=true}
\pagenumbering{roman}
\section*{Purpose and scope}
\addcontentsline{toc}{section}{Purpose and scope}
This glossary is a semantic inventory of the mathematical vocabulary used in
Lean comments and module docstrings, Markdown documentation, and LaTeX sources
under \repo.  It includes both standard concepts and repository-specific
objects whose names carry mathematical content.  Ordinary words such as
``number,'' ``function,'' ``set,'' ``sum,'' and ``proof'' are omitted unless a
qualified term has a distinct technical meaning.

Closely synonymous forms are consolidated under one headword and listed on the
\emph{Also} line.  For example, \emph{$q$-binomial coefficient} and
\emph{Gaussian binomial coefficient} share one definition.  Homonyms are
kept distinct when the corpus uses them differently: a probability distribution
is not a Schwartz distribution, and the bounded Fabius CDF is not the signed
global extension.  Each entry gives an encyclopedic definition first and then
explains its role, normalization, or formalization-sensitive boundary in this
development.

\begin{boxedremark}
\textbf{Coverage convention.}  Build-system language, tactic names, file
organization, and raw Lean syntax are excluded unless they denote a mathematical
interface---for example \texttt{HasSum}, filter convergence, or the
\texttt{IsFabius} specification.  Spelling variants and plural forms do not
receive duplicate entries.  The alias index at the end points alternative names
to their consolidated headwords.
\end{boxedremark}

\section*{Notation and object conventions}
\addcontentsline{toc}{section}{Notation and object conventions}
\begin{itemize}
  \item $F$ denotes the bounded, CDF-style Fabius function unless the context
        explicitly says \emph{signed global extension}.
  \item $\mathcal F$ denotes that signed global extension, agreeing with $F$ on
        $[0,1]$ but not clamped outside it.
  \item $\Up$ denotes Rvachev's compactly supported even up-function.
  \item $s_2(n)$ or $\wt(n)$ denotes binary digit sum; $v_2$ or $\nu_2$ denotes
        the $2$-adic valuation.
  \item The Fourier convention is
        $\widehat f(\xi)=\int_\R f(x)e^{-2\pi\ii x\xi}\dd x$.
  \item Empty sums are $0$, empty products are $1$, and natural powers use
        $0^0=1$ where the finite identities require the order-zero case.
\end{itemize}

\section*{Subject guide}
\addcontentsline{toc}{section}{Subject guide}
\begin{tabularx}{\textwidth}{@{}>{\raggedright\arraybackslash\bfseries}p{0.25\textwidth}>{\raggedright\arraybackslash}X@{}}
\toprule
Core objects &
\glslink{fabius-function}{bounded Fabius function},
\glslink{signed-global-fabius}{signed global extension},
\glslink{rvachev-up-function}{Rvachev's up-function},
\glslink{is-fabius-specification}{\texttt{IsFabius}},
\glslink{fixed-point}{fixed point}, and
\glslink{refinement-equation}{refinement equation}.\\
\addlinespace
Probability and Fourier theory &
\glslink{weighted-sum}{weighted random sum},
\glslink{cumulative-distribution}{CDF},
\glslink{density}{density},
\glslink{infinite-convolution}{infinite convolution},
\glslink{fourier-transform}{Fourier transform},
\glslink{sinc-function}{sinc}, and
\glslink{poisson-summation}{Poisson summation}.\\
\addlinespace
Exact arithmetic &
\glslink{dyadic-rational}{dyadic rational},
\glslink{moment}{moment},
\glslink{half-moment}{half moment},
\glslink{moment-numerators}{moment numerators},
\glslink{reshetnikov-numbers}{Reshetnikov numbers},
\glslink{dyadic-denominator}{common denominator}, and
\glslink{two-adic-valuation}{$2$-adic valuation}.\\
\addlinespace
Binary and $q$-combinatorics &
\glslink{thue-morse}{Thue--Morse},
\glslink{prouhet-thue-morse}{Prouhet cancellation},
\glslink{q-pochhammer}{$q$-Pochhammer symbol},
\glslink{gaussian-binomial}{Gaussian binomial coefficient}, and
\glslink{geometric-nodes}{geometric interpolation nodes}.\\
\addlinespace
Endpoint asymptotics &
\glslink{negative-laplace-product}{negative-Laplace product},
\glslink{mellin-transform}{Mellin transform},
\glslink{finite-part}{finite part},
\glslink{periodic-correction}{periodic correction},
\glslink{lambert-w}{Lambert $W$},
\glslink{bromwich-inversion}{Bromwich inversion}, and
\glslink{saddle-point}{saddle point}.\\
\addlinespace
Approximation and computability &
\glslink{centered-finite-spline}{centered finite spline},
\glslink{fourier-legendre-series}{Fourier--Legendre series},
\glslink{least-squares}{least squares},
\glslink{computable-real-sequence}{fast dyadic name}, and
\glslink{grzegorczyk-computability}{Grzegorczyk computability}.\\
\bottomrule
\end{tabularx}

\tableofcontents
\clearpage
\pagenumbering{arabic}

\section*{Alphabetical glossary}
\addcontentsline{toc}{section}{Alphabetical glossary}
\alphabetbar
\clearpage
\hypertarget{letter:A}{}
\section*{\color{linkblue}A}
\addcontentsline{toc}{section}{A}
\markboth{A}{A}
\alphabetbar
\begin{glossaryentry}{absolute-continuity}{Absolute continuity}{absolutely continuous function; absolutely continuous measure}
A function $f:[a,b]\to\mathbb R$ is \emph{absolutely continuous} if for every $\varepsilon>0$ there is $\delta>0$ such that, for every finite disjoint family of intervals $(x_j,y_j)$ with total length less than $\delta$, one has $\sum_j|f(y_j)-f(x_j)|<\varepsilon$.  Equivalently, $f$ has an integrable derivative almost everywhere and satisfies $f(x)=f(a)+\int_a^x f'(t)\,dt$.  A finite measure is absolutely continuous with respect to Lebesgue measure when every Lebesgue-null set has measure zero; the Radon--Nikodym theorem then represents it by a density.  The bounded Fabius function is far smoother than merely absolutely continuous: it is $C^\infty$, and its derivative is a compactly supported density obtained from Rvachev's up-function.  Nevertheless, absolute continuity is the correct measure-theoretic bridge between its cumulative-distribution and density descriptions.
\end{glossaryentry}
\begin{glossaryentry}{absolute-convergence}{Absolute convergence}{absolutely convergent series}
A numerical or function series $\sum_{n\ge 0} a_n$ is \emph{absolutely convergent} when $\sum_{n\ge 0}|a_n|<\infty$.  Absolute convergence implies ordinary convergence and, in complete normed spaces, allows many operations that conditional convergence does not: terms may be rearranged, finite sums may be interchanged with the series, and dominated-convergence arguments often justify passage of limits through integrals.  For a series of functions, one usually asks for pointwise absolute convergence, or the stronger uniform estimate $\sum_n\sup_{x\in E}|a_n(x)|<\infty$.  The latter is called normal convergence and implies absolute and uniform convergence on $E$.  In the Fabius development, absolute convergence is essential for the signed Thue--Morse continuation, the Fourier--Legendre expansion, the Gamma--zeta Fourier series of the periodic correction, and the binary-reduction series.  It is also what makes the order of several nested sums mathematically harmless rather than merely suggestive.
\end{glossaryentry}
\begin{glossaryentry}{admissible-class}{Admissible class}{invariant class; constraint set}
An \emph{admissible class} is a set of candidate objects on which a construction or variational argument is carried out.  Its defining conditions are chosen so that the operator under study maps the class into itself.  In a fixed-point proof one also wants the class to be closed in a complete metric space, ensuring that a contraction has a fixed point inside it.  For the bounded Fabius function, the natural class consists of continuous functions on $[0,1]$ taking values in $[0,1]$, satisfying the reflection law $f(1-x)=1-f(x)$ and the relevant endpoint normalization.  The integral transform used in the existence proof preserves these properties.  Calling the class admissible is not mere bookkeeping: positivity, symmetry, and the half-integral normalization are built into the space before differentiability is recovered from the fixed-point equation.
\end{glossaryentry}
\begin{glossaryentry}{almost-everywhere}{Almost everywhere}{a.e.; almost-everywhere equality}
A property holds \emph{almost everywhere} with respect to a measure $\mu$ if the set on which it fails has $\mu$-measure zero.  Two measurable functions that agree almost everywhere define the same Lebesgue integral and the same element of an $L^p$ space.  Almost-everywhere language is indispensable when integrability is proved by domination or when derivatives exist only in the measure-theoretic sense.  Lean's measure theory expresses this through filters such as $\mathsf{ae}$ and restricted measures.  In the Fabius corpus, many integrability arguments identify kernels pointwise on a positive interval and then promote the identity to an almost-everywhere equality under the restricted Lebesgue measure.  The final Fabius and up-functions are continuous or smooth, so their principal identities are pointwise; the a.e. layer is mainly the machinery that licenses termwise integration and changes of variables.
\end{glossaryentry}
\begin{glossaryentry}{almost-sure-convergence}{Almost-sure convergence}{convergence with probability one}
A sequence of random variables $X_n$ converges \emph{almost surely} to $X$ when $\mathbb P\{\omega:X_n(\omega)\to X(\omega)\}=1$.  This is stronger than convergence in probability and usually stronger than convergence in distribution.  For the Fabius probability model, the random series $X=\sum_{k\ge1}2^{-k}U_k$, with independent $U_k$ uniform on $[0,1]$, converges for every outcome because $0\le U_k\le1$ and $\sum 2^{-k}=1$.  Thus its convergence is actually pointwise, and hence automatically almost sure.  The distinction still matters conceptually: the law of the limiting weighted sum is obtained from finite partial sums, whose distributions are finite convolutions.  Almost-sure convergence explains the existence of the random variable, while weak convergence or convergence of characteristic functions identifies its distribution.
\end{glossaryentry}
\begin{glossaryentry}{analytic-continuation}{Analytic continuation}{holomorphic continuation}
Analytic continuation extends an analytic function beyond its original domain by requiring agreement on overlaps.  In complex analysis, two holomorphic functions that agree on a set with an accumulation point agree throughout every connected common domain.  The Fourier transform of a compactly supported smooth function is initially defined by an integral on the real axis but extends to an entire function of a complex variable; this is a basic form of analytic continuation and is quantified by the Paley--Wiener theorem.  In the Fabius development, the sinc product gives the continuation explicitly.  By contrast, the Fabius function itself has no nontrivial real-analytic continuation through points of $[0,1]$.  Thus the transform is entire while the original function is nowhere analytic on its support---a useful reminder that Fourier transformation exchanges localization in one variable for analytic growth in the dual variable.
\end{glossaryentry}
\begin{glossaryentry}{analytic-function}{Analytic function and analytic locus}{real analytic; analytic at a point; \texttt{AnalyticAt}}
A real function is \emph{analytic at} $x_0$ if it agrees in some neighborhood of $x_0$ with a convergent power series $\sum_{n\ge0}a_n(x-x_0)^n$.  Its \emph{analytic locus} is the set of points where this holds.  Smoothness is necessary but not sufficient: a $C^\infty$ function may have a Taylor series that converges yet fails to reproduce the function.  The bounded Fabius function is analytic exactly outside $[0,1]$, where it is locally constant, and nowhere analytic on the interval where its nontrivial behavior lives.  Rvachev's up-function is analytic off its support and nowhere analytic on $[-1,1]$.  The obstruction is especially transparent at dyadic points: sufficiently high derivatives vanish there, so the Taylor series is a polynomial, whereas the function is not locally polynomial.  At the endpoints every derivative vanishes, making the Taylor series identically zero despite nonzero nearby values.
\end{glossaryentry}
\begin{glossaryentry}{approximate-identity}{Approximate identity}{smoothing kernel; averaging kernel}
An \emph{approximate identity} is a family of integrable kernels $K_\varepsilon$ whose mass is one and which concentrate near the origin as $\varepsilon\to0$.  Convolution $f*K_\varepsilon$ then approaches $f$ in an appropriate topology.  A compactly supported smooth function whose integer translates form a partition of unity is also an exact local averaging kernel and can serve in approximation schemes.  Rvachev's up-function has this flavor: it is nonnegative, compactly supported, smooth, and its integer translates sum to one.  Dilated versions reproduce polynomials to increasing degree.  The Fabius corpus therefore connects the up-function to approximation theory, splines, and wavelet-like constructions, even though its primary development concerns exact identities rather than a general theory of approximate identities.
\end{glossaryentry}
\begin{glossaryentry}{approximation-theory}{Approximation theory}{polynomial approximation; numerical approximation}
Approximation theory studies how complicated functions or data can be represented by simpler objects---polynomials, splines, trigonometric sums, rational functions, wavelets, or finite-dimensional subspaces---and how the error behaves.  The Fabius development contains several distinct approximation mechanisms.  Contraction iterates converge geometrically to the fixed point; centered finite splines approximate the Fabius distribution with an explicit uniform error; Fourier--Legendre partial sums are least-squares-optimal polynomials; and saddle expansions approximate the logarithm near an endpoint.  These approximations answer different questions and use different metrics.  Uniform approximation controls the worst pointwise error, $L^2$ approximation controls mean-square error, and asymptotic approximation controls a limit regime.  Treating them as interchangeable is a common source of false claims---for example, a visually excellent compact-interval fit need not have the correct endpoint asymptotic scale.
\end{glossaryentry}
\begin{glossaryentry}{asymptotic-equivalent}{Asymptotic equivalence}{asymptotic equivalent; $f\sim g$}
For functions $f$ and $g$ defined near a limit point, $f(x)\sim g(x)$ means $f(x)/g(x)\to1$.  This is stronger than having the same logarithmic order or the same leading exponent.  If $f$ and $g$ are positive and exponentially small, even a bounded additive error in $\log f-\log g$ may prevent asymptotic equivalence after exponentiation.  The corrected Fabius endpoint formula first establishes a sharp expansion for $\log F(x)$ in the exact Lambert phase, including a nonconstant periodic term.  Exponentiating the controlled logarithmic remainder then yields an asymptotic equivalent for $F(x)$.  By contrast, the quotient-of-exponentials approximation audited in the repository has the wrong boundary scale: it may track the graph well away from the endpoint but is little-$o$ of the displaced Fabius bump there, so their ratio tends to zero rather than one.
\end{glossaryentry}
\begin{glossaryentry}{asymptotic-expansion}{Asymptotic expansion}{Poincar\'e expansion; all-orders expansion}
A sequence of functions $\phi_0,\phi_1,\ldots$ with $\phi_{n+1}=o(\phi_n)$ is an \emph{asymptotic scale}.  A function $f$ has expansion $f\sim\sum_{n\ge0}a_n\phi_n$ when, for every $N$, subtracting the first $N$ terms leaves $o(\phi_{N-1})$ or a specified $O(\phi_N)$ remainder.  The series need not converge for any fixed argument.  The Fabius small-argument expansion is naturally written in powers of $1/\lambda(x)$, where $\lambda$ is the lower-Lambert phase, not directly in powers of $1/|\log x|$.  Its coefficients are smooth one-periodic functions, so the expansion is a periodic-coefficient Poincar\'e expansion.  ``All orders'' means that a remainder theorem exists for every finite truncation; it does not assert convergence of the infinite formal series.
\end{glossaryentry}
\begin{glossaryentry}{atomic-function}{Atomic function}{Rvachev atomic function; finite function}
In Rvachev's terminology, an \emph{atomic function} is a compactly supported infinitely differentiable solution of a functional--differential or refinement equation, used as a basic building block for approximation and numerical methods.  Rvachev's up-function is the canonical example.  Its derivative is a finite linear combination of translated and dilated copies of itself,
\[
  \Up'(x)=2\bigl(\Up(2x+1)-\Up(2x-1)\bigr),
\]
and its support is exactly $[-1,1]$.  Atomic functions combine features that are usually in tension: compact support, arbitrary smoothness, exact self-similarity, and polynomial-reproduction identities.  The word ``atomic'' here has nothing to do with atoms of a measure; it refers to a fundamental compactly supported element from which larger approximation spaces may be assembled.
\end{glossaryentry}
\begin{glossaryentry}{attained-bound}{Attained bound and sharp constant}{optimal bound; least constant; exact supremum}
A bound is \emph{sharp} if its constant cannot be improved.  One strong way to prove sharpness is to exhibit a point where equality holds; then the bound equals the supremum of the quantity being estimated.  The Fabius formalization distinguishes carefully between a valid upper estimate and an optimal one.  For example, the global Lipschitz constant of the bounded Fabius function is $2$, and $2$ is the least possible constant.  Likewise,
\[
 \sup_x |F^{(m)}(x)|=2^{\binom{m+1}{2}},
\]
with equality at an explicit dyadic point, and the analogous derivative bound for the up-function is attained.  This distinction matters because words such as ``sharp,'' ``exact,'' and ``optimal'' require a lower-bound or equality argument, not merely a good-looking one-sided estimate.
\end{glossaryentry}
\begin{glossaryentry}{automatic-sequence}{Automatic sequence}{$2$-automatic sequence}
A sequence is $k$-automatic if a finite automaton can produce its $n$th term from the base-$k$ digits of $n$.  The Thue--Morse sequence is the archetypal $2$-automatic sequence: its bit is the parity of the binary digit sum, and its sign is $(-1)^{\operatorname{wt}_2(n)}$.  Automaticity explains the exact block recursions
$t_{2n}=t_n$ and $t_{2n+1}=1-t_n$ for the zero-one version, or $\varepsilon_{2n}=\varepsilon_n$ and $\varepsilon_{2n+1}=-\varepsilon_n$ for the sign version.  The Fabius signed extension and its dyadic formulas use these block laws repeatedly.  Automatic sequences often have low symbolic complexity but highly nonperiodic behavior; that combination is precisely what makes Thue--Morse signs useful for structured cancellation.
\end{glossaryentry}
\clearpage
\hypertarget{letter:B}{}
\section*{\color{linkblue}B}
\addcontentsline{toc}{section}{B}
\markboth{B}{B}
\alphabetbar
\begin{glossaryentry}{b-spline}{B-spline}{box spline; discrete B-spline law}
A univariate B-spline is a compactly supported piecewise polynomial obtained by repeated convolution of interval indicators.  The order-$m$ cardinal B-spline is the density of a sum of $m$ independent uniform random variables, up to translation and scaling.  It has degree $m-1$, smoothness increasing with $m$, and a refinement relation.  Finite truncations of the random series defining the Fabius law therefore have spline densities and spline cumulative distributions.  The centered finite splines in the repository are rescaled, shifted versions tailored to the dyadic grid.  Iterated prefix sums of the Thue--Morse sequence satisfy a discrete B-spline convolution law, giving a combinatorial counterpart of the same smoothing mechanism.
\end{glossaryentry}
\begin{glossaryentry}{banach-fixed-point}{Banach fixed-point theorem}{contraction mapping theorem}
The Banach fixed-point theorem states that a strict contraction $T$ on a nonempty complete metric space has a unique fixed point $x_*$, and the iterates $T^n x_0$ converge to it from every starting point.  If $d(Tx,Ty)\le qd(x,y)$ with $0<q<1$, then
\[
 d(T^n x_0,x_*)\le \frac{q^n}{1-q}d(Tx_0,x_0).
\]
The bounded Fabius function is constructed this way on a closed class of continuous symmetric $[0,1]$-valued functions.  The operator integrates a rescaled candidate on each half of the unit interval.  The contraction estimate proves existence and uniqueness simultaneously and also yields a certified numerical scheme.  Smoothness is not assumed in the metric space; it is bootstrapped afterward from the integral fixed-point equation.
\end{glossaryentry}
\begin{glossaryentry}{bell-polynomial}{Bell-polynomial and exponential-coefficient recurrence}{formal exponential recurrence; exponential formula}
When a formal power series $E(z)=\sum_{n\ge1}E_nz^n$ is exponentiated, the coefficients of $\exp E(z)$ are universal polynomials in $E_1,\ldots,E_n$.  They can be written with complete Bell polynomials or computed recursively from
\[
 b_0=1,\qquad (n+1)b_{n+1}=\sum_{j=0}^{n}(j+1)E_{j+1}b_{n-j}.
\]
The Fabius saddle expansion uses this recurrence under the name of an exponential-coefficient operator.  It converts the local exponent expansion near the saddle into mass coefficients before Gaussian integration.  The procedure is algebraic: once the exponent jets are known, no further analysis is hidden in it.  Bell-polynomial language clarifies why the coefficients are finite universal polynomials and why only finitely many jets can contribute at each asymptotic order.
\end{glossaryentry}
\begin{glossaryentry}{bernoulli-numbers}{Bernoulli numbers}{Bernoulli recurrence}
The Bernoulli numbers $B_n$ are defined by the exponential generating function
\[
  \frac{z}{e^z-1}=\sum_{n=0}^{\infty}B_n\frac{z^n}{n!},
\]
with $B_0=1$, $B_1=-\tfrac12$, $B_2=\tfrac16$, and $B_{2m+1}=0$ for $m\ge1$.  They occur throughout analysis and number theory---in Euler--Maclaurin summation, special values of the zeta function, and expansions of trigonometric kernels.  In the Fabius arithmetic layer, the functional equations for the moment and half-moment generating functions are combined with the power series for $z/(e^z-1)$ and $\pi z/\sin\pi z$.  Equating coefficients yields Bernoulli-number recurrences for both rational sequences.  These formulas are equivalent to the original triangular recurrences but expose a classical special-number structure.
\end{glossaryentry}
\begin{glossaryentry}{big-o-little-o}{Big-$O$ and little-$o$}{asymptotic order; Landau notation}
The statement $f(x)=O(g(x))$ near a limit means that $|f(x)|\le C|g(x)|$ eventually for some constant $C$.  The statement $f(x)=o(g(x))$ means $f(x)/g(x)\to0$.  Big-$O$ is a bound; little-$o$ is a strict asymptotic comparison.  In Lean these are expressed through filter-based relations such as \texttt{IsBigO} and \texttt{IsLittleO}, which make the limiting filter and scalar norm explicit.  The Fabius corpus uses these notions to state flatness, compare proposed approximants, control minor arcs, and formulate all-orders remainders.  An $O(1/\lambda)$ error cannot be replaced by $o(1)$ without proof, and an omitted bounded periodic term is not $O(1/\lambda)$.
\end{glossaryentry}
\begin{glossaryentry}{bijection-inverse}{Bijection and inverse function}{one-to-one correspondence; totalized inverse}
A map is a \emph{bijection} when it is both injective and surjective; it then has a two-sided inverse on its codomain.  The bounded Fabius function is continuous and strictly increasing from $0$ to $1$ on the unit interval, hence is a homeomorphism $[0,1]\to[0,1]$.  Its inverse can therefore be defined intrinsically on this interval.  A \emph{totalized inverse} extends the inverse to all real inputs by assigning endpoint values outside the range, for example $0$ below $0$ and $1$ above $1$.  The development proves quantitative endpoint behavior of this inverse: because $F$ is flatter than every power at zero, $F^{-1}(y)$ tends to zero more slowly than every positive power of $y$---equivalently, the inverse is steeper than every root near the endpoint.
\end{glossaryentry}
\begin{glossaryentry}{bilateral-laplace-transform}{Bilateral Laplace transform}{two-sided Laplace transform}
The bilateral Laplace transform of an integrable function $f$ is
\[
  \mathcal L_b f(s)=\int_{-\infty}^{\infty}e^{-sx}f(x)\,dx,
\]
for complex $s$ in the region where the integral converges.  It is a moment-generating function with a sign convention and reduces to the Fourier transform on the imaginary axis.  For a compactly supported density it is entire.  In the Fabius analysis, a centered or signed version of the transform factors into an infinite product of elementary terms.  Evaluating its logarithm on the negative real axis produces the ``negative-Laplace'' product from which the endpoint asymptotic and the exact periodic correction are extracted.  The adjective bilateral distinguishes this transform from the one-sided transform over $[0,\infty)$ used in many differential-equation problems.
\end{glossaryentry}
\begin{glossaryentry}{binary-expansion}{Binary expansion, set bit, and highest set bit}{base-two representation; most significant bit}
Every nonnegative integer has a unique finite expansion $n=\sum_j\varepsilon_j2^j$ with digits $\varepsilon_j\in\{0,1\}$.  A position with $\varepsilon_j=1$ is a \emph{set bit}; the largest such $j$ is the highest set bit or most significant bit.  Several Fabius algorithms follow the binary representation rather than the magnitude of the numerator.  The exact dyadic evaluator removes one highest set bit at each recursive step and evaluates a finite Taylor polynomial around the corresponding inverse power of two.  Consequently the number of recursive calls is the binary weight of the numerator.  Binary block decompositions also organize Thue--Morse cancellation, because adding a leading bit either preserves or reverses the signs of a whole block.
\end{glossaryentry}
\begin{glossaryentry}{binary-reduction}{Binary-reduction formula}{global binary recursion; bit-removal expansion}
A \emph{binary-reduction formula} evaluates the Fabius function by successively separating the largest dyadic scale present in an argument.  At exact dyadic inputs, each step removes one set bit of the numerator and adds a finite Taylor contribution; termination occurs when no set bits remain.  For arbitrary nonnegative real inputs, the analogous infinite series sums scale-by-scale floor terms and converges absolutely to the signed global extension.  The outer index must begin at zero for an all-endpoint theorem.  Binary reduction explains the efficient evaluator's dependence on binary weight rather than on the numerical size of the numerator.
\end{glossaryentry}
\begin{glossaryentry}{binary-weight}{Binary weight}{Hamming weight; sum of binary digits; $\operatorname{wt}_2(n)$}
The binary or Hamming weight $\operatorname{wt}_2(n)$ is the number of $1$ digits in the base-two expansion of $n$.  It controls the Thue--Morse sign $\varepsilon_n=(-1)^{\operatorname{wt}_2(n)}$ and appears in parity counts for binomial coefficients.  Lucas's theorem implies that the number of odd entries in the $n$th row of Pascal's triangle is $2^{\operatorname{wt}_2(n)}$.  In the Fabius corpus the same statistic has an algorithmic meaning: evaluating $F(a/2^n)$ by the bit recursion takes one recursive descent per set bit of $a$.  Thus binary weight links combinatorial parity, automatic sequences, and exact-computation complexity.
\end{glossaryentry}
\begin{glossaryentry}{binomial-coefficient}{Binomial coefficient}{combinatorial number; Pascal coefficient}
For integers $0\le k\le n$, the binomial coefficient is
\[
  \binom nk=\frac{n!}{k!(n-k)!},
\]
and is extended by zero outside that range.  It counts $k$-element subsets of an $n$-element set and is the coefficient of $x^k$ in $(1+x)^n$.  Binomial coefficients permeate the Fabius development: they occur in moment recurrences, Taylor expansions, Legendre polynomials, dyadic closed forms, and the $q$-binomial analogues.  Their parity is analyzed through Lucas's theorem.  The formalization also repeatedly uses triangular identities such as $\binom{n+1}{2}=\binom n2+n$, because powers of two indexed by triangular numbers arise when the differential scaling law is iterated.
\end{glossaryentry}
\begin{glossaryentry}{borel-radon-measure}{Borel and Radon measure}{Radon probability measure}
A Borel measure is defined on the sigma-algebra generated by open sets.  On a locally compact Hausdorff space, a Radon measure is locally finite and regular: its values can be approximated from inside by compact sets and from outside by open sets.  Probability laws on $\mathbb R$ are naturally Radon measures.  The early construction of Rvachev's function considers finite convolution measures whose Fourier transforms are finite sinc products and then passes to a weak limit.  The limiting Radon probability measure has a smooth compactly supported density.  Radon regularity is what allows measure-theoretic convergence statements to be connected with pointwise distribution functions and ordinary integrals.
\end{glossaryentry}
\begin{glossaryentry}{bounded-variation}{Bounded variation}{total variation}
A function $f:[a,b]\to\mathbb C$ has bounded variation if
\[
  \sup_{a=x_0<\cdots<x_m=b}\sum_{j=1}^m|f(x_j)-f(x_{j-1})|<\infty.
\]
The supremum is its total variation.  Such functions are differences of two monotone real functions in the real-valued case, have one-sided limits everywhere, and define finite signed Stieltjes measures.  Bounded variation is also a robust hypothesis for summation and averaging arguments.  In the Fabius discrete-limit proof, a Toeplitz-type averaging lemma controls finite shifted spline rows by the variation of the limiting profile.  The hypothesis is strong enough to bound the cumulative effect of small coefficient changes but weak enough to include the compactly supported piecewise-polynomial approximants used there.
\end{glossaryentry}
\begin{glossaryentry}{bromwich-inversion}{Bromwich inversion formula}{inverse Laplace integral; Bromwich contour}
The Bromwich inversion formula recovers a function from its one-sided Laplace transform by integration along a vertical line in the complex plane:
\[
 f(x)=\frac{1}{2\pi i}\int_{c-i\infty}^{c+i\infty}e^{sx}\mathcal Lf(s)\,ds,
\]
where $c$ lies to the right of the transform's singularities.  For the Fabius endpoint problem, an analogous contour integral expresses the small value $F(x)$ in terms of the negative-Laplace product.  The contour is moved or rescaled so that it passes through the positive saddle radius.  A central neighborhood is approximated by a Gaussian integral, while vertical tails and minor arcs are bounded quantitatively.  ``Bromwich saddle'' in module names refers to this combination of inverse-Laplace representation and saddle-point analysis.
\end{glossaryentry}
\begin{glossaryentry}{bump-function}{Bump function}{smooth compactly supported function}
A \emph{bump function} is a $C^\infty$ function with compact support.  Standard examples are built from $e^{-1/(1-x^2)}$ inside $(-1,1)$ and zero outside, but their derivatives have no comparably simple self-similarity.  Rvachev's up-function is a distinguished bump: it is nonnegative, even, supported exactly on $[-1,1]$, equals one at the origin, and satisfies an exact functional--differential equation.  Its Fourier transform is an explicit infinite sinc product, and its values and derivatives at dyadic points are rational.  The existence of such a highly structured bump is the original motivation behind the historical construction formalized in the repository.
\end{glossaryentry}
\begin{glossaryentry}{bose-kernel}{Logarithmic Bose kernel}{Bose--Mellin kernel; $\log(1-e^{-x})$}
The logarithmic Bose kernel is the function
\[
  B(x)=\log(1-e^{-x}),\qquad x>0.
\]
It is named for its relation to Bose--Einstein integrals and the geometric-series expansion of $1/(e^x-1)$.  Near zero, $B(x)=\log x+O(x)$; at infinity it decays like $-e^{-x}$.  Its Mellin transform is
\[
  \int_0^\infty x^{a-1}B(x)\,dx=-\Gamma(a)\zeta(a+1),\qquad a>0.
\]
The Fabius asymptotic analysis uses this identity to compute the mean and Fourier coefficients of the periodic correction.  Because the Mellin transform has a double pole at $a=0$, the constant term is recovered through a finite-part regularization that subtracts the local $\log x$ singularity and splits the integral at a positive cutoff.
\end{glossaryentry}
\clearpage
\hypertarget{letter:C}{}
\section*{\color{linkblue}C}
\addcontentsline{toc}{section}{C}
\markboth{C}{C}
\alphabetbar
\begin{glossaryentry}{centered-finite-spline}{Centered finite spline}{centered uniform spline; finite-spline approximant}
A centered finite spline in this development is the cumulative distribution or density associated with a finite partial sum of centered, dyadically weighted uniform random variables.  Equivalently, it is built from a finite convolution of interval indicators and then translated so that the support is symmetric.  It is piecewise polynomial on a dyadic grid, takes values in $[0,1]$, and converges uniformly to the bounded Fabius function with an explicit error controlled by the omitted tail of the random series.  Because its values on dyadic inputs can be computed by natural-number operations, the centered spline is also the bridge to primitive-recursive computability.  The degree-zero case is a step function and requires separate endpoint conventions; higher degrees are ordinary continuous splines.
\end{glossaryentry}
\begin{glossaryentry}{centering}{Centering and mean-zero normalization}{centered periodic correction; zero-mean function}
To \emph{center} a random variable one subtracts its expectation; to center a periodic function one subtracts its average over one period.  If $P$ is one-periodic, its mean is $\bar P=\int_0^1P(t)\,dt$ and the centered function $\Psi=P-\bar P$ satisfies $\int_0^1\Psi=0$.  Centering separates the constant term from genuine oscillation.  In the Fabius endpoint asymptotic, $P$ arises from an exact dyadic decomposition of the logarithmic Laplace product, and $\Psi$ is the nonconstant periodic correction appearing in the sharp main term.  Its zero Fourier coefficient vanishes by construction, while every nonzero mode has an explicit Gamma--zeta coefficient.
\end{glossaryentry}
\begin{glossaryentry}{characteristic-function}{Characteristic function}{Fourier transform of a probability law}
The characteristic function of a real random variable $X$ is $\varphi_X(t)=\mathbb E[e^{itX}]$.  It always exists, uniquely determines the law of $X$, and turns sums of independent random variables into products: $\varphi_{X+Y}=\varphi_X\varphi_Y$.  With a $2\pi$ Fourier convention the argument is rescaled accordingly.  For a uniform random variable, the centered characteristic function is a sinc factor.  Hence the infinite weighted uniform sum underlying the Fabius distribution has an infinite sinc-product characteristic function.  This probability interpretation and the Fourier-transform description of the up-function are the same identity in different normalizations.
\end{glossaryentry}
\begin{glossaryentry}{coefficient-extraction}{Coefficient extraction}{$[z^n]f(z)$}
For a formal or convergent power series $f(z)=\sum_{n\ge0}a_nz^n$, the notation $[z^n]f(z)$ denotes the coefficient $a_n$.  Coefficient extraction converts functional identities into recurrences.  Comparing the sinc-product expansion with a moment series yields the rational moment recurrence; comparing a generating-function dilation equation yields the half-moment recurrence; and extracting coefficients from formal exponentials and logarithms produces the saddle corrections.  The operation is algebraic in formal power series, so convergence is irrelevant there, but analytic coefficient extraction additionally requires a justified series expansion.  The glossary uses ``coefficient'' broadly for Fourier, Taylor, Legendre, generating-function, and asymptotic coefficients; their ambient algebra and convergence meaning must be kept distinct.
\end{glossaryentry}
\begin{glossaryentry}{compact-support}{Compact support and topological support}{support; \texttt{support}; \texttt{tsupport}}
The pointwise support of $f$ is $\{x:f(x)\ne0\}$, while its topological support is the closure of that set.  A function has compact support if its topological support is compact.  For the up-function, positivity holds exactly on $(-1,1)$, so the pointwise support is open and the topological support is $[-1,1]$.  For a continuous function these notions are related by closure, but formal statements must specify which one is meant.  Compact support makes every polynomial moment finite and makes the Fourier transform entire.  It also permits exact local finiteness of the Thue--Morse translate sum: at any fixed argument only finitely many translated bumps are nonzero.
\end{glossaryentry}
\begin{glossaryentry}{complex-egf}{Complex exponential generating function}{entire moment-generating series}
An exponential generating function has the form $G(z)=\sum_{n\ge0}a_nz^n/n!$.  When the series converges for every complex $z$, it is an entire function.  Moment-generating functions are exponential generating functions of moments whenever they exist near the origin.  The Fabius half moments are encoded by an entire function satisfying a dyadic functional equation such as $G(2z)=((e^z-1)/z)G(z)$.  This equation yields a triangular recurrence for the coefficients and, after logarithmic or Bernoulli expansion, alternative recurrences.  Complexifying the generating function is useful because entire-function identities can be proved by power-series comparison and then specialized to real arguments.
\end{glossaryentry}
\begin{glossaryentry}{computable-real-function}{Computable real function}{Grzegorczyk computability; sequential computability}
In the Grzegorczyk-style definition used by the repository, a real function is computable when two conditions hold.  First, it is \emph{sequentially computable}: every computable real sequence is mapped to a computable real sequence, uniformly in index and precision.  Second, it is effectively uniformly continuous, with a recursive positive modulus converting an input tolerance into an output tolerance.  For the bounded Fabius function, a primitive-recursive centered-spline evaluator supplies output names, while the global Lipschitz bound supplies the modulus $d(n)=2n$ in the source's reciprocal convention.  Both ingredients are necessary: an algorithm on names without an effective continuity guarantee can depend pathologically on the chosen representation.
\end{glossaryentry}
\begin{glossaryentry}{computable-real-sequence}{Computable real number and computable real sequence}{fast Cauchy name; fast dyadic name}
A real number is computable when an algorithm produces rational approximations with a prescribed effective error.  A common representation is a fast dyadic name: at precision $p$, the algorithm returns a dyadic $a_p/2^p$ within $2^{-p}$ of the target.  A real sequence is computable when one algorithm works uniformly in both the sequence index and the requested precision.  The Fabius development uses a signed-natural pair $(a,b)$ to represent $(a-b)/2^p$, avoiding hidden integer coding.  This redundancy is extensionally harmless and makes primitive-recursive arithmetic explicit.  Computable names concern approximation of arbitrary reals; they are distinct from the exact rational evaluator at dyadic arguments.
\end{glossaryentry}
\begin{glossaryentry}{congruence-lucas}{Congruence and Lucas's theorem}{binomial coefficients modulo a prime}
Integers $a$ and $b$ are congruent modulo $m$, written $a\equiv b\pmod m$, when $m$ divides $a-b$.  Lucas's theorem computes a binomial coefficient modulo a prime $p$ digitwise.  If $n=\sum n_jp^j$ and $k=\sum k_jp^j$, then
\[
 \binom nk\equiv\prod_j\binom{n_j}{k_j}\pmod p.
\]
For $p=2$, $\binom nk$ is odd exactly when every set bit of $k$ is also a set bit of $n$.  Hence row $n$ of Pascal's triangle contains $2^{\operatorname{wt}_2(n)}$ odd entries.  The Fabius arithmetic proof uses this parity count to show that certain moment numerators and Reshetnikov numbers are odd, yielding the exact $2$-adic valuation of inverse-dyadic values.
\end{glossaryentry}
\begin{glossaryentry}{constant-polynomial}{Constant polynomial and translation independence}{common-shift invariance}
A polynomial is constant when all positive-degree coefficients vanish.  Several finite Fabius formulas contain a seemingly arbitrary translation parameter $q$ inside powers such as $(r+q)^N$.  The formalization proves that the whole finite sum is a constant polynomial in $q$.  The mechanism is cancellation of all positive-degree coefficients, ultimately supplied by Thue--Morse moment identities and a Gaussian $q$-binomial row.  Therefore the expression may be evaluated over any field containing $\mathbb Q$, and rational, real, complex, or Gaussian-rational shifts all give the same value.  This is stronger than checking equality for a few special shifts: it explains structurally why differently centered published formulas agree.
\end{glossaryentry}
\begin{glossaryentry}{contour-integral}{Contour integral}{complex line integral}
A contour integral integrates a complex function along a parametrized curve $\gamma$:
\[
 \int_\gamma f(z)\,dz=\int_a^b f(\gamma(t))\gamma'(t)\,dt.
\]
Holomorphicity allows contours to be deformed without changing the integral, provided singularities and endpoint behavior are controlled.  In the Fabius saddle analysis, the inverse-Laplace integral runs along a vertical contour through or near a positive real saddle.  The contour is divided into a central region, where a Taylor expansion of the complex phase is accurate, and tails or minor arcs, where real-part estimates force decay.  The principal logarithm is used only in a neighborhood where the contour avoids its branch cut.
\end{glossaryentry}
\begin{glossaryentry}{contraction}{Contraction}{strict contraction; contraction constant}
A map $T$ on a metric space is a contraction if there is $q<1$ such that $d(Tx,Ty)\le qd(x,y)$ for all $x,y$.  The smallest usable $q$ is a contraction constant.  Contractions are uniformly continuous, have at most one fixed point, and on complete invariant spaces have exactly one by the Banach theorem.  The Fabius integral operator contracts the sup norm because each output value integrates a difference over an interval of length at most one half.  The same estimate not only establishes uniqueness but bounds the error after finitely many iterations.  This is conceptually different from the later statement that the Fabius function itself is Lipschitz with optimal constant two.
\end{glossaryentry}
\begin{glossaryentry}{convexity-concavity}{Convexity, concavity, and inflection point}{strictly convex; strictly concave}
A function is convex on an interval if its graph lies below every chord, equivalently $f''\ge0$ when twice differentiable; it is concave if $f''\le0$.  Strict inequalities give strict convexity or concavity.  An inflection point is a point where the curvature changes sign.  The bounded Fabius function is strictly convex on the left half of $(0,1)$ and strictly concave on the right half, with its unique inflection point at $1/2$.  This follows from the derivative scaling law and symmetry.  The up-function has the corresponding single central maximum and opposite curvature pattern after folding.  Formal proofs must distinguish local sign statements for derivatives from global geometric conclusions, especially at endpoints where derivatives vanish.
\end{glossaryentry}
\begin{glossaryentry}{convolution}{Convolution}{self-convolution; convolution power}
For integrable functions on $\mathbb R$, convolution is
\[
 (f*g)(x)=\int_{\mathbb R}f(t)g(x-t)\,dt.
\]
For finite measures, $(\mu*\nu)(A)=\int\nu(A-x)\,d\mu(x)$; it is the law of the sum of independent random variables with laws $\mu$ and $\nu$.  Fourier transformation converts convolution into multiplication.  Finite truncations of the Fabius random series therefore have densities obtained by convolving scaled uniform densities, hence splines, while the limiting law is an infinite convolution.  The up-function also satisfies an exact self-convolution identity at an appropriate scale.  Repeated discrete convolution appears independently in the iterated-prefix-sum treatment of Thue--Morse.
\end{glossaryentry}
\begin{glossaryentry}{cumulant}{Cumulant}{cumulant coefficient}
Cumulants are the coefficients of the logarithm of a moment-generating or characteristic function.  If $K(t)=\log\mathbb E[e^{tX}]$, then the $n$th cumulant is $\kappa_n=K^{(n)}(0)$.  The first two are the mean and variance; higher cumulants measure non-Gaussian shape and add under independent sums.  In saddle-point analysis, expanding a logarithmic transform is essentially a cumulant expansion: the quadratic term produces the Gaussian core, and cubic and higher terms produce corrections.  The Fabius dyadic and saddle modules use ``cumulant'' for normalized logarithmic derivatives or coefficient combinations that isolate this additive structure.
\end{glossaryentry}
\begin{glossaryentry}{cumulant-generating}{Cumulant-generating function}{log moment-generating function; log transform}
The cumulant-generating function is $K(t)=\log M(t)$, where $M(t)=\mathbb E[e^{tX}]$ is the moment-generating function.  Taking a logarithm converts products associated with independent sums into sums and exposes additive scale recurrences.  The Fabius negative-Laplace analysis studies a logarithm $\Lambda(s)$ of an infinite product.  After subtracting a quadratic drift in $\log_2 s$ and a forward tail, the remaining part is exactly periodic.  Derivatives of $\Lambda$ determine the saddle location, local Gaussian width, and higher exponent jets.  A branch of the logarithm must be chosen consistently when the argument is complex; on the positive real axis the ordinary real logarithm suffices.
\end{glossaryentry}
\begin{glossaryentry}{cumulative-distribution}{Cumulative distribution function}{distribution function; CDF}
For a real random variable $X$, its cumulative distribution function is $F_X(x)=\mathbb P(X\le x)$.  Every CDF is nondecreasing, right-continuous, and tends to $0$ and $1$ at the two infinities.  If the law has a continuous density $f$, then $F_X'(x)=f(x)$.  The bounded Fabius function is the CDF of the weighted uniform sum $\sum_{k\ge1}2^{-k}U_k$.  It is unusually regular: it is $C^\infty$, constant outside $[0,1]$, symmetric by $F(x)+F(1-x)=1$, and strictly increasing in the interior.  The term ``bounded Fabius function'' is used in the repository to distinguish this genuine CDF from the signed global continuation, which can be negative and is not a probability distribution function.
\end{glossaryentry}
\begin{glossaryentry}{convolution-measures}{Finite and infinite convolution of measures}{convolution measure; convolution product}
A finite convolution $\mu_1*\cdots*\mu_n$ is the distribution of $X_1+\cdots+X_n$ for independent $X_j$ with laws $\mu_j$.  An infinite convolution is a weak limit of these finite products, when it exists.  Its characteristic function is the corresponding infinite product of characteristic functions.  In the original construction of the up-function, each factor is a centered uniform measure on a shrinking interval.  The supports have summable radii, so the sums converge and the limiting measure is compactly supported.  Increasing convolution smoothness ultimately gives a $C^\infty$ density.  The repository formalizes both the weak convergence of measures and the pointwise convergence of corrected step or spline densities.
\end{glossaryentry}
\begin{glossaryentry}{compact-uniform-convergence}{Uniform convergence on compact sets}{locally uniform convergence; compact convergence}
A sequence $f_n$ converges uniformly on compact sets to $f$ if, for every compact $K$, $\sup_{x\in K}|f_n(x)-f(x)|\to0$.  This topology is weaker than global uniform convergence but strong enough to preserve continuity and, under suitable derivative bounds, analyticity.  Infinite products and series defining the Fabius Fourier transform and the periodic correction are often handled locally on compact positive intervals.  A summable majorant independent of $x\in K$ gives uniform convergence by the Weierstrass $M$-test.  Local uniform convergence is particularly natural when the behavior at zero or infinity is singular and no single global bound can cover the whole domain.
\end{glossaryentry}
\clearpage
\hypertarget{letter:D}{}
\section*{\color{linkblue}D}
\addcontentsline{toc}{section}{D}
\markboth{D}{D}
\alphabetbar
\begin{glossaryentry}{delay-equation}{Delay and dilation equation}{functional--differential equation; pantograph-type equation}
A \emph{delay equation} relates the value or derivative of an unknown function at one argument to values at shifted or rescaled arguments.  The equations governing the Fabius and Rvachev functions are functional--differential rather than ordinary differential equations because they contain compositions such as $F(2x)$ and $\Up(2x\pm1)$.  The basic laws are
\[
  F'(x)=2F(2x),\qquad
  \Up'(x)=2\bigl(\Up(2x+1)-\Up(2x-1)\bigr),
\]
with the first interpreted on the appropriate domain or for the signed global extension.  Dilation makes the equation recursive across scales: information near an endpoint determines derivatives at larger arguments, while compact support eventually terminates many derivative formulas.  Unlike a conventional initial-value ODE, uniqueness depends on global support, symmetry, normalization, or boundedness conditions in addition to local derivative data.
\end{glossaryentry}
\begin{glossaryentry}{derivative-jet}{Derivative jet}{jet; finite jet; saddle jet}
The \emph{$n$-jet} of a smooth function at a point is the finite list of derivatives through order $n$, usually encoded as $(f(x),f'(x),\ldots,f^{(n)}(x))$.  Jets retain exactly the local data needed for a Taylor polynomial while forgetting the remainder.  In saddle-point analysis, the term also describes normalized derivatives of an exponent at the moving saddle.  The Fabius asymptotic development introduces periodic and bounded saddle jets whose differential recurrence packages derivatives of the periodic correction.  Closed forms express the $n$th jet as a harmonic-number constant plus a signed-Stirling combination of $\Psi',\ldots,\Psi^{(n+1)}$.  These jets are then algebraic inputs to the Gaussian contraction and formal logarithm that produce every coefficient of the endpoint expansion.
\end{glossaryentry}
\begin{glossaryentry}{differential-equation}{Differential equation}{ordinary differential equation; differential identity}
A \emph{differential equation} constrains an unknown function through its derivatives.  In the present corpus, ordinary differential identities arise after the functional--differential equations are specialized to intervals on which one translated term vanishes.  For instance, support properties reduce the up-equation to a one-sided identity, and repeated differentiation yields
\[
  F^{(k)}(x)=2^{\binom{k+1}{2}}F(2^k x)
\]
for the signed global extension.  Such formulas are stronger than abstract smoothness: they explicitly transport values to derivatives and are the engine behind dyadic rationality, flatness, and exact derivative maxima.  Endpoint conventions matter.  An identity valid for the signed extension on all of $\mathbb R$ may hold for the bounded CDF only before clamping becomes active.
\end{glossaryentry}
\begin{glossaryentry}{dilation}{Dilation and translation}{scaling; affine change of variable; homothetic copy}
A \emph{dilation} replaces $x$ by $ax$, while a translation replaces it by $x-b$; together they form an affine change of variable.  The Fabius theory is organized around the dyadic dilation $x\mapsto2x$.  Under Fourier transformation, dilation becomes inverse dilation of the frequency variable together with a normalization factor; translation becomes multiplication by a complex exponential.  These rules turn the up-function's differential equation into a product recurrence for its Fourier transform.  In probability, dilation rescales uniform summands and convolution factors.  In exact arithmetic, repeated dyadic dilation moves a rational argument through the binary tree.  The special role of $2$ is forced by the support and normalization conditions rather than chosen merely for convenience.
\end{glossaryentry}
\begin{glossaryentry}{dirac-measure}{Dirac measure}{point mass; delta measure}
The \emph{Dirac measure} $\delta_a$ at a point $a$ assigns mass one to every measurable set containing $a$ and zero otherwise.  Integration against it evaluates a function: $\int f\,d\delta_a=f(a)$.  Finite discrete probability measures built from Dirac masses model the partial random sums used in the probabilistic construction of the up-function.  Their convolutions describe sums of independent discrete coordinates, and weak convergence sends these finite laws to the continuous limiting distribution.  A Dirac measure should not be confused with the Dirac delta as an ordinary function; it is either a measure or, in distribution theory, a generalized function.  Its Fourier transform is the pure phase $e^{-2\pi i a\xi}$.
\end{glossaryentry}
\begin{glossaryentry}{dirichlet-series}{Dirichlet series}{zeta-type series}
A \emph{Dirichlet series} has the form $\sum_{n\ge1}a_n n^{-s}$ for a complex variable $s$.  Its half-plane of absolute convergence often admits meromorphic continuation far beyond the defining region.  The Riemann zeta function is the basic example, $\zeta(s)=\sum_{n\ge1}n^{-s}$ for $\Re s>1$.  In the Mellin analysis of the logarithmic Bose kernel, expanding $\log(1-e^{-x})=-\sum_{n\ge1}e^{-nx}/n$ and integrating termwise produces the factor $\zeta(a+1)$.  Thus the Gamma--zeta product is not an ornamental special-function appearance: it is the exact Mellin image of the exponential Dirichlet expansion.
\end{glossaryentry}
\begin{glossaryentry}{discrete-limit}{Discrete limit}{continuum limit of finite sums; Wolfram \texttt{DiscreteLimit}}
A \emph{discrete limit} is a limit of values indexed by integers, often of finite sums whose mesh tends to zero or whose combinatorial depth tends to infinity.  In the Fabius development, finite $q$-binomial/Thue--Morse expressions converge to the signed global Fabius function.  A typical row depends on a shift parameter even though the limiting value does not.  The proof does not infer this from formal cancellation at finite depth; it reindexes each row as a uniformly bounded Toeplitz average of centered finite splines, proves convergence of the underlying distribution functions, and controls complex shifts by Taylor estimates.  This distinction prevents a common error: equality of limits need not imply termwise equality of the approximants.
\end{glossaryentry}
\begin{glossaryentry}{distribution-generalized}{Distribution in the sense of generalized functions}{Schwartz distribution; generalized function}
A \emph{Schwartz distribution} is a continuous linear functional on the space $C_c^\infty$ of smooth compactly supported test functions.  Locally integrable functions and Radon measures define distributions by integration, while derivatives are defined by duality and therefore exist for objects that are not classically differentiable.  The word ``distribution'' is also used for a probability law; the two meanings are related but not identical.  The Fabius corpus mainly uses probability distributions and classical smooth functions, yet its source literature places compactly supported test functions, Fourier transforms, and Radon measures in the broader distribution-theoretic setting.  Because the up-function is itself a test function, it can serve both as an ordinary function and as a regular distribution.
\end{glossaryentry}
\begin{glossaryentry}{dominated-convergence}{Dominated convergence theorem}{Lebesgue dominated convergence}
The \emph{dominated convergence theorem} states that if measurable functions $f_n$ converge almost everywhere to $f$ and satisfy $|f_n|\le g$ for one integrable function $g$, then $f$ is integrable and $\int f_n\to\int f$.  It is the standard device for interchanging a limit with an integral when uniform convergence is unavailable.  In the Fabius asymptotic layer it justifies limiting procedures for split finite-part integrals and logarithmic kernels.  The difficult step is rarely the theorem itself; it is constructing a domination uniform in the parameter, with separate estimates near zero and at infinity.  Exponential decay controls the large-variable tail, while subtraction of the singular logarithmic term regularizes the small-variable part.
\end{glossaryentry}
\begin{glossaryentry}{double-factorial}{Double factorial}{odd factorial; $(2n+1)!!$}
The \emph{double factorial} multiplies every second positive integer: $(2n+1)!!=1\cdot3\cdot5\cdots(2n+1)$ and $(2n)!!=2\cdot4\cdots2n$.  It satisfies $(2n+1)!!=(2n+1)!/(2^n n!)$.  Odd double factorials occur in Gaussian moments, since $\mathbb E[Z^{2n}]=(2n-1)!!$ for a standard normal variable, and in the rational normalization of the Fabius moment sequence.  Their oddness also makes them transparent to the $2$-adic valuation.  The convention $(-1)!!=1$ corresponds to the empty product and gives the zeroth Gaussian moment.
\end{glossaryentry}
\begin{glossaryentry}{dyadic-denominator}{Dyadic common denominator}{$D_n$; denominator sequence}
The \emph{dyadic common denominator} $D_n$ is the least common multiple of the reduced denominators of the up- or Fabius-function values at the odd grid points of denominator $2^n$ in the relevant fundamental interval.  It measures the simultaneous arithmetic complexity of an entire dyadic row rather than one special value.  Explicit moment formulas give a universal integer divisible by $D_n$, while inverse-dyadic values provide lower divisibility constraints.  The arithmetic paper formulates further patterns for the normalized odd part of $D_n$.  Because equivalent grid presentations can repeat values, the definition uses the reduced finite family appropriate to order $n$.
\end{glossaryentry}
\begin{glossaryentry}{dyadic-grid}{Dyadic grid}{binary mesh; grid of order $n$}
The \emph{dyadic grid of order $n$} consists of points $k/2^n$, usually in a fixed interval.  Refining the grid replaces $n$ by $n+1$ and inserts midpoints.  Dyadic grids are natural for the Fabius function because its equations rescale by two.  Exact tables of values on one grid refine consistently to the next, and finite spline approximants are evaluated on grids whose order is tied to the requested output precision.  The binary numerator records a path through nested dyadic cells, so removing its highest set bit corresponds to one step of the exact evaluator.  Dyadic-grid arguments should be distinguished from statements about arbitrary rationals, whose denominators may contain odd factors and for which the exact dyadic evaluator deliberately returns no value.
\end{glossaryentry}
\begin{glossaryentry}{dyadic-rational}{Dyadic rational}{binary rational; dyadic point}
A \emph{dyadic rational} is a rational number of the form $a/2^n$ with $a\in\mathbb Z$ and $n\in\mathbb N$.  It has a finite binary expansion, although its representation is not unique because $a/2^n=(2a)/2^{n+1}$.  The Fabius and up-functions take rational values at every dyadic point.  The development proves this both by an explicit finite Thue--Morse/moment formula and by a terminating binary recursion.  Correctness requires \emph{representation independence}: changing to an unreduced dyadic presentation must not change the computed value.  Dyadic points also have exceptional local structure: after enough differentiations, the scaled argument leaves the support, so sufficiently high derivatives vanish there.
\end{glossaryentry}
\begin{glossaryentry}{density}{Probability density}{density function; Radon--Nikodym density}
A \emph{probability density} $p$ with respect to Lebesgue measure is a nonnegative integrable function satisfying $\int_{\mathbb R}p(x)\,dx=1$.  It represents a probability law by $\mathbb P(X\in A)=\int_Ap(x)\,dx$, and its cumulative distribution function is $F(x)=\int_{-\infty}^x p(t)\,dt$.  In the Fabius normalization, the bounded function is a cumulative distribution function, whereas a translated and rescaled copy of Rvachev's up-function supplies its derivative: globally, $F'(x)=2\Up(2x-1)$.  The density is $C^\infty$, compactly supported, symmetric about $1/2$, and strictly positive in the interior.  This density/CDF distinction is crucial because the literature sometimes calls either object ``the Fabius function,'' leading to apparent discrepancies in support, normalization, and functional equations.
\end{glossaryentry}
\begin{glossaryentry}{derivative-bound}{Uniform derivative bound}{bounded derivatives; derivative supremum}
A \emph{uniform derivative bound} controls $|f^{(k)}(x)|$ by a constant depending only on the order $k$, uniformly in $x$.  Repeated use of the Fabius dilation law gives the exact scale $2^{\binom{k+1}{2}}$.  For the signed global Fabius function and for Rvachev's up-function,
\[
  |f^{(k)}(x)|\le 2^{\binom{k+1}{2}},
\]
and the repository proves that this bound is attained, hence is the exact supremum rather than a convenient estimate.  The rapid growth of these norms is compatible with $C^\infty$ smoothness but incompatible with analyticity on the active interval: analytic derivatives would obey a factorial-geometric Cauchy bound locally.  Sharp derivative bounds also feed into Taylor remainders, Lipschitz estimates, and effective computability.
\end{glossaryentry}
\clearpage
\hypertarget{letter:E}{}
\section*{\color{linkblue}E}
\addcontentsline{toc}{section}{E}
\markboth{E}{E}
\alphabetbar
\begin{glossaryentry}{effective-flatness}{Effective flatness}{quantitative infinite-order vanishing}
A smooth function is \emph{flat} at the origin when every derivative there vanishes, equivalently $f(x)=o(x^n)$ for every fixed $n$ under mild regularity.  An \emph{effective flatness bound} supplies explicit constants and a stated region.  The Fabius development proves, whenever $x\ge0$ and $2^n x\le1$,
\[
 F(x)\le 2^{\binom{n+1}{2}}x^n,
 \qquad
 F(x)\le \frac{2^{\binom{n+1}{2}}}{n!}x^n.
\]
The first follows from an iterated mean-value estimate; the sharper factorial version iterates the fundamental theorem of calculus.  Such estimates quantify how rapidly $F$ disappears near zero, yield computable stopping rules, and transfer to both ends of the support of the up-function.
\end{glossaryentry}
\begin{glossaryentry}{effective-uniform-continuity}{Effective uniform continuity}{recursive modulus of continuity; computable modulus}
A function is \emph{effectively uniformly continuous} when there is an explicit computable modulus converting an output precision into a sufficient input precision.  In a dyadic convention one may require a recursive function $d$ such that $|x-y|\le2^{-d(n)}$ implies $|f(x)-f(y)|\le2^{-n}$.  The Fabius development uses the global Lipschitz constant $2$ to obtain a simple primitive-recursive modulus, presented in the source's reciprocal precision convention as $d(n)=2n$.  The exact formula is less important than uniformity and recursiveness: one algorithm must work for every input and precision.  Together with sequential computability, this is one of the two clauses in the Grzegorczyk-style definition of a computable real function used in the repository.
\end{glossaryentry}
\begin{glossaryentry}{empty-conventions}{Empty sum, empty product, and $0^0$ conventions}{boundary conventions; neutral-element convention}
Finite formulas often remain valid at order zero only after the neutral-element conventions are stated explicitly: an empty sum is $0$, an empty product is $1$, the unique composition of $0$ is the empty composition, and natural-number exponentiation takes $0^0=1$.  These are not cosmetic choices.  They determine endpoint terms in $q$-binomial identities, the initial value in the composition formula for $F(1)$, and the $m=0$ row in global binary-reduction series.  In formal proofs the conventions are part of the definitions, so an omitted boundary term can change a theorem at exactly one endpoint.  Several corrections in the corpus amount to restoring the neutral term that an informal one-indexed sum silently discarded.
\end{glossaryentry}
\begin{glossaryentry}{endpoint-asymptotic}{Endpoint asymptotic}{small-argument asymptotic; boundary-layer asymptotic}
An \emph{endpoint asymptotic} describes a function as its argument approaches the boundary of its active domain.  For the bounded Fabius function the central regime is $x\to0^+$; symmetry transfers the same result to $1-F(1-x)$ as $x\to0^+$.  Because $F$ is flat to every algebraic order, no nonzero power of $x$ is a useful leading term.  Instead, $\log F(x)$ has a quadratic expression in a lower-Lambert phase, a nonconstant periodic correction, and an all-orders inverse-phase expansion.  Endpoint asymptotics are much more sensitive than numerical fits on compact subintervals: a tiny logarithmic or periodic omission can destroy asymptotic equivalence even when plots look indistinguishable.
\end{glossaryentry}
\begin{glossaryentry}{entire-function}{Entire function}{entire complex function}
An \emph{entire function} is holomorphic on all of $\mathbb C$.  The Fourier transform of a compactly supported $C^\infty$ function extends to an entire function, with growth controlled by the support.  For Rvachev's up-function the transform is explicitly
\[
 \widehat{\Up}(z)=\prod_{j\ge0}\frac{\sin(\pi z/2^j)}{\pi z/2^j},
\]
an infinite product converging uniformly on compact subsets.  Its zeros and power-series coefficients encode arithmetic information about moments.  Entirety of the transform contrasts with nowhere analyticity of the original bump on its support.  Fourier inversion recovers the bump from the restriction of this entire function to the real axis.
\end{glossaryentry}
\begin{glossaryentry}{euler-maclaurin}{Euler--Maclaurin summation formula}{sum--integral expansion}
The \emph{Euler--Maclaurin formula} compares a sum with an integral and supplies boundary corrections involving Bernoulli numbers:
\[
 \sum_{k=m}^n f(k)=\int_m^n f(x)\,dx+\frac{f(m)+f(n)}2
 +\sum_{r=1}^{p}\frac{B_{2r}}{(2r)!}
   \bigl(f^{(2r-1)}(n)-f^{(2r-1)}(m)\bigr)+R_p.
\]
It is a classical source of asymptotic expansions for sums over a smooth scale.  In dyadically self-similar problems it captures the smooth drift but can miss the exact periodic fluctuation unless the discrete scale is retained, for example through Mellin analysis.  The Fabius asymptotic discussion uses this contrast to explain why a constant correction is insufficient.
\end{glossaryentry}
\begin{glossaryentry}{euler-mascheroni}{Euler--Mascheroni constant}{Euler's constant; $\gamma$}
The \emph{Euler--Mascheroni constant} is
\[
 \gamma=\lim_{n\to\infty}\left(\sum_{k=1}^n\frac1k-\log n\right).
\]
It is the constant term in the Laurent expansion $\zeta(s)=1/(s-1)+\gamma+\cdots$ and appears in the expansion of the Gamma function at the origin.  In the Fabius Mellin calculation, the finite part of $-\Gamma(a)\zeta(a+1)$ at $a=0$ combines $\gamma$, the first Stieltjes constant, and powers of $\pi$.  Euler's constant also enters the closed saddle-jet constants through harmonic-number limits.  Its role is therefore structural: it records the mismatch between a harmonic sum and its logarithmic continuum approximation.
\end{glossaryentry}
\begin{glossaryentry}{expectation-variance}{Expectation and variance}{mean of a random variable; second central moment}
The \emph{expectation} of an integrable random variable $X$ is $\mathbb E[X]=\int X\,d\mathbb P$.  Its variance is $\operatorname{Var}(X)=\mathbb E[(X-\mathbb E X)^2]$, measuring quadratic spread.  For independent summands, expectations and variances add, with variances scaled by the square of the coefficients.  The weighted-uniform representation of the Fabius law therefore gives its mean, symmetry center, and variance directly.  Raw moments $\mathbb E[X^n]$ and centered moments play different roles: raw moments feed generating functions and dyadic formulas, while centered moments expose symmetry and eliminate odd terms.  In the saddle expansion, Gaussian moments appear after a separate local normalization and should not be confused with moments of the Fabius distribution itself.
\end{glossaryentry}
\begin{glossaryentry}{exponential-generating-function}{Exponential generating function}{EGF}
For a sequence $(a_n)$, its \emph{exponential generating function} is $A(z)=\sum_{n\ge0}a_n z^n/n!$.  The factorial denominator makes differentiation, binomial convolution, and moment identities especially natural.  The half-moment sequence in the Fabius arithmetic satisfies a functional equation through its EGF, while moments appear in the even power series of the Fourier transform.  Formal exponentiation and formal logarithms of EGFs generate Bell-polynomial and cumulant recurrences.  An EGF may be treated analytically when the series converges or formally when only coefficient identities are needed; the corpus carefully separates those two uses and later proves analytic convergence where required.
\end{glossaryentry}
\begin{glossaryentry}{exponential-type}{Exponential type}{entire function of exponential type}
An entire function $G$ is of \emph{exponential type} at most $A$ if roughly $|G(z)|\le C_\varepsilon e^{(A+\varepsilon)|z|}$ for every $\varepsilon>0$, or in directional refinements if its growth is controlled by $e^{A|\operatorname{Im}z|}$.  The Paley--Wiener theorem relates such growth to compact support of the inverse Fourier transform.  Since the up-function is supported on $[-1,1]$, its Fourier transform has corresponding exponential-type bounds.  Rapid decay on the real axis and exponential growth off it coexist: the first reflects smoothness, the second records the finite support radius.  This is why an entire Fourier transform can still be absolutely integrable along the real axis.
\end{glossaryentry}
\begin{glossaryentry}{extended-nonnegative-reals}{Extended nonnegative reals}{$\mathbb R_{\ge0}\cup\{\infty\}$; $\mathbb R_{\ge0\infty}$; \texttt{ENNReal}}
The \emph{extended nonnegative real numbers} adjoin $+\infty$ to $[0,\infty)$.  Measure theory naturally takes values in this ordered space because a measurable set may have infinite measure and because countable sums of nonnegative quantities always have a supremum there.  Lean denotes the type by \texttt{\ensuremath{\mathbb R_{\ge0\infty}}} or \texttt{ENNReal}.  Probability measures have total mass one, so their final values are finite, but many intermediate statements are most canonical in the extended type.  The Fabius probability API exposes both native measure-valued identities and real-valued corollaries obtained after proving finiteness and coercing to $\mathbb R$.
\end{glossaryentry}
\clearpage
\hypertarget{letter:F}{}
\section*{\color{linkblue}F}
\addcontentsline{toc}{section}{F}
\markboth{F}{F}
\alphabetbar
\begin{glossaryentry}{fabius-function}{Bounded Fabius function}{canonical Fabius CDF; \texttt{fabiusReal}; bounded Fabius function}
The \emph{bounded Fabius function} $F:\mathbb R\to[0,1]$ is the unique smooth cumulative distribution function that equals $0$ on $(-\infty,0]$, equals $1$ on $[1,\infty)$, satisfies $F(x)+F(1-x)=1$, and obeys $F'(x)=2F(2x)$ on the left half of the unit interval.  Equivalently, it is the distribution function of $\sum_{k\ge1}2^{-k}U_k$ with independent $U_k$ uniform on $[0,1]$.  It is strictly increasing on $[0,1]$, convex up to $1/2$, concave after $1/2$, flat at both endpoints, and nowhere analytic on $[0,1]$.  Every dyadic value is rational and exactly computable.  This bounded object must be distinguished from the signed global extension, which agrees with it on $[0,1]$ but is not clamped outside that interval.
\end{glossaryentry}
\begin{glossaryentry}{factorial}{Factorial}{$n!$}
The \emph{factorial} $n!=1\cdot2\cdots n$, with $0!=1$, counts permutations and is the natural normalization for derivatives, Taylor coefficients, exponential generating functions, and simplex volumes.  In the Fabius theory factorials enter exact dyadic values, moment normalizations, Taylor recurrences, and saddle coefficients.  Their $2$-adic valuation is given by Legendre's formula $v_2(n!)=n-s_2(n)$, where $s_2(n)$ is the binary digit sum.  This identity links factorial denominators to Thue--Morse parity and is one reason binary weight appears in exact valuation formulas.  Factorial growth also separates analyticity from mere smoothness through Cauchy derivative estimates.
\end{glossaryentry}
\begin{glossaryentry}{filter-convergence}{Filter convergence}{tendsto; neighborhood filter; asymptotic filter}
A \emph{filter} abstracts the sets that count as ``eventually'' in a limiting process.  A function $f$ tends to $y$ along a filter $\mathcal F$ when inverse images of every neighborhood of $y$ belong to $\mathcal F$.  Sequences use the cofinite filter on $\mathbb N$; $x\to0^+$ uses a punctured right-neighborhood filter; almost-everywhere statements use the measure filter.  Lean's asymptotic relations, continuity, series, and weak convergence are all formulated through filters.  This uniform language prevents the limiting direction or domain restriction from being implicit, which is especially important for the Fabius left endpoint, lower-Lambert phase, and contour parameters.
\end{glossaryentry}
\begin{glossaryentry}{finite-difference}{Finite difference}{discrete derivative; forward difference}
The \emph{forward difference} of a function on an arithmetic grid is $\Delta f(x)=f(x+h)-f(x)$, and higher differences are defined iteratively.  For a polynomial of degree less than $n$, the $n$th finite difference vanishes.  Alternating binomial sums are finite differences in disguise, and Thue--Morse signed sums provide stronger cancellations adapted to dyadic blocks.  In the Fabius corpus, these cancellations annihilate low-degree polynomial terms, isolate leading coefficients, and make certain translated expressions independent of their common shift.  Finite-difference reasoning is the discrete counterpart of differentiation and underlies polynomial reproduction, prefix-sum identities, and the $q$-binomial formulas.
\end{glossaryentry}
\begin{glossaryentry}{finite-part}{Finite part of a singular expression}{Hadamard finite part; constant term at a pole; regularized value}
When an integral or meromorphic expression diverges in a controlled way, its \emph{finite part} is the constant remaining after explicitly subtracting the singular terms.  If $M(a)=c_{-2}a^{-2}+c_{-1}a^{-1}+c_0+O(a)$, the finite part at $a=0$ is $c_0$.  The Mellin transform of $\log(1-e^{-x})$ equals $-\Gamma(a)\zeta(a+1)$ for $a>0$ and has a double pole at the origin.  The Fabius development realizes its finite part both algebraically from Laurent expansions and analytically as a convergent split integral: near zero it removes the $\log x$ singularity, and near infinity exponential decay makes the remaining quotient integrable.  This gives the constant needed for the mean of the periodic correction.
\end{glossaryentry}
\begin{glossaryentry}{fixed-point}{Fixed point}{invariant point; fixed-point equation}
A \emph{fixed point} of an operator $T$ is an object $f$ satisfying $T(f)=f$.  Many self-similar functions are most naturally constructed this way rather than by guessing a closed formula.  The canonical bounded Fabius function is obtained as the unique fixed point of an integral operator on a complete class of continuous symmetric unit-interval-valued functions.  Contractivity supplies existence, uniqueness, and convergence of iterates.  Differentiating the integral fixed-point equation then recovers the functional--differential law and bootstraps continuity to $C^\infty$ regularity.  This order is logically useful: existence is established before one assumes the derivatives that the final object will possess.
\end{glossaryentry}
\begin{glossaryentry}{flat-function}{Flat function}{infinitely flat; vanishing to infinite order}
A smooth function $f$ is \emph{flat at $x_0$} when $f^{(n)}(x_0)=0$ for every $n\ge0$.  Equivalently, its Taylor series at $x_0$ is identically zero.  Flatness does not imply local vanishing: the classical example $e^{-1/x^2}$ for $x>0$ is positive arbitrarily close to zero.  The Fabius function is flat at $0$ and $1$, and Rvachev's up-function is flat at the endpoints of its support.  Consequently these functions cannot be analytic there.  Their decay is even more unusual than the standard exponential bump: $F(x)$ is smaller than every power of $x$ but larger than $e^{-c/x}$ for every fixed $c>0$, with the precise intermediate scale governed by a quadratic logarithmic law.
\end{glossaryentry}
\begin{glossaryentry}{floor-function}{Floor and ceiling functions}{greatest-integer function; least-integer function}
The \emph{floor} $\lfloor x\rfloor$ is the greatest integer not exceeding $x$, and the ceiling $\lceil x\rceil$ is the least integer not smaller than $x$.  They convert real arguments into finite summation ranges and encode dyadic block boundaries.  In the global binary-reduction and discrete-limit formulas, expressions such as $\lfloor2^{m-1}x\rfloor$ or a truncated natural floor determine how many translated compact-support terms can contribute.  Endpoint identities depend on whether a bound is inclusive, so replacing a successor of a floor by an apparently equivalent range can create an empty-sum error.  Lean's typed floor operations make the nonnegativity and conversion conventions explicit.
\end{glossaryentry}
\begin{glossaryentry}{formal-power-series}{Formal power series}{FPS; coefficientwise series}
A \emph{formal power series} $A(X)=\sum_{n\ge0}a_nX^n$ is an algebraic object whose coefficients are manipulated without asking whether the series converges for any numerical $X$.  Addition and multiplication are coefficientwise and Cauchy-product operations; composition, exponentiation, and logarithm are defined under suitable constant-term conditions.  The Fabius saddle calculation uses formal series to expand a normalized exponent, exponentiate it, contract against Gaussian moments, and take a formal logarithm.  This isolates finite algebra at each asymptotic order from the analytic theorem that controls the remainder.  Confusing formal identity with analytic convergence would leave a gap; the corpus supplies both layers separately.
\end{glossaryentry}
\begin{glossaryentry}{fourier-coefficient}{Fourier coefficient}{harmonic coefficient; Fourier mode}
For a one-periodic integrable function $f$, its $k$th \emph{Fourier coefficient} is
\[
 \widehat f(k)=\int_0^1 f(t)e^{-2\pi i kt}\,dt.
\]
The coefficients describe the amplitudes and phases of the integer-frequency modes.  Smoothness forces rapid decay of $\widehat f(k)$, while analyticity in a strip gives exponential decay.  The centered periodic correction in the Fabius asymptotic has explicit nonzero coefficients
\[
 \widehat\Psi(k)=-\frac{\Gamma(-\chi_k)\zeta(1-\chi_k)}{\log2},
 \qquad \chi_k=\frac{2\pi i k}{\log2},\quad k\ne0.
\]
Their nonvanishing proves that the correction is genuinely nonconstant.  Differentiation multiplies the $k$th coefficient by $(2\pi i k)^m$, explaining amplification in higher saddle coefficients.
\end{glossaryentry}
\begin{glossaryentry}{fourier-inversion}{Fourier inversion}{inversion formula}
\emph{Fourier inversion} reconstructs a function from its Fourier transform.  With the convention $\widehat f(\xi)=\int f(x)e^{-2\pi i x\xi}\,dx$, a sufficiently regular function satisfies
\[
 f(x)=\int_{\mathbb R}\widehat f(\xi)e^{2\pi i x\xi}\,d\xi.
\]
Schwartz functions satisfy the theorem pointwise.  In the original construction of the up-function, the infinite sinc product first defines an entire rapidly decreasing transform; inversion then defines the compactly supported smooth function.  Injectivity of the Fourier transform also proves uniqueness once the transform is forced by the functional--differential equation and normalization.  Exact normalization of the $2\pi$ factors is essential because it changes the scale of every sinc factor and exponential phase.
\end{glossaryentry}
\begin{glossaryentry}{fourier-legendre-series}{Fourier--Legendre series}{Legendre expansion; orthogonal-polynomial expansion}
A \emph{Fourier--Legendre series} expands a function on $[-1,1]$ in the orthogonal basis of Legendre polynomials:
\[
 f(x)\sim\sum_{n\ge0}a_nP_n(x),\qquad
 a_n=\frac{2n+1}{2}\int_{-1}^1f(x)P_n(x)\,dx.
\]
Evenness of Rvachev's up-function eliminates all odd coefficients.  The repository evaluates each even coefficient as an exact finite sum of dyadic Fabius values and proves absolute, uniform convergence on the closed interval, endpoints included.  Its finite even partial sums are not merely approximations: they are the unique least-squares minimizers among all polynomials of degree at most one higher than their visible even degree, because the missing odd component is orthogonal to the target.
\end{glossaryentry}
\begin{glossaryentry}{fourier-series}{Fourier series}{trigonometric expansion}
A \emph{Fourier series} represents a periodic function as $\sum_{k\in\mathbb Z}\widehat f(k)e^{2\pi i kx}$.  Convergence may be pointwise, uniform, absolute, or in $L^2$; these are distinct claims.  The up-function admits a cosine series on its support whose coefficients are samples of its Fourier transform at half-integers.  The asymptotic periodic correction has an absolutely convergent Gamma--zeta Fourier series.  Absolute convergence permits rearrangement and yields uniform convergence; smoothness then permits termwise differentiation when the differentiated coefficient series also converges.  Fourier series should not be confused with the continuous Fourier transform, although Poisson summation links them by sampling one side and periodizing the other.
\end{glossaryentry}
\begin{glossaryentry}{fourier-transform}{Fourier transform}{frequency transform; characteristic transform}
The \emph{Fourier transform} converts a function of position into a function of frequency.  Under the $e^{-2\pi i x\xi}$ convention, differentiation becomes multiplication by $2\pi i\xi$, translation becomes a phase factor, dilation rescales the frequency, and convolution becomes multiplication.  These rules turn the up-function's differential equation into the recurrence
\[
 \widehat\Up(z)=\sinc(\pi z)\widehat\Up(z/2),
\]
whose iteration gives the infinite sinc product.  The transform is simultaneously a characteristic function of the associated probability law, an entire function encoding all moments, and the central uniqueness tool.  The normalization convention must be stated because alternative definitions distribute the $2\pi$ factors differently.
\end{glossaryentry}
\begin{glossaryentry}{fundamental-theorem-calculus}{Fundamental theorem of calculus}{FTC; integral--derivative theorem}
The \emph{fundamental theorem of calculus} links differentiation and integration.  If $f'$ is continuous, then $f(b)-f(a)=\int_a^b f'(t)\,dt$; conversely, an integral of a continuous function differentiates to its integrand.  Iterating the theorem with vanishing endpoint derivatives gives the integral Taylor formula
\[
 f(x)=\frac1{(n-1)!}\int_0^x(x-t)^{n-1}f^{(n)}(t)\,dt.
\]
The sharp Fabius flatness estimate uses this formula recursively, preserving the factorial that a pointwise mean-value bound loses.  Integral fixed-point definitions also use the converse direction to obtain the functional--differential equation after existence has been established.
\end{glossaryentry}
\clearpage
\hypertarget{letter:G}{}
\section*{\color{linkblue}G}
\addcontentsline{toc}{section}{G}
\markboth{G}{G}
\alphabetbar
\begin{glossaryentry}{gamma-function}{Gamma function}{Euler gamma function; $\Gamma$}
The \emph{Gamma function} extends the factorial to complex arguments.  For $\Re s>0$,
\[
 \Gamma(s)=\int_0^\infty x^{s-1}e^{-x}\,dx,
\]
and it satisfies $\Gamma(s+1)=s\Gamma(s)$, hence $\Gamma(n+1)=n!$.  It extends meromorphically to $\mathbb C$ with simple poles at the nonpositive integers and no zeros.  Mellin transforms of exponentially decaying kernels naturally produce Gamma factors.  In the Fabius endpoint theory, the Fourier coefficients of the periodic correction contain $\Gamma(-2\pi i k/\log2)$, and its nonvanishing helps prove that every nonzero Fourier mode is present.  Stirling's formula for $\Gamma$ also gives rapid decay of these coefficients as $|k|\to\infty$.
\end{glossaryentry}
\begin{glossaryentry}{gamma-zeta-coefficient}{Gamma--zeta Fourier coefficient}{Gamma--zeta mode; Mellin--Fourier coefficient}
A \emph{Gamma--zeta Fourier coefficient} is the explicit product of special-function values that arises when a Mellin transform is sampled at imaginary frequencies.  For the centered Fabius periodic correction,
\[
 \widehat\Psi(k)=-\frac{\Gamma(-\chi_k)\zeta(1-\chi_k)}{\log2},
 \qquad \chi_k=\frac{2\pi i k}{\log2},\quad k\ne0.
\]
The Gamma factor comes from the Mellin transform of an exponential, while the zeta factor comes from summing the logarithmic Bose series over positive integers.  Sampling on the lattice $2\pi i k/\log2$ converts multiplicative dyadic periodicity into ordinary period one.  Since Gamma has no zeros and zeta has no zeros on $\Re s=1$, every nonzero mode survives; the endpoint oscillation is therefore intrinsic, not a numerical artifact.
\end{glossaryentry}
\begin{glossaryentry}{gaussian-binomial}{Gaussian binomial coefficient}{$q$-binomial coefficient; Gaussian coefficient}
The \emph{Gaussian binomial coefficient} is the polynomial
\[
 \qbinom{n}{k}{q}=\frac{(1-q^n)(1-q^{n-1})\cdots(1-q^{n-k+1})}
 {(1-q)(1-q^2)\cdots(1-q^k)}
\]
for $0\le k\le n$, with value zero outside that range.  It counts $k$-dimensional subspaces of $\mathbb F_q^n$ when $q$ is a prime power, but it is meaningful as a polynomial or at any nonsingular scalar base.  The Fabius identities use the fixed base $q=1/2$.  At this base a Gaussian-binomial row is the leading-coefficient functional for Lagrange interpolation at geometric nodes, and it combines with Thue--Morse cancellation to produce exact dyadic values.  The term ``Gaussian'' here predates and is unrelated to the normal probability distribution.
\end{glossaryentry}
\begin{glossaryentry}{gaussian-contraction}{Gaussian contraction}{Wick contraction; Gaussian expectation operator}
A \emph{Gaussian contraction} applies expectation with respect to a standard Gaussian variable to a polynomial or formal series.  Because $\mathbb E[Z^m]=0$ for odd $m$ and $(m-1)!!$ for even $m$, it deletes odd monomials and replaces even powers by double factorials.  In saddle-point asymptotics, the contour is rescaled so that the leading exponent is $-z^2/2$; integrating correction polynomials against this Gaussian kernel is exactly a contraction.  The Fabius all-orders construction first forms a formal exponential of the higher saddle terms and then contracts each coefficient.  This produces the ``mass coefficients'' before a formal logarithm converts them into coefficients for $\log F$.
\end{glossaryentry}
\begin{glossaryentry}{gaussian-moment}{Gaussian moment}{normal moment; Wick moment}
For a standard normal random variable $Z\sim N(0,1)$, the \emph{Gaussian moments} are
\[
 \mathbb E[Z^{2m+1}]=0,
 \qquad
 \mathbb E[Z^{2m}]=(2m-1)!!=\frac{(2m)!}{2^m m!}.
\]
Odd moments vanish by symmetry, and even moments follow from the moment-generating function $e^{t^2/2}$ or integration by parts.  These values are the universal coefficients in one-dimensional saddle expansions.  They should be distinguished from the Fabius moment sequence: Gaussian moments describe the local quadratic approximation at a saddle, whereas Fabius moments integrate powers against the up-function itself.  Their appearance explains why only finitely many exponent jets contribute to each asymptotic order.
\end{glossaryentry}
\begin{glossaryentry}{gaussian-rational}{Gaussian rational}{rational complex number; $\mathbb Q(i)$}
A \emph{Gaussian rational} is a complex number $a+bi$ with $a,b\in\mathbb Q$, that is, an element of the quadratic field $\mathbb Q(i)$.  Gaussian rationals provide an exact coefficient domain for complex specializations without introducing arbitrary real or complex approximation.  The Fabius translated $q$-binomial and discrete-limit APIs include Gaussian-rational shift parameters as convenient exact cases between rational and fully complex statements.  The name is unrelated to Gaussian probability; it refers to the Gaussian integers $\mathbb Z[i]$ and their fraction field.  Arithmetic equality is decidable and exact, while analytic theorems use the canonical embedding into $\mathbb C$.
\end{glossaryentry}
\begin{glossaryentry}{generating-function}{Generating function}{ordinary generating function; sequence transform}
A \emph{generating function} encodes a sequence as coefficients of a series.  The ordinary generating function is $A(z)=\sum_{n\ge0}a_nz^n$, while the exponential generating function uses $a_nz^n/n!$.  Recurrences become algebraic or functional equations for $A$, and products encode convolution of sequences.  In the Fabius corpus, generating functions organize moment and half-moment recurrences, Thue--Morse prefix sums, parity-power identities, and the recurrence sequence attached to inverse dyadic values.  One must specify whether a series is formal or analytic: coefficient manipulations may be valid without convergence, whereas evaluation at a numerical $z$ requires a radius-of-convergence argument.
\end{glossaryentry}
\begin{glossaryentry}{geometric-nodes}{Geometric interpolation nodes}{geometric progression of nodes}
\emph{Geometric nodes} are interpolation points whose successive ratios are constant, for example $x_j=-2^j$.  Their pairwise differences factor into powers of two and factors $1-q^j$, so Lagrange weights are naturally expressed through $q$-Pochhammer symbols and Gaussian binomial coefficients.  In the Fabius half-base $q$-calculus, evaluation at these nodes turns a translated-power sum into a leading-coefficient extractor.  This is the structural reason niche $q$-special functions arise in exact dyadic formulas: they are compact notation for the barycentric denominators of a geometric interpolation problem, not an externally imposed analogy.
\end{glossaryentry}
\begin{glossaryentry}{germ}{Germ of a function}{local germ; equality near a point}
The \emph{germ} of a function at a point records only its behavior in some unspecified neighborhood of that point.  Two functions have the same germ if they agree on a neighborhood, even if they differ globally.  Germ language is useful when transferring local derivatives or analyticity between the bounded Fabius function and its signed extension.  Below the clamping point they have the same germ, so every iterated derivative identity valid for the signed extension transfers locally to the bounded function.  At and beyond the endpoint this ceases to be automatic because the bounded function becomes locally constant.  A germ is stronger than agreement at a point or agreement of finitely many derivatives.
\end{glossaryentry}
\begin{glossaryentry}{grzegorczyk-computability}{Grzegorczyk computability of a real function}{computable real function in the Grzegorczyk sense}
A real function is \emph{computable in the Grzegorczyk sense} when it maps computable real sequences uniformly to computable real sequences and possesses an effective modulus of uniform continuity on its domain.  Real inputs are represented by fast rational or dyadic names rather than exact bit strings.  The Fabius development constructs a primitive-recursive centered-spline evaluator, rounds it to a requested dyadic precision, and proves a uniform error estimate.  This establishes sequential computability.  A recursive modulus derived from the global Lipschitz bound supplies the second clause.  The theorem concerns the function as an operation on represented reals, not merely the existence of an algorithm for special arguments such as dyadic rationals.
\end{glossaryentry}
\clearpage
\hypertarget{letter:H}{}
\section*{\color{linkblue}H}
\addcontentsline{toc}{section}{H}
\markboth{H}{H}
\alphabetbar
\begin{glossaryentry}{half-moment}{Half moment}{one-sided moment; $d_n$}
A \emph{half moment} integrates a power over one half of a symmetric density's support.  For Rvachev's up-function the rational sequence customarily denoted $d_n$ is tied to
\[
 \int_0^1 t^{n-1}\Up(t)\,dt
\]

after a simple normalization.  These values also determine $F(2^{-n})$ exactly and satisfy triangular recurrences, Bernoulli recurrences, product identities, and precise denominator formulas.  Full odd moments vanish by evenness, so half moments retain information that full moments discard.  The normalized sequence $d_n/n!$ is the Fabius recurrence sequence studied in the asymptotic layer.
\end{glossaryentry}
\begin{glossaryentry}{half-endpoint-indicator}{Half-weighted interval indicator}{endpoint-corrected indicator; half-endpoint convention}
A \emph{half-weighted interval indicator} equals one in the interior of a closed interval, one half at each endpoint, and zero outside.  It is useful when adjacent intervals share endpoints: ordinary closed-interval indicators would double-count boundary points, while half weights preserve exact normalization under summation.  In the finite convolution approximants to the up-function, this convention corrects a printed source formula and keeps the central value and partition identities consistent.  Measure-theoretically the endpoints have Lebesgue measure zero, so the choice does not affect integrals; pointwise convergence and exact endpoint values, however, do depend on it.
\end{glossaryentry}
\begin{glossaryentry}{harmonic-number}{Harmonic number}{$H_n$}
The \emph{$n$th harmonic number} is $H_n=\sum_{k=1}^n1/k$, with $H_0=0$.  Harmonic numbers satisfy $H_n=\log n+\gamma+O(1/n)$ and occur whenever logarithmic derivatives of factorial-like products are expanded.  In the Fabius all-orders endpoint analysis, the closed solution of the saddle-jet recurrence contains harmonic-number constants together with signed Stirling coefficients acting on derivatives of the periodic correction.  Their appearance reflects repeated differentiation of a reciprocal or logarithmic factor in the moving saddle equation.  Generalized harmonic numbers $H_n^{(r)}=\sum_{k=1}^n k^{-r}$ can similarly occur in higher logarithmic expansions.
\end{glossaryentry}
\begin{glossaryentry}{has-sum-tsum}{\texorpdfstring{\texttt{HasSum}}{HasSum} and infinite sum}{summability witness; \texttt{tsum}; topological sum}
In Lean, \texttt{HasSum f a} states that the finite partial sums of a family $f$ converge to $a$ in the relevant topology.  The expression \texttt{tsum f} chooses the sum when $f$ is summable and a default value otherwise.  A theorem of the form \texttt{HasSum} therefore retains the convergence witness, while an equality with \texttt{tsum} is a convenient numeric corollary.  The Fabius naming scheme deliberately distinguishes these forms for Fourier products, translate series, Legendre expansions, and binary-reduction formulas.  Mathematically, this separation prevents an equality from hiding the nontrivial proof that the infinite sum exists and has the stated mode of convergence.
\end{glossaryentry}
\begin{glossaryentry}{horner-method}{Horner's method}{nested polynomial evaluation}
\emph{Horner's method} evaluates a polynomial
\[
 a_0+a_1x+\cdots+a_nx^n
\]

after rewriting it as $a_0+x(a_1+x(a_2+\cdots+x a_n))$.  It uses $n$ multiplications and $n$ additions, minimizes intermediate powers, and works over any semiring.  The exact dyadic Fabius evaluator uses Horner form for the Taylor polynomial generated at each highest-set-bit reduction step.  Since the coefficients and argument are rational, evaluation is exact; efficiency concerns the growth and number of rational operations rather than floating-point rounding.  Horner form is also well suited to formal verification because the recursive syntax mirrors the evaluation proof.
\end{glossaryentry}
\clearpage
\hypertarget{letter:I}{}
\section*{\color{linkblue}I}
\addcontentsline{toc}{section}{I}
\markboth{I}{I}
\alphabetbar
\begin{glossaryentry}{intervals}{Closed, open, and half-open intervals}{$Icc$, $Ioo$, $Ioc$, $Ico$}
For real $a<b$, the closed interval is $[a,b]$, the open interval is $(a,b)$, and the half-open intervals are $(a,b]$ and $[a,b)$.  Mathlib names these sets $\mathrm{Icc}$, $\mathrm{Ioo}$, $\mathrm{Ioc}$, and $\mathrm{Ico}$.  Endpoint choice matters in this corpus.  The pointwise support of the up-function is the open interval $(-1,1)$, while its topological support is the closed interval $[-1,1]$.  Half-open cells avoid double-counting shared boundaries in step approximants.  Likewise, a binary-reduction identity valid on $[0,1)$ may fail at $1$ by a single missing zeroth term.  These are not cosmetic distinctions: formalization forces every endpoint convention to be explicit.
\end{glossaryentry}
\begin{glossaryentry}{improper-integral}{Improper integral}{integral over an unbounded or singular interval}
An \emph{improper integral} is defined as a limit when an integration interval is unbounded or the integrand is singular at an endpoint.  For example,
\[
 \int_0^\infty f(x)\,dx=\lim_{A\to\infty}\int_0^A f(x)\,dx
\]
when the limit exists.  Absolute integrability is stronger and usually permits safer rearrangements.  Mellin and Gamma integrals in the Fabius development are improper at both zero and infinity.  Near zero, powers and logarithms determine convergence; at infinity, exponential decay dominates.  Finite-part integrals are not ordinary improper integrals until the singular terms have been subtracted so that each resulting piece is genuinely integrable.
\end{glossaryentry}
\begin{glossaryentry}{indicator-function}{Indicator function}{characteristic function of a set; set indicator}
For a set $A$, its \emph{indicator function} $\mathbf1_A$ is one on $A$ and zero outside.  It converts set operations into algebra and represents uniform densities on intervals after division by interval length.  Finite convolutions of interval indicators produce compactly supported splines, which converge to the up-function.  The phrase ``characteristic function'' is ambiguous: in probability it usually means the Fourier transform of a law, while in elementary analysis it often means an indicator.  The glossary follows the probability convention for characteristic functions and uses ``indicator'' for sets.  Endpoint values of an indicator are a pointwise convention and can matter even though they do not affect its Lebesgue integral.
\end{glossaryentry}
\begin{glossaryentry}{infinite-convolution}{Infinite convolution}{convolution product of measures; infinite sum law}
An \emph{infinite convolution} is the limiting probability law obtained by convolving an increasing family of measures, typically the laws of partial sums of independent random variables.  If $X=\sum_{k\ge1}X_k$ converges and the $X_k$ are independent, the law of $X$ is the weak limit of $\mu_1*\cdots*\mu_n$.  For the Fabius distribution, the factors are rescaled uniform laws.  Their characteristic functions multiply, yielding the infinite sinc product, while their densities are finite splines that converge to the smooth up-function.  Infinite convolution is thus the bridge between probability, Fourier products, and spline approximation.  Existence requires tightness or convergence of the random sum; it is not merely formal multiplication of infinitely many measures.
\end{glossaryentry}
\begin{glossaryentry}{infinite-product}{Infinite product}{convergent product}
An \emph{infinite product} $\prod_{n\\ge0}(1+a_n)$ converges to a nonzero limit when its partial products converge and, under standard hypotheses, when $\sum|a_n|<\infty$.  Products containing zeros require a slightly broader convention.  Uniform convergence on compact sets makes a product of holomorphic factors holomorphic.  The up-function's Fourier transform is an infinite product of sinc factors, equivalently a product of cosine powers or a canonical zero product.  Taking logarithms converts the negative-Laplace product into a sum, but only where every factor is positive and branch choices are controlled.  Product recurrences are often the multiplicative form of the dyadic refinement law.
\end{glossaryentry}
\begin{glossaryentry}{integrability}{Integrability}{Lebesgue integrability; absolute integrability}
A measurable real or complex function is \emph{integrable} when the integral of its norm is finite.  This absolute notion guarantees that the positive and negative parts are both finite and that standard linearity and convergence theorems apply.  \emph{Integrable on} a set means integrable with respect to the measure restricted to that set.  The Fabius proof corpus repeatedly establishes integrability of Mellin kernels, Gamma kernels after rescaling, Fourier transforms, and large- or small-argument finite-part kernels.  Continuity on a compact interval gives integrability automatically, whereas unbounded intervals require decay estimates.  A conditionally convergent improper integral need not satisfy the same rearrangement rules as a Lebesgue-integrable function.
\end{glossaryentry}
\begin{glossaryentry}{integration-by-parts}{Integration by parts}{partial integration}
\emph{Integration by parts} is the integral form of the product rule:
\[
 \int_a^b u(x)v'(x)\,dx=[u(x)v(x)]_a^b-\int_a^b u'(x)v(x)\,dx.
\]
It transfers derivatives between factors and is valid under several regularity hypotheses.  In the Fabius theory, repeated integration by parts relates moments to inverse-dyadic values, evaluates Fourier coefficients, derives Gaussian moments, and computes period averages of products of derivatives of the periodic correction.  Compact support and infinite-order endpoint vanishing often eliminate all boundary terms, which is one of the principal advantages of the up-function as a test function.  In formalization, each disappearance of a boundary term must be justified rather than inferred from a picture.
\end{glossaryentry}
\begin{glossaryentry}{inverse-dyadic-value}{Inverse-dyadic value}{$F(2^{-n})$; reciprocal-power-of-two value}
An \emph{inverse-dyadic value} is a value of the Fabius function at $2^{-n}$.  These values form the arithmetic spine of the theory.  They are proportional to normalized half moments,
\[
 F(2^{-n})=2^{-\binom n2}\frac{d_n}{n!},
\]
under the repository's convention, and satisfy a triangular recurrence involving earlier inverse-dyadic values.  A finite weighted-path expansion solves that recurrence as a sum over ordered compositions of $n$, with the empty composition supplying $F(1)=1$.  Their exact $2$-adic valuation, denominator structure, and endpoint asymptotic connect arithmetic and analysis in an unusually explicit way.
\end{glossaryentry}
\begin{glossaryentry}{irwin-hall-distribution}{Irwin--Hall distribution}{sum of independent uniforms; uniform-sum spline}
The \emph{Irwin--Hall distribution} is the law of a finite sum of independent uniform $[0,1]$ random variables.  Its density is a compactly supported piecewise polynomial spline, and its cumulative distribution is a piecewise polynomial of one higher degree.  Rescaled and centered variants arise as the finite approximants to the Fabius law: each additional weighted uniform coordinate convolves the existing density with a smaller interval indicator.  As the number of factors grows, the knots become dyadically finer and the densities converge to the infinitely smooth up-function.  Thus the Fabius density may be viewed as an infinite-order limit of nonuniformly scaled Irwin--Hall splines.
\end{glossaryentry}
\begin{glossaryentry}{iterated-derivative}{Iterated derivative}{higher derivative; $f^{(n)}$}
The \emph{$n$th iterated derivative} is obtained by applying differentiation $n$ times, with the zeroth derivative defined as the original function.  Repeated differentiation of the Fabius scaling equation gives an exact closed identity rather than merely a recursion:
\[
 \mathcal F^{(n)}(x)=2^{\binom{n+1}{2}}\mathcal F(2^n x).
\]
For the up-function, an analogous formula involves a Thue--Morse translate combination.  These identities imply exact derivative bounds, endpoint flatness, and eventual vanishing at dyadic points.  In Lean, iterated derivatives are often expressed through a library operator whose indexing and differentiability hypotheses must be tracked explicitly, especially at order zero.
\end{glossaryentry}
\begin{glossaryentry}{interpolation}{Polynomial interpolation}{Lagrange interpolation; interpolation polynomial}
Given distinct nodes $x_0,\ldots,x_n$ and values $y_0,\ldots,y_n$, \emph{polynomial interpolation} constructs the unique polynomial $p$ of degree at most $n$ with $p(x_j)=y_j$.  Lagrange's formula is
\[
 p(x)=\sum_{j=0}^n y_j\prod_{m\ne j}\frac{x-x_m}{x_j-x_m}.
\]
The leading coefficient is therefore a weighted sum of the sampled values.  At geometric nodes the weights simplify to $q$-Pochhammer and Gaussian-binomial factors.  The Fabius $q$-binomial identities exploit exactly this leading-coefficient functional, while Thue--Morse signs annihilate lower-degree components inside dyadic blocks.  Translation invariance follows because translating the argument changes only lower coefficients, not the leading one.
\end{glossaryentry}
\begin{glossaryentry}{is-fabius-specification}{The \texorpdfstring{\texttt{IsFabius}}{IsFabius} specification}{bounded Fabius axioms; abstract Fabius candidate}
In the Lean development, \texttt{IsFabius F} is the reusable predicate asserting that a bounded candidate $F$ has the endpoint behavior, smoothness, symmetry, and differential law characteristic of the Fabius function.  Most analytic theorems are first proved for arbitrary $(F,h_F)$ satisfying this specification and are only then specialized to the canonical fixed point.  This separates mathematical consequences of the axioms from the construction proving that such a function exists uniquely.  It also makes dependencies visible: evenness of the folded up-function follows structurally from the definition, while exact support, strict positivity, or derivative identities require selected parts of the full specification.
\end{glossaryentry}
\clearpage
\hypertarget{letter:L}{}
\section*{\color{linkblue}L}
\addcontentsline{toc}{section}{L}
\markboth{L}{L}
\alphabetbar
\begin{glossaryentry}{levy-continuity}{L\'evy's continuity theorem}{continuity theorem for characteristic functions}
\emph{L\'evy's continuity theorem} says that if characteristic functions $\varphi_n(t)$ converge pointwise to a function $\varphi(t)$ that is continuous at zero, then $\varphi$ is the characteristic function of a probability law and the corresponding laws converge weakly.  Conversely, weak convergence implies pointwise convergence of characteristic functions.  In the Fabius probability construction, finite convolution characteristic functions converge to the infinite sinc product, providing a natural route to the limiting law.  Tightness and normalization are encoded in the theorem through continuity at zero and the value $\varphi(0)=1$.
\end{glossaryentry}
\begin{glossaryentry}{lambert-w}{Lambert $W$ function}{product logarithm; $W_0$; $W_{-1}$}
The \emph{Lambert $W$ function} is the multivalued inverse of $w\mapsto we^w$: $W(z)e^{W(z)}=z$.  On $[-e^{-1},0)$ it has two real branches.  The principal branch $W_0$ takes values in $[-1,0)$, while the lower branch $W_{-1}$ takes values in $(-\infty,-1]$.  The Fabius saddle phase is
\[
 \lambda(x)=-\frac{W_{-1}(-x\log2)}{\log2},
\]
which is the large positive solution of $\lambda2^{-\lambda}=x$.  Choosing $W_0$ would produce the small solution and the wrong asymptotic regime.  Expansions of $W_{-1}$ translate the exact inverse-phase result into nested powers of $-\log x$ and $\log(-\log x)$, but periodic coefficients must still be evaluated at the exact phase when sharp errors are claimed.
\end{glossaryentry}
\begin{glossaryentry}{laplace-method}{Laplace's method}{steepest descent; saddle-point method}
\emph{Laplace's method} estimates integrals with a large parameter by expanding the exponent near the point where its real part is maximal.  For a real integral $\int e^{-n\phi(t)}a(t)\,dt$, a nondegenerate minimum of $\phi$ produces a Gaussian main term and an inverse-power expansion.  Complex contour versions are called steepest descent or saddle-point methods.  The Fabius endpoint problem starts from a Bromwich inversion integral whose saddle moves to infinity as $x\to0^+$.  After solving the saddle equation with Lambert $W$, the contour is split into central and minor-arc regions; the central exponent is normalized to a Gaussian, and higher jets generate the all-orders correction.
\end{glossaryentry}
\begin{glossaryentry}{laplace-transform}{Laplace transform}{one-sided Laplace transform; moment-generating transform}
The \emph{Laplace transform} of a function supported on the nonnegative half-line is $\mathcal Lf(s)=\int_0^\infty e^{-sx}f(x)\,dx$.  For a probability distribution, replacing $s$ by $-t$ gives the moment-generating function where finite.  The Fabius density has compact support, so its transform is entire in the parameter.  The ``negative-Laplace'' product evaluates the transform at a large negative real argument, where factors take the form $(1-e^{-s/2^j})/(s/2^j)$ after normalization.  Functional equations in the original variable become scaling recurrences for the transform, and Bromwich inversion returns the density or CDF.
\end{glossaryentry}
\begin{glossaryentry}{laurent-series}{Laurent series and pole}{meromorphic expansion; principal part}
A \emph{Laurent series} about $z_0$ allows finitely or infinitely many negative powers:
\[
 f(z)=\sum_{n=-m}^\infty a_n(z-z_0)^n.
\]
If only finitely many negative powers occur, $z_0$ is a pole; their sum is the principal part, and $a_{-1}$ is the residue.  The constant coefficient $a_0$ is the finite part.  In the Mellin analysis of the Fabius periodic mean, both $\Gamma(a)$ and $\zeta(a+1)$ have simple poles at $a=0$, so their product has a double pole.  Expanding through second order isolates the regularized constant and introduces Euler's constant and the first Stieltjes constant.
\end{glossaryentry}
\begin{glossaryentry}{leading-coefficient-functional}{Leading-coefficient functional}{top-coefficient extractor}
On polynomials of degree at most $n$, the map $p\mapsto[x^n]p(x)$ is a linear functional.  Interpolation expresses it as a weighted sum of values at any $n+1$ distinct nodes:
\[
 [x^n]p(x)=\sum_{j=0}^n\frac{p(x_j)}{\prod_{m\ne j}(x_j-x_m)}.
\]
For geometric nodes these weights become Gaussian-binomial factors.  The Fabius translated-power formulas use this functional to prove that a complicated finite sum is independent of a common shift: replacing $p(x)$ by $p(x+q)$ leaves its leading coefficient unchanged.  The argument works over every field containing the rational coefficients, which explains the uniform real, complex, and Gaussian-rational specializations.
\end{glossaryentry}
\begin{glossaryentry}{least-common-multiple}{Least common multiple}{LCM; $\operatorname{lcm}$}
The \emph{least common multiple} of positive integers is the smallest positive integer divisible by all of them.  Prime by prime, its valuation is the maximum of the corresponding valuations.  The dyadic denominator $D_n$ in the Fabius arithmetic is the LCM of the reduced denominators of all relevant values on a grid of order $n$.  Bounding $D_n$ requires a common integer multiplier for every moment entering the explicit dyadic formula.  An LCM statement is stronger than a separate denominator bound only when it is shown uniformly over the entire finite family.  The convention for the LCM of an empty family is typically one.
\end{glossaryentry}
\begin{glossaryentry}{least-squares}{Least-squares approximation}{$L^2$ approximation; orthogonal projection}
A \emph{least-squares approximation} minimizes the integral of the squared error.  In a Hilbert space, the best approximation from a finite-dimensional subspace is the orthogonal projection, characterized by the residual being orthogonal to that subspace.  The Pythagorean identity
\[
 \|f-q\|_2^2=\|f-P_Vf\|_2^2+\|P_Vf-q\|_2^2
\]
proves both optimality and uniqueness.  The even Fourier--Legendre partial sum of the up-function is the orthogonal projection onto the even polynomial modes through degree $2N$.  Because the target is even, the odd mode of degree $2N+1$ is also orthogonal, so the same partial sum is optimal among all polynomials of degree at most $2N+1$.
\end{glossaryentry}
\begin{glossaryentry}{legendre-polynomial}{Legendre polynomial}{$P_n$; classical orthogonal polynomial}
The \emph{Legendre polynomials} are defined by Rodrigues' formula
\[
 P_n(x)=\frac1{2^n n!}\frac{d^n}{dx^n}(x^2-1)^n.
\]
They satisfy $P_n(1)=1$, $P_n(-x)=(-1)^nP_n(x)$, the Sturm--Liouville equation
\[
 ((1-x^2)P_n'(x))'+n(n+1)P_n(x)=0,
\]
and orthogonality $\int_{-1}^1P_mP_n=0$ for $m\ne n$, with squared norm $2/(2n+1)$.  The repository develops these facts from the polynomial definition before using them to expand Rvachev's up-function.  The bound $|P_n(x)|\le1$ on $[-1,1]$ helps upgrade coefficient summability to uniform convergence.
\end{glossaryentry}
\begin{glossaryentry}{legendre-factorial-formula}{Legendre's formula for factorial valuations}{de Polignac formula; $v_p(n!)$}
For a prime $p$, \emph{Legendre's formula} states
\[
 v_p(n!)=\sum_{j\ge1}\left\lfloor\frac{n}{p^j}\right\rfloor.
\]
Equivalently, $v_p(n!)=(n-s_p(n))/(p-1)$, where $s_p(n)$ is the base-$p$ digit sum.  At $p=2$ this becomes $v_2(n!)=n-s_2(n)$, directly linking factorial divisibility to binary weight and Thue--Morse parity.  The Fabius $2$-adic analysis uses this identity to simplify the exact valuation of inverse-dyadic values.  This theorem is unrelated to Legendre polynomials, despite sharing Legendre's name.
\end{glossaryentry}
\begin{glossaryentry}{lipschitz-continuity}{Lipschitz continuity}{Lipschitz bound; least Lipschitz constant}
A function $f$ is \emph{$L$-Lipschitz} when $|f(x)-f(y)|\le L|x-y|$ for all $x,y$.  A bounded derivative implies a Lipschitz bound by the mean value theorem.  The bounded Fabius function has least global Lipschitz constant $2$: the derivative is everywhere at most $2$, and the maximum is attained.  Lipschitz continuity gives quantitative stability of evaluation, propagates input approximation error in the computability proof, and supplies a simple modulus of uniform continuity.  ``Least constant'' is stronger than exhibiting one valid $L$; it requires an example or supremum argument showing that every smaller constant fails.
\end{glossaryentry}
\begin{glossaryentry}{local-finiteness}{Local finiteness}{locally finite family; pointwise finite sum}
A family of supported functions is \emph{locally finite} if every point has a neighborhood meeting the supports of only finitely many members.  Then its pointwise sum is locally a finite sum and inherits continuity or smoothness without a convergence theorem.  The Thue--Morse translate definition of the signed global Fabius extension uses compact support so that, for a fixed argument, only finitely many translated up-functions contribute.  Local finiteness also underlies partitions of unity.  It is stronger than pointwise finiteness and is precisely what permits differentiation term by term by reducing the issue to an ordinary finite sum near each point.
\end{glossaryentry}
\begin{glossaryentry}{logarithmic-scale}{Logarithmic scale and log-periodicity}{base-two logarithmic coordinate; multiplicative periodicity}
Passing from a positive variable $s$ to $t=\log_2s$ converts multiplication by two into translation by one.  A function $Q(s)$ satisfying $Q(2s)=Q(s)$ is therefore \emph{multiplicatively periodic}, while $P(t)=Q(2^t)$ is one-periodic.  The Fabius negative-Laplace product has a quadratic drift in $\log s$ plus such a periodic remainder.  This logarithmic scale explains why the endpoint asymptotic contains oscillations in the Lambert phase: dyadic self-similarity leaves a residue periodic in the fractional part of a base-two logarithm.  Averaging over one period extracts the constant term; centering removes it.
\end{glossaryentry}
\clearpage
\hypertarget{letter:M}{}
\section*{\color{linkblue}M}
\addcontentsline{toc}{section}{M}
\markboth{M}{M}
\alphabetbar
\begin{glossaryentry}{major-minor-arc}{Major arc and minor arc}{central arc; tail of a saddle contour}
In a contour or oscillatory integral, a \emph{major arc} is the neighborhood in which the principal contribution is extracted, while the \emph{minor arc} is the complementary region shown to be smaller.  The terminology comes from the circle method but is used more generally.  In the Fabius Bromwich analysis, the vertical contour through the saddle is split into a central window, where a Taylor expansion gives a Gaussian integral, and an outer region, where the real part of the exponent decreases enough to yield a uniform bound.  An all-orders saddle expansion requires both pieces: local algebra alone computes candidate coefficients, while the minor-arc estimate proves that contributions outside the central scaling are absorbed by the stated remainder.
\end{glossaryentry}
\begin{glossaryentry}{mean-value-theorem}{Mean value theorem}{Lagrange mean value theorem}
The \emph{mean value theorem} states that if $f$ is continuous on $[a,b]$ and differentiable on $(a,b)$, then $f(b)-f(a)=f'(c)(b-a)$ for some $c\in(a,b)$.  Its estimate form gives $|f(b)-f(a)|\le\sup|f'|\,|b-a|$.  The global Fabius derivative bound yields the sharp Lipschitz constant by this route.  Iterating a one-sided mean-value estimate also gives the first effective flatness bound $F(x)\le2xF(2x)$ and its higher-order consequences.  The theorem loses integral averaging information, which is why the fundamental-theorem-of-calculus proof recovers an additional factorial in the sharper flatness estimate.
\end{glossaryentry}
\begin{glossaryentry}{measurable-function}{Measurable function}{Borel measurable; strongly measurable}
A function is \emph{measurable} when inverse images of measurable sets are measurable; equivalently for real-valued functions, sets of the form $\{x:f(x)<a\}$ are measurable.  Continuity implies Borel measurability.  Bochner integration in normed spaces uses strong or almost-everywhere strong measurability together with integrability of the norm.  Most named Fabius functions are continuous or smooth, so measurability is automatic, but series kernels and restricted-domain expressions still need explicit measurable representatives before dominated convergence can be invoked.  Lean often packages this obligation as \texttt{AEMeasurable} or \texttt{AEStronglyMeasurable}, reflecting that altering a function on a null set does not change its integral.
\end{glossaryentry}
\begin{glossaryentry}{mellin-transform}{Mellin transform}{multiplicative Fourier transform}
The \emph{Mellin transform} of $f$ is
\[
 \mathcal Mf(s)=\int_0^\infty x^{s-1}f(x)\,dx,
\]
initially on a vertical strip where the integral converges.  It is the Fourier transform after the logarithmic substitution $x=e^t$, so it diagonalizes multiplicative scaling.  This makes it exceptionally well suited to dyadic self-similarity.  For the logarithmic Bose kernel, termwise integration gives $\mathcal M\log(1-e^{-x})=-\Gamma(s)\zeta(s+1)$.  The Fabius development then studies the meromorphic expansion at $s=0$ to compute the mean and Fourier modes of the log-periodic correction.  Poles of the Mellin transform encode powers and logarithms in asymptotic expansions.
\end{glossaryentry}
\begin{glossaryentry}{modulus-of-continuity}{Modulus of continuity}{continuity modulus}
A \emph{modulus of continuity} is a function $\omega(\delta)$ such that $|f(x)-f(y)|\le\omega(|x-y|)$ and $\omega(\delta)\to0$ as $\delta\to0$.  Lipschitz continuity corresponds to $\omega(\delta)=L\delta$.  A recursive modulus converts this analytic control into an effective precision transformer.  For the Fabius function, the sharp derivative bound gives the global modulus $2\delta$.  Higher-order bounds yield moduli for derivatives as well.  A modulus is uniform over the domain; pointwise continuity alone supplies neighborhoods that may depend noncomputably on both the point and the desired precision.
\end{glossaryentry}
\begin{glossaryentry}{moment-numerators}{Moment and half-moment numerators}{$F_n$ and $G_n$ sequences; division-free numerator sequences}
The rational moment and half-moment sequences admit normalized forms
\[
 c_n=\frac{F_n}{(2n+1)!!\prod_{j=1}^n(2^{2j}-1)},
 \qquad
 d_n=\frac{G_n}{(n+1)!\prod_{j=1}^n(2^j-1)}
\]
under the source conventions.  The natural numbers $F_n$ and $G_n$ are the corresponding \emph{moment numerators}.  They can be defined by division-free triangular recurrences, making integrality part of the construction rather than a post hoc claim about reduced fractions.  Their parity controls $2$-adic valuations, and their growth feeds denominator and asymptotic calculations.  The notation $F_n$ here denotes a sequence of integers, not the Fabius function itself.
\end{glossaryentry}
\begin{glossaryentry}{moment-generating-function}{Moment-generating function}{MGF; bilateral Laplace transform of a law}
The \emph{moment-generating function} of a random variable $X$ is $M_X(t)=\mathbb E[e^{tX}]$ wherever finite.  Its derivatives at zero are the raw moments: $M_X^{(n)}(0)=\mathbb E[X^n]$.  Independent sums turn into products of MGFs.  Since the Fabius law has compact support, its MGF is entire, and the weighted-uniform representation gives an infinite product.  A related exponential generating function for half moments satisfies a dyadic functional equation.  The characteristic function is obtained by setting $t=2\pi i\xi$ under the Fourier convention and is always bounded, whereas an MGF may diverge for general unbounded laws.
\end{glossaryentry}
\begin{glossaryentry}{moment}{Moment of a distribution}{raw moment; moment sequence; $c_n$}
The \emph{$n$th moment} of a random variable or density is $\mathbb E[X^n]=\int x^n\,d\mu(x)$ when the integral exists.  Moments encode location, spread, and higher shape information and appear as derivatives of characteristic or moment-generating functions at the origin.  For the even up-function, all odd full moments vanish and the even moments form a rational sequence $c_n$.  They satisfy a triangular recurrence, a Bernoulli recurrence, and a numerator/product formula.  These moments enter exact dyadic evaluations through finite Taylor integrals.  Compact support guarantees existence of moments of every order, but a moment sequence determines a measure only under additional determinacy hypotheses.
\end{glossaryentry}
\begin{glossaryentry}{monotonicity}{Monotonicity}{increasing function; strictly increasing; monotone map}
A function is \emph{nondecreasing} if $x\le y$ implies $f(x)\le f(y)$ and \emph{strictly increasing} if strict inequality is preserved.  A nonnegative derivative implies nondecreasing behavior; a positive derivative on every interior point gives strict increase.  The bounded Fabius function is constant outside $[0,1]$ and strictly increasing inside it.  The up-function increases on $[-1,0]$ and decreases symmetrically on $[0,1]$, hence is strictly unimodal.  Formal proofs separate interval-restricted monotonicity from global monotonicity and distinguish positivity at every point from positivity almost everywhere, since only the former immediately gives strictness by elementary calculus.
\end{glossaryentry}
\begin{glossaryentry}{mean-periodic}{Period mean and centered periodic function}{average over one period; zero-mean normalization}
For an integrable one-periodic function $P$, its \emph{period mean} is $\overline P=\int_0^1P(t)\,dt$.  The centered function $\Psi(t)=P(t)-\overline P$ has mean zero and the same nonzero Fourier modes.  In the Fabius negative-Laplace decomposition, separating the mean is analytically useful because the mean joins the constant term of the sharp asymptotic, while $\Psi$ carries the genuine log-periodic fluctuation.  The repository computes the mean by a Mellin finite-part integral and reconstructs the centered part by its Gamma--zeta Fourier series.  A zero-mean periodic function need not be small pointwise; centering removes only the constant Fourier coefficient.
\end{glossaryentry}
\clearpage
\hypertarget{letter:N}{}
\section*{\color{linkblue}N}
\addcontentsline{toc}{section}{N}
\markboth{N}{N}
\alphabetbar
\begin{glossaryentry}{numerical-error}{Absolute, relative, and truncation error}{approximation error; evaluator error}
The \emph{absolute error} of an approximation $\widetilde y$ to $y$ is $|\widetilde y-y|$; the relative error is $|\widetilde y-y|/|y|$ when $y\ne0$.  A truncation error is the contribution omitted by stopping a series, recursion, or asymptotic expansion.  Near the Fabius endpoint, values are extraordinarily small, so relative error and logarithmic error are often more informative than absolute error.  The computable-spline proof instead needs an absolute dyadic bound uniform in the input.  Exact rational dyadic evaluation has no rounding error, although numerator and denominator growth affect cost.  Numerical checks of saddle coefficients compare the asymptotic residual at several orders against exact rational values.
\end{glossaryentry}
\begin{glossaryentry}{newton-iteration}{Fixed-point and Newton iteration for the Lambert phase}{iterative solution of the saddle equation}
An \emph{iteration} solves an implicit equation by repeatedly applying a map whose fixed point is the desired root.  For dyadic $x=2^{-n}$, the Lambert phase satisfies $\lambda=n+\log\lambda/\log2$, suggesting the simple fixed-point iteration $\lambda_{m+1}=n+\log\lambda_m/\log2$.  Newton's method instead uses the derivative of $h(\lambda)=\lambda2^{-\lambda}-x$ and generally converges quadratically near the correct root.  The lower branch must be selected by an initial value in the large-$\lambda$ basin.  Numerical iteration is used only for verification or evaluation; the formal asymptotic statements define the phase exactly through Lambert $W$ or its implicit equation.
\end{glossaryentry}
\begin{glossaryentry}{negative-laplace-product}{Negative-Laplace product}{two-sided Laplace product; logarithmic product}
The \emph{negative-Laplace product} is the product representation obtained by evaluating the Laplace or moment-generating transform of the Fabius density at a large positive growth parameter.  In normalized form it contains factors
\[
 \frac{1-e^{-s/2^j}}{s/2^j}.
\]
Taking logarithms separates a quadratic drift in $\log_2s$, a forward exponential tail, and an exactly multiplicatively periodic correction.  The small-scale factors require estimates uniform in $j$, while the forward tail is exponentially small for large $s$.  This decomposition is the analytic starting point for the corrected sharp endpoint asymptotic and reveals the periodic term that a purely continuous approximation of the dyadic sum misses.
\end{glossaryentry}
\begin{glossaryentry}{normalization}{Normalization}{choice of scale; canonical normalization}
A \emph{normalization} fixes multiplicative, additive, or coordinate ambiguities so that an object is uniquely specified.  The up-function is normalized by $\Up(0)=1$ and support $[-1,1]$; the bounded Fabius CDF is normalized by endpoint values $F(0)=0$ and $F(1)=1$; its probability density has integral one.  Fourier-transform conventions add another normalization through the placement of $2\pi$.  The Fabius literature contains several functions denoted by the same symbol but related by translation, reflection, differentiation, or signed continuation.  Stating the normalization is therefore part of the definition, not a notational afterthought.  Many apparent source errors are actually mismatches between normalizations or domains.
\end{glossaryentry}
\begin{glossaryentry}{nowhere-analytic}{Nowhere analytic on a set}{nonanalytic at every point}
A smooth function is \emph{nowhere analytic on} a set if it is analytic at none of its points.  The Fabius function provides a canonical example: it is $C^\infty$ everywhere but nowhere analytic on $[0,1]$, while being analytic outside that interval because it is locally constant there.  At dyadic points, all sufficiently high derivatives vanish, so the Taylor series is a polynomial, but the function is not locally polynomial.  Nondyadic points are handled by self-similarity and propagation of analyticity to a contradiction with endpoint flatness or support.  ``Nowhere analytic'' is much stronger than ``not entire'' or ``not analytic at the endpoints.''
\end{glossaryentry}
\clearpage
\hypertarget{letter:O}{}
\section*{\color{linkblue}O}
\addcontentsline{toc}{section}{O}
\markboth{O}{O}
\alphabetbar
\begin{glossaryentry}{ordered-composition}{Ordered composition of an integer}{composition; weighted composition}
An ordered composition of $n$ is a finite sequence of positive integers $(r_1,\ldots,r_m)$ with $r_1+\cdots+r_m=n$.  Order matters: $(1,2)$ and $(2,1)$ are different compositions.  The empty sequence is the unique composition of zero.  Iterating a triangular recurrence naturally produces a sum over compositions, because each recursive descent chooses a positive decrement.  In the inverse-dyadic formula for $F(2^{-n})$, the partial sums $s_j=r_1+\cdots+r_j$ label the successive states, and each edge contributes a factor involving $(2^{s_j}-1)^{-1}$ and $(r_j+1)!^{-1}$.  The resulting finite weighted-path expansion is a closed form rather than another recursion.
\end{glossaryentry}
\begin{glossaryentry}{orthogonality}{Orthogonality}{orthogonal functions; zero inner product}
Two functions are \emph{orthogonal} in $L^2([-1,1])$ when $\int_{-1}^1f(x)g(x)\,dx=0$.  Orthogonality generalizes perpendicular vectors and turns coefficient extraction into an inner-product calculation.  Legendre polynomials of different degrees are orthogonal, and even functions are orthogonal to every odd polynomial on a symmetric interval.  These facts yield the exact Fourier--Legendre coefficients of the up-function and the Pythagorean identity behind least-squares optimality.  Orthogonality is weaker than statistical independence and does not imply pointwise separation; it is a statement about averaged products.
\end{glossaryentry}
\clearpage
\hypertarget{letter:P}{}
\section*{\color{linkblue}P}
\addcontentsline{toc}{section}{P}
\markboth{P}{P}
\alphabetbar
\begin{glossaryentry}{paley-wiener}{Paley--Wiener theorem}{Fourier support--growth theorem}
A \emph{Paley--Wiener theorem} relates compact support of a function or distribution to holomorphic extension and exponential growth of its Fourier transform.  Roughly, a smooth function supported in $[-A,A]$ has an entire transform that decays rapidly on the real axis and grows no faster than $e^{2\pi A|\operatorname{Im}z|}$ in the complex plane; conversely, suitable entire functions with those bounds arise from compactly supported smooth functions.  The original up-function proof mentions this theorem but chooses a probabilistic convolution argument to establish support.  The theorem still explains the compatibility of the entire sinc product with a compactly supported inverse transform.
\end{glossaryentry}
\begin{glossaryentry}{parity}{Parity, evenness, and oddness}{even integer; odd integer; parity argument}
The \emph{parity} of an integer records whether it is divisible by two.  Modulo $2$, signs and binomial coefficients often simplify dramatically.  In the Fabius arithmetic, Lucas's theorem counts odd binomial coefficients from binary digits, proving oddness of moment numerators and exact $2$-adic valuations.  For functions, ``even'' means $f(-x)=f(x)$ and ``odd'' means $f(-x)=-f(x)$; evenness of the up-function eliminates odd moments and odd Legendre coefficients.  The same word thus describes integer congruence and reflection symmetry, two related but distinct uses that interact through Thue--Morse signs.
\end{glossaryentry}
\begin{glossaryentry}{parity-power-series}{Parity-power form}{binomial-parity power series; parity-power summand}
The \emph{parity-power form} is a representation of the global Fabius function whose coefficients or exponents are built from parities of binomial coefficients and binary-weight signs.  It is equivalent in the limit to the binary-reduction and $q$-binomial series but has delicate indexing.  The all-real formula begins with the $m=0$ term; omitting it happens to preserve the formula on a half-open interval but fails at the right endpoint.  Integer exponents are used so the order-zero exponent can genuinely be negative.  The formalization proves convergence and the exact domain of the one-indexed variant rather than treating the indexing correction as harmless notation.
\end{glossaryentry}
\begin{glossaryentry}{partial-sum}{Partial sum and remainder}{truncation; tail}
The \emph{$N$th partial sum} of a series $\sum_{n\ge0}a_n$ is $S_N=\sum_{n<N}a_n$ or, under another convention, the sum through $n=N$; the indexing convention must be stated.  The remainder or tail is the difference between the full sum and $S_N$.  Fourier--Legendre partial sums are finite polynomials with an exact least-squares interpretation.  Saddle partial sums are finite asymptotic corrections whose remainder is $O(\lambda^{-N})$.  Binary-reduction telescopes have a finite remainder term that converges to zero.  A partial sum may have properties, such as dependence on a translation parameter, that disappear only in the limit.
\end{glossaryentry}
\begin{glossaryentry}{partition-of-unity}{Partition of unity}{translate partition; cardinal partition}
A \emph{partition of unity} is a family of functions $(\phi_j)$ satisfying $\sum_j\phi_j(x)=1$ at every point, usually with nonnegative values and local finiteness.  Integer translates of Rvachev's up-function form such a partition.  Because the support has length two, on a fundamental interval the global sum reduces to an elementary two-term identity.  Partitions of unity allow local data to be blended into global approximations and imply reproduction of constants.  Higher polynomial reproduction identities strengthen this to exact recovery of low-degree polynomials from weighted translates.  Endpoint weighting matters for finite approximants even when the limiting smooth partition is unambiguous.
\end{glossaryentry}
\begin{glossaryentry}{periodic-correction}{Periodic correction in a self-similar asymptotic}{log-periodic fluctuation; $\Psi$}
A \emph{periodic correction} is a bounded periodic term that remains after the smooth drift of a discrete-scale asymptotic has been removed.  For the Fabius endpoint problem, dyadic scaling produces a one-periodic function of the lower-Lambert phase.  Its centered version $\Psi$ is smooth, has mean zero, and has explicit nonzero Gamma--zeta Fourier coefficients.  It appears at order one in the logarithm, not as a negligible higher correction.  Replacing it by a constant leaves an oscillatory residual bounded away from the claimed inverse-logarithmic error scale.  The phenomenon is common in digital trees, divide-and-conquer recurrences, and multiplicatively periodic renewal problems.
\end{glossaryentry}
\begin{glossaryentry}{periodic-function}{Periodic function}{one-periodic function; period}
A function $p$ is \emph{periodic} with period $T\ne0$ when $p(x+T)=p(x)$ for every $x$.  A fundamental period is the smallest positive period when one exists.  One-periodic functions are naturally expanded in integer Fourier modes.  In the Fabius asymptotic, periodicity arises after converting multiplicative dyadic scaling into additive translation by the logarithm.  Every saddle coefficient $A_j(\lambda)$ is one-periodic because it is a polynomial in derivatives of the basic periodic correction.  Periodicity alone does not imply boundedness unless continuity or another regularity condition is supplied; continuous periodic functions are bounded on the real line by compactness of one period.
\end{glossaryentry}
\begin{glossaryentry}{pointwise-convergence}{Pointwise convergence}{convergence at each point}
A sequence of functions $f_n$ converges \emph{pointwise} to $f$ when $f_n(x)\to f(x)$ for each fixed $x$.  The rate and the index after which an approximation is good may depend on $x$.  Pointwise convergence does not in general preserve continuity, permit interchange with integration, or justify differentiation.  The finite polynomial step approximants to the up-function have a corrected pointwise limit, while stronger uniform estimates are developed separately for computability and spline approximation.  Weak convergence of measures can imply convergence of cumulative distribution functions only at continuity points of the limiting CDF; the Fabius law is continuous everywhere, simplifying that passage.
\end{glossaryentry}
\begin{glossaryentry}{poisson-summation}{Poisson summation formula}{periodization--sampling identity}
The \emph{Poisson summation formula} relates samples of a function to samples of its Fourier transform:
\[
 \sum_{n\in\mathbb Z}f(x+n)=\sum_{k\in\mathbb Z}\widehat f(k)e^{2\pi ikx}
\]
under suitable decay and regularity hypotheses.  Applying it to the up-function yields its integer-translate partition of unity because the sinc product vanishes at every nonzero integer frequency.  Sampling at half-integers gives the cosine expansion on $[-1,1]$, and other lattice spacings produce exact identities involving values of the transform.  The formula is sensitive to Fourier normalization, lattice scale, and the phase $x$; several source corrections in the corpus restore omitted variables or endpoint restrictions.
\end{glossaryentry}
\begin{glossaryentry}{polynomial-reproduction}{Polynomial reproduction}{Strang--Fix-type identity; exact reproduction}
A translate family \emph{reproduces polynomials} through degree $m$ if weighted sums of its shifts equal every polynomial of degree at most $m$.  Reproduction of constants is a partition of unity.  Higher reproduction follows from moment cancellations or zeros of the Fourier transform at lattice frequencies.  The up-function and its finite spline relatives possess exact identities of this kind, making them useful in approximation theory.  Thue--Morse signed sums furnish a complementary cancellation principle: they annihilate low-degree polynomials over dyadic blocks.  Reproduction is an exact algebraic statement, stronger than approximating polynomials with small error.
\end{glossaryentry}
\begin{glossaryentry}{polynomial-step-approximant}{Polynomial step approximant}{finite convolution approximant; piecewise-polynomial approximant}
A \emph{polynomial step approximant} is a finite-stage function obtained by repeated integration or convolution of step functions, hence piecewise polynomial on a finite grid.  The first source paper constructs such approximants from finite convolutions of interval measures.  Their knot set is dyadic, their degree increases with the stage, and they converge pointwise to the up-function.  A half-weight convention at shared endpoints corrects the literal closed-interval indicator formula.  Centered uniform splines used later for computability are a quantitatively controlled version of the same basic idea, with explicit uniform error instead of only qualitative convergence.
\end{glossaryentry}
\begin{glossaryentry}{power-series}{Power series and radius of convergence}{Taylor series about the origin}
A \emph{power series} is $\sum_{n\ge0}a_n(z-z_0)^n$.  It converges inside a disk of radius $R$, diverges outside, and defines an analytic function in the interior.  The sinc product and moment-generating functions have entire power series, while a formal saddle series is not asserted to converge.  The up-transform is even, so only even powers occur and their coefficients are the moment sequence after normalization.  Coefficient comparison in functional equations yields exact recurrences.  For the Fabius function itself, the Taylor series at an endpoint is identically zero because every derivative vanishes, illustrating that a smooth function need not equal its formal Taylor series.
\end{glossaryentry}
\begin{glossaryentry}{prefix-sum}{Prefix sum and iterated prefix sum}{cumulative sum; summatory sequence}
For a sequence $(a_n)$, its \emph{prefix sum} is $A(N)=\sum_{n<N}a_n$.  Iterating the operation produces discrete antiderivatives and turns finite differences back into the original sequence.  Thue--Morse prefix sums exhibit exact block cancellations and long zero runs; $k$-fold prefix sums give polygonal or spline-like interpolants after rescaling.  The audited draft in the repository formalizes these identities and corrects normalization errors in a proposed approximation scheme.  Prefix-sum conventions differ on whether the upper bound is included, so order-zero and endpoint formulas require explicit indexing.
\end{glossaryentry}
\begin{glossaryentry}{primitive-recursion}{Primitive recursion}{primitive-recursive function}
A \emph{primitive-recursive function} is built from zero, successor, projections, composition, and recursion over a natural-number counter with a previously computed value.  This class is total and effectively computable, though smaller than the class of all computable functions.  The Fabius computability proof keeps the evaluator at the natural-number level: it computes Thue--Morse bits, finite spline sums, integer powers, and dyadic rounding by primitive-recursive operations.  Proving a real function computable requires more than labeling its mathematical formula algorithmic; one must give a uniform encoding and prove that every component and error modulus belongs to the accepted recursion class.
\end{glossaryentry}
\begin{glossaryentry}{probability-law}{Probability law}{distribution of a random variable; pushforward probability measure}
The \emph{probability law} of a random variable $X:\Omega\to\mathbb R$ is the measure $\mu_X(A)=\mathbb P(X\in A)$, the pushforward of the underlying probability measure through $X$.  Different random variables on different sample spaces can have the same law.  The Fabius law is that of the convergent weighted sum $\sum_{k\ge1}2^{-k}U_k$.  Its CDF is the bounded Fabius function, its density is a translated up-function, and its characteristic function is the sinc product.  These are three representations of one law.  Independence belongs to a chosen realization of the law; convolution and the transform product are the law-level consequences.
\end{glossaryentry}
\begin{glossaryentry}{product-measure-independence}{Product measure and independence}{independent coordinates; product probability space}
A \emph{product measure} combines measures on coordinate spaces so that rectangular events have product mass.  Coordinate projections on the resulting product probability space are independent.  The Fabius random variable is constructed from an infinite sequence of independent uniform coordinates, each scaled by $2^{-k}$.  Independence implies that the law of a finite partial sum is the convolution of the coordinate laws and that its characteristic function is the product of the coordinate characteristic functions.  Infinite product spaces require a measure-extension theorem, but boundedness of the weighted series makes the resulting random variable straightforward once the coordinate space exists.
\end{glossaryentry}
\begin{glossaryentry}{prouhet-thue-morse}{Prouhet--Thue--Morse cancellation}{Prouhet identity; equal-power partition}
The \emph{Prouhet--Thue--Morse cancellation} states that the Thue--Morse signs annihilate low powers over a complete dyadic block:
\[
 \sum_{n=0}^{2^m-1}(-1)^{s_2(n)}n^k=0
 \qquad(0\le k<m).
\]
Equivalently, numbers with even and odd binary digit sum have equal sums of like powers through degree $m-1$.  The first nonzero degree has an explicit product value.  This triangular cancellation is central to derivative formulas and exact dyadic evaluations: after expanding a translated power, all lower monomials disappear.  It is independent of, but meshes perfectly with, the geometric-node $q$-binomial interpolation cancellation.
\end{glossaryentry}
\begin{glossaryentry}{pushforward-measure}{Pushforward measure}{image measure; map of a measure}
Given a measurable map $T:X\to Y$ and a measure $\mu$ on $X$, the \emph{pushforward} $T_*\mu$ is defined by $(T_*\mu)(A)=\mu(T^{-1}(A))$.  It is the law of $T(X)$ when $\mu$ is a probability measure.  Affine pushforwards rescale and translate interval measures; convolution can be expressed as the pushforward of a product measure under addition.  The finite and infinite weighted-uniform models of the Fabius function are naturally formulated with pushforwards.  Integration obeys the change-of-variables identity $\int_Y f\,d(T_*\mu)=\int_X f\circ T\,d\mu$ whenever either side is defined.
\end{glossaryentry}
\clearpage
\hypertarget{letter:Q}{}
\section*{\color{linkblue}Q}
\addcontentsline{toc}{section}{Q}
\markboth{Q}{Q}
\alphabetbar
\begin{glossaryentry}{q-binomial-theorem}{Finite $q$-binomial theorem}{Gaussian binomial theorem}
The \emph{finite $q$-binomial theorem} states
\[
 \prod_{j=0}^{n-1}(1+zq^j)
   =\sum_{k=0}^n q^{\binom k2}\qbinom{n}{k}{q} z^k.
\]
Variants with minus signs or reversed powers are obtained by substitution.  It is the $q$-analogue of $(1+z)^n=\sum_k\binom nkz^k$.  In the Fabius development it organizes finite products and proves identities among half-base interpolation weights.  The full translated Thue--Morse formula is not merely an application of this theorem; it also needs leading-coefficient extraction and dyadic cancellation.  The theorem nevertheless explains the recurring factor $q^{\binom k2}$ and the triangular power of two in exact formulas.
\end{glossaryentry}
\begin{glossaryentry}{q-calculus}{$q$-calculus}{quantum calculus; basic hypergeometric calculus}
\emph{$q$-calculus} is a deformation of ordinary calculus and combinatorics in which integers, factorials, binomial coefficients, derivatives, and products acquire a parameter $q$.  Typical objects are the $q$-integer $[n]_q=(1-q^n)/(1-q)$, the $q$-factorial, the $q$-Pochhammer symbol, and the Gaussian binomial coefficient.  As $q\to1$, many recover their classical counterparts.  The Fabius formulas use the fixed half-base $q=1/2$, so the theory is finite algebra rather than a limiting deformation.  Geometric interpolation nodes and dyadic scaling force this base.  Calling the identities ``half-base $q$-calculus'' highlights the shared mechanism behind their product denominators, interpolation weights, and finite $q$-binomial sums.
\end{glossaryentry}
\begin{glossaryentry}{q-integer-factorial}{$q$-integer and $q$-factorial}{quantum integer; basic factorial}
The \emph{$q$-integer} is $[n]_q=(1-q^n)/(1-q)=1+q+\cdots+q^{n-1}$, and the \emph{$q$-factorial} is $[n]_q!=\prod_{j=1}^n[j]_q$.  They interpolate ordinary integers and factorials as $q\to1$.  Gaussian binomial coefficients can be written $\qbinom{n}{k}{q}=[n]_q!/([k]_q![n-k]_q!)$.  At $q=1/2$, powers of two may be extracted to rewrite these factors as products of $2^j-1$, which are precisely the odd denominators prominent in Fabius moment arithmetic.  Different authors normalize $q$-factorials with or without powers of $1-q$; formulas must therefore state the convention rather than rely on the notation alone.
\end{glossaryentry}
\begin{glossaryentry}{q-pochhammer}{$q$-Pochhammer symbol}{shifted $q$-factorial; $(a;q)_n$}
The \emph{$q$-Pochhammer symbol} is
\[
 (a;q)_n=\prod_{j=0}^{n-1}(1-aq^j),
 \qquad (a;q)_0=1,
\]
and, when convergent, $(a;q)_\infty=\prod_{j\ge0}(1-aq^j)$.  It is the basic product notation of $q$-series.  Gaussian binomial coefficients satisfy $\qbinom{n}{k}{q}=(q;q)_n/((q;q)_k(q;q)_{n-k})$.  In the Fabius half-base formulas, $q=1/2$ packages products of geometric-node differences and factors $1-2^{-j}$.  The notation is unrelated to the ordinary rising Pochhammer symbol $(a)_n=a(a+1)\cdots(a+n-1)$ unless a limiting or conversion formula is explicitly given.
\end{glossaryentry}
\begin{glossaryentry}{quadratic-logarithmic-law}{Quadratic logarithmic decay law}{log-squared asymptotic; de Bruijn-type scale}
A \emph{quadratic logarithmic law} has leading logarithm proportional to $-(\log(1/x))^2$.  The Fabius endpoint satisfies such a law: after using the exact Lambert phase $\lambda$, the dominant term in $\log F(x)$ is $-(\log2)\lambda^2/2$, with lower linear, logarithmic, constant, and periodic terms.  Since $\lambda\sim\log_2(1/x)$, this places $F(x)$ between algebraic flatness and $e^{-c/x}$ decay.  The quadratic term arises by summing contributions across the roughly $\log_2 s$ dyadic scales active in the negative-Laplace product.  A coarse quadratic-log law captures the scale but not an asymptotic equivalent; the periodic and lower-order terms remain essential.
\end{glossaryentry}
\clearpage
\hypertarget{letter:R}{}
\section*{\color{linkblue}R}
\addcontentsline{toc}{section}{R}
\markboth{R}{R}
\alphabetbar
\begin{glossaryentry}{rapid-decrease}{Rapid decrease and Schwartz space}{rapidly decreasing function; Schwartz function; test function}
A smooth function $f$ is \emph{rapidly decreasing} if $x^m f^{(n)}(x)$ is bounded for every $m,n\ge0$.  Such functions form the Schwartz space $\mathcal S(\mathbb R)$, which is preserved by the Fourier transform.  Every compactly supported smooth function is Schwartz.  In the original construction, the infinite sinc product is shown to be rapidly decreasing on the real axis, so Fourier inversion produces a smooth rapidly decreasing function; a separate support argument then proves it is actually compactly supported.  Rapid decrease permits unrestricted integration by parts and ensures absolute convergence of many Fourier integrals.  The notation $\mathcal D(\mathbb R)$ usually denotes the smaller space $C_c^\infty(\mathbb R)$ of compactly supported test functions.
\end{glossaryentry}
\begin{glossaryentry}{rational-recurrence}{Rational recurrence}{exact arithmetic recurrence}
A \emph{rational recurrence} defines each term from earlier terms using exact rational operations.  The moment and half-moment sequences satisfy triangular recurrences with denominators involving $2^n-1$ or $2^{2n}-1$.  The repository also defines division-free natural numerator recurrences and proves that their quotients agree with the rational sequences.  This separation supports integrality, parity, and valuation arguments without repeatedly reasoning about reduced fractions.  A recurrence may be executable yet inefficient if it recomputes a large prefix; tables, memoization, and triangular dynamic programming turn it into a practical exact evaluator.
\end{glossaryentry}
\begin{glossaryentry}{recurrence-relation}{Recurrence relation}{recursive sequence equation; triangular recurrence}
A \emph{recurrence relation} expresses a sequence term in terms of earlier terms.  It is \emph{triangular} when the $n$th equation contains only indices below $n$ after the leading term is isolated, so initial data determine the sequence uniquely.  Fabius moments, half moments, numerator sequences, saddle exponent coefficients, formal logarithm coefficients, and inverse-dyadic values all obey such recurrences.  Different recurrences reveal different structures: Bernoulli recurrences connect to classical special numbers; division-free recurrences expose integrality; composition expansions solve a recurrence by summing over all paths.  Boundary cases matter when the nominal leading coefficient vanishes, as happens in an inverse-dyadic recurrence at $n=0$.
\end{glossaryentry}
\begin{glossaryentry}{refinement-invariance}{Refinement and representation invariance}{dyadic refinement invariance}
A dyadic formula has \emph{refinement invariance} when replacing $(a,n)$ by $(2a,n+1)$ leaves its value unchanged, reflecting $a/2^n=2a/2^{n+1}$.  This is the fundamental representation-independence check for an exact dyadic evaluator.  The Fabius development proves table-prefix stability and refinement invariance before bridging the rational algorithm to the analytic function.  Similar invariance holds for the finite $q$-binomial expression.  Without it, an algorithm would compute a property of a numeral presentation rather than a well-defined function of the rational number.
\end{glossaryentry}
\begin{glossaryentry}{refinement-equation}{Refinement equation}{two-scale relation; dilation equation}
A \emph{refinement equation} expresses a function as a finite linear combination of dilated translates of itself,
\[
 \phi(x)=\sum_k h_k\phi(2x-k),
\]
or gives an equivalent relation for its derivative.  Such equations are central to subdivision schemes, splines, and wavelets.  Rvachev's up-function is governed by a differential refinement equation, while its Fourier transform satisfies a multiplicative two-scale relation.  The dilation factor determines the natural grid and product frequencies.  Refinement equations alone may have many distributional or nonnormalized solutions; support, positivity, smoothness, and a value such as $\phi(0)=1$ select the canonical up-function.
\end{glossaryentry}
\begin{glossaryentry}{regularity}{Regularity of a function}{continuity class; differentiability class; $C^k$; $C^\infty$}
\emph{Regularity} describes how many derivatives a function has and how those derivatives behave.  $C^k$ means derivatives through order $k$ exist and are continuous; $C^\infty$ means this for every finite order.  Analyticity is stronger and requires local equality with a power series.  The Fabius fixed-point construction begins with continuity, recovers one derivative from the integral equation, and bootstraps indefinitely through dyadic scaling.  Formal libraries distinguish finite differentiability, smoothness, and analyticity by different indices or predicates.  In Lean's current conventions, visually similar top elements in extended index types can encode different regularity notions, so source comments must be checked against the exact type.
\end{glossaryentry}
\begin{glossaryentry}{representation-independence}{Representation independence}{well-definedness on equivalence classes}
A construction is \emph{representation independent} when equivalent encodings produce the same mathematical object.  Rational numbers may be written by many numerator/denominator pairs; dyadic rationals have infinitely many refinements; a finite translated formula may use different common shifts.  The Fabius arithmetic proves invariance under reduction and dyadic refinement, while the $q$-binomial identities prove invariance under a common translation of inner powers.  Such theorems are not implementation niceties.  They establish that a computation descends from raw syntax to the quotient object---the rational number, polynomial functional, or analytic value---that the notation purports to denote.
\end{glossaryentry}
\begin{glossaryentry}{reshetnikov-numbers}{Reshetnikov numbers}{$R_n$}
The \emph{Reshetnikov numbers} $R_n$ are integer normalizations of the Fabius half moments introduced in the arithmetic paper to study Vladimir Reshetnikov's denominator conjecture.  They remove the explicit odd product and double-factorial denominators from $d_n$.  The formal development first defines the associated rational expression, proves its integrality, and then proves $R_n$ is odd for $n\ge1$.  That oddness is equivalent to the exact $2$-adic valuation of $F(2^{-n})$.  The sequence is therefore a bridge between rational moment recurrences and a clean integer-parity statement.
\end{glossaryentry}
\begin{glossaryentry}{riemann-zeta}{Riemann zeta function}{$\zeta(s)$}
The \emph{Riemann zeta function} is $\zeta(s)=\sum_{n\ge1}n^{-s}$ for $\Re s>1$.  It extends meromorphically to $\mathbb C$ with a single simple pole at $s=1$ of residue one and satisfies a functional equation relating $s$ to $1-s$.  In the Fabius Mellin transform, zeta appears because the Bose logarithm expands as a sum over $n^{-1}e^{-nx}$.  Values on the vertical line $1-2\pi ik/\log2$ determine Fourier modes of the periodic correction.  The classical zero-free theorem on $\Re s=1$ proves those modes nonzero.  Bernoulli numbers encode zeta values at nonpositive integers and even positive integers, explaining another link between the arithmetic and analytic layers.
\end{glossaryentry}
\begin{glossaryentry}{rodrigues-formula}{Rodrigues' formula}{derivative definition of Legendre polynomials}
\emph{Rodrigues' formula} defines a classical orthogonal polynomial by differentiating a simple weight expression.  For Legendre polynomials,
\[
 P_n(x)=\frac{1}{2^n n!}\frac{d^n}{dx^n}(x^2-1)^n.
\]
It immediately gives polynomiality, degree, parity, and endpoint information, while repeated integration by parts proves orthogonality.  The repository constructs Legendre polynomials from this formula rather than importing every property as a black box.  This makes the coefficient normalization used in the up-function expansion explicit and aligns the formal proof with the later derivative and binomial formulas.
\end{glossaryentry}
\begin{glossaryentry}{rvachev-up-function}{Rvachev's up-function}{$\Up$; \texttt{rvachevUp}; finite function; Rvachev function}
\emph{Rvachev's up-function} is the unique even, nonnegative, compactly supported $C^\infty$ function with support $[-1,1]$, value $\Up(0)=1$, and differential refinement law
\[
 \Up'(x)=2\bigl(\Up(2x+1)-\Up(2x-1)\bigr).
\]
It is positive on $(-1,1)$, strictly increases to its maximum at zero, and then decreases symmetrically.  On the unit interval it is a translate/reflection of the bounded Fabius function: $F(x)=\Up(x-1)$.  It is also the density, after affine rescaling, of the Fabius probability law.  Its Fourier transform is the infinite sinc product, its integer translates form a partition of unity, and its dyadic values and derivatives are rational.
\end{glossaryentry}
\clearpage
\hypertarget{letter:S}{}
\section*{\color{linkblue}S}
\addcontentsline{toc}{section}{S}
\markboth{S}{S}
\alphabetbar
\begin{glossaryentry}{smoothness}{$C^k$ and $C^\infty$ smoothness}{continuously differentiable; infinitely differentiable; \texttt{ContDiff}}
A function is of class $C^k$ when derivatives through order $k$ exist and are continuous; it is $C^\infty$ or smooth when this holds for every finite $k$.  Smoothness is a local regularity condition and does not imply analyticity.  In Lean, the regularity index is represented by an extended natural, and the symbols used for ``all finite orders'' and for the analytic exponent must not be conflated.  The Fabius and up-functions are smooth on all of $\mathbb R$, including the boundary of their nonconstant regions.  Their exterior constancy forces all positive-order derivatives to vanish at the endpoints.  Iterating the scaling equation gives exact formulas for every derivative and sharp global derivative bounds.
\end{glossaryentry}
\begin{glossaryentry}{saddle-correction}{Saddle correction}{first saddle correction; higher-order steepest-descent term}
A \emph{saddle correction} is a term beyond the leading Gaussian approximation in a steepest-descent expansion.  It combines cubic and higher derivatives of the exponent, derivatives of the amplitude, and Gaussian moments.  For the Fabius function the first correction is a periodic function of the exact Lambert phase divided by $\lambda$.  Higher corrections are generated recursively or by explicit finite combinatorial formulas.  Adding one correction generally improves the logarithmic residual from $O(\lambda^{-1})$ to $O(\lambda^{-2})$, provided the analytic remainder theorem is proved at the corresponding order.  A formally computed correction without a uniform contour estimate is only a candidate coefficient.
\end{glossaryentry}
\begin{glossaryentry}{saddle-coefficients}{Saddle mass and logarithmic coefficients}{all-orders saddle coefficient; $a_j$; $A_j$}
After normalizing a saddle integral by its leading Gaussian factor, the remaining multiplicative expansion has \emph{mass coefficients} $a_j$.  Taking the formal logarithm produces \emph{logarithmic coefficients} $A_j$, so that
\[
 \log F(x)=\mathcal M(x)+\sum_{j=1}^{N-1}\frac{A_j(\lambda)}{\lambda^j}
   +O(\lambda^{-N}).
\]
Each $A_j$ is a bounded smooth one-periodic polynomial in derivatives of the basic correction $\Psi$ through order $2j$.  The highest derivative appears with coefficient $(-1)^j/(2^j j!(\log2)^{2j})$.  The distinction between $a_j$ and $A_j$ is the usual distinction between moments and cumulants: logarithms remove disconnected products.
\end{glossaryentry}
\begin{glossaryentry}{saddle-point}{Saddle point}{stationary point of an exponent; steepest-descent point}
A \emph{saddle point} of an exponent $\Phi(z)$ is a point $z_0$ with $\Phi'(z_0)=0$.  Along suitable directions the real part decreases away from $z_0$, so a contour through it captures the dominant contribution to an integral.  In the Fabius Bromwich integral, the saddle radius is $r=2^{\lambda(x)}$ and moves to infinity as $x\to0^+$.  Its defining equation is solved by the lower branch of Lambert $W$.  The second derivative determines the Gaussian scale; higher derivatives form saddle jets.  A quantitative saddle theorem must also control contour deformation, the central window, and the outer tail uniformly in the moving parameter.
\end{glossaryentry}
\begin{glossaryentry}{self-similarity}{Self-similarity}{discrete scale invariance; dyadic self-similarity}
A function or measure is \emph{self-similar} when rescaled pieces reproduce the whole according to a fixed rule.  The Fabius system has discrete scale invariance under factor two, visible in the derivative equation, the weighted-uniform decomposition, the sinc product, and exact dyadic recurrences.  Discrete rather than continuous scale invariance commonly creates periodic fluctuations in logarithmic coordinates.  Self-similarity also allows proofs to propagate local information across scales: support gives a base region, and dilation determines derivatives, signs, or values elsewhere.  The term does not imply fractal nonsmoothness; the up-function is infinitely smooth despite its exact dyadic recursion.
\end{glossaryentry}
\begin{glossaryentry}{signed-global-fabius}{Signed global Fabius extension}{extended Fabius function; global Fabius function; $\mathcal F$}
The \emph{signed global Fabius extension} is a real-valued continuation of the bounded Fabius function defined by a locally finite or absolutely convergent Thue--Morse translate sum of Rvachev's up-function.  It agrees with the bounded CDF on $[0,1]$ but may be negative or oscillatory outside that interval.  Unlike the clamped CDF, it satisfies the global dilation identity
\[
 \mathcal F'(x)=2\mathcal F(2x)
\]
for every real $x$, and therefore supports formulas involving arbitrary nonnegative dyadic arguments, translated $q$-binomial sums, and global binary-reduction series.  A statement using ``global'' in the probability API can instead mean ``valid for every real argument'' and need not refer to this signed extension; the repository explicitly separates the two meanings.
\end{glossaryentry}
\begin{glossaryentry}{signed-stirling}{Signed Stirling numbers of the first kind}{first-kind Stirling numbers; permutation-cycle coefficients}
The \emph{signed Stirling numbers of the first kind} $s(n,k)$ are defined by
\[
 x(x-1)\cdots(x-n+1)=\sum_{k=0}^n s(n,k)x^k.
\]
Their absolute values count permutations of $n$ elements with $k$ cycles.  Shifted versions occur when expanding $\prod_{j=1}^n(z-j)$.  The Fabius saddle-jet recurrence is solved by such shifted signed Stirling coefficients: repeated application of a first-order differential operator creates exactly the falling-factorial transition matrix.  They should not be confused with Stirling numbers of the second kind, which expand ordinary powers into falling factorials and count set partitions.
\end{glossaryentry}
\begin{glossaryentry}{sinc-function}{Sinc function}{cardinal sine; $\operatorname{sinc}$}
The \emph{sinc function} is $\sinc z=\sin z/z$ with the removable value $\sinc0=1$; some engineering conventions instead use $\sin(\pi z)/(\pi z)$.  It is the Fourier transform of a uniform interval density up to scaling and phase.  The up-function transform is an infinite product of sinc factors at dyadically shrinking frequencies.  Zeros at nonzero integers drive the partition-of-unity identity through Poisson summation.  The product identity $\sin z/z=\prod_{n\ge1}\cos(z/2^n)$ gives an equivalent cosine-power representation and exposes the multiplicity of transform zeros through $2$-adic valuations.
\end{glossaryentry}
\begin{glossaryentry}{standard-gaussian}{Standard Gaussian distribution}{standard normal law; $N(0,1)$}
The \emph{standard Gaussian distribution} has density $(2\pi)^{-1/2}e^{-x^2/2}$, mean zero, and variance one.  Its characteristic function is $e^{-t^2/2}$ and its moments are given by double factorials.  In saddle-point analysis it is the universal local model obtained after translating to a nondegenerate saddle and scaling by the square root of the second derivative.  The Fabius saddle expansion uses Gaussian expectation as an algebraic contraction of higher exponent terms.  This auxiliary Gaussian is not the probability distribution defining the Fabius function; it emerges from local asymptotic geometry of the inversion integral.
\end{glossaryentry}
\begin{glossaryentry}{stieltjes-constant}{Stieltjes constants}{generalized Euler constants; $\gamma_n$}
The \emph{Stieltjes constants} are the coefficients in the Laurent expansion of zeta at its pole,
\[
 \zeta(1+s)=\frac1s+\sum_{n=0}^\infty\frac{(-1)^n\gamma_n}{n!}s^n.
\]
Here $\gamma_0=\gamma$ is Euler's constant and $\gamma_1$ is the first Stieltjes constant.  When the Mellin identity $-\Gamma(s)\zeta(s+1)$ is expanded at $s=0$, the finite part depends on $\gamma_1$ as well as $\gamma$ and $\pi^2$.  The repository develops the required local expansion explicitly because the constant fixes the mean term in the sharp Fabius asymptotic.  Stieltjes constants should not be confused with Stirling numbers despite the similar names.
\end{glossaryentry}
\begin{glossaryentry}{stirling-asymptotic}{Stirling's asymptotic formula}{Stirling approximation}
\emph{Stirling's formula} gives
\[
 n!=\sqrt{2\pi n}\,(n/e)^n\bigl(1+O(1/n)\bigr),
\]
with a full asymptotic expansion in inverse powers of $n$.  Its logarithmic form is $\log n!=n\log n-n+\tfrac12\log(2\pi n)+O(1/n)$.  The corpus uses precise Stirling bounds to audit growth claims in a $k$-fold Thue--Morse draft and to estimate factorials and Gamma values.  Omitting a linear term in a logarithmic Stirling calculation can change an exponential-scale conclusion, so the formalization records the exact $O(\log n)$ remainder at the level needed by the argument.
\end{glossaryentry}
\begin{glossaryentry}{sturm-liouville}{Sturm--Liouville equation}{self-adjoint second-order eigenvalue problem}
A \emph{Sturm--Liouville equation} has the form $-(p(x)y')'+q(x)y=\lambda w(x)y$ and defines a self-adjoint eigenvalue problem under suitable boundary conditions.  Eigenfunctions for distinct eigenvalues are orthogonal with respect to the weight $w$.  Legendre polynomials satisfy
\[
 -\bigl((1-x^2)P_n'(x)\bigr)'=n(n+1)P_n(x)
\]
on $[-1,1]$.  The repository proves this equation and uses it as part of the classical Legendre package supporting the up-function expansion.  The vanishing factor $1-x^2$ at the endpoints makes the natural boundary terms disappear in orthogonality arguments.
\end{glossaryentry}
\begin{glossaryentry}{summability}{Summability of a family}{summable series; unconditional convergence}
A family $(a_i)_{i\in I}$ in a normed space is \emph{summable} when its finite subsums converge to a limit independent of enumeration.  For countable scalar families, absolute convergence implies summability; in finite-dimensional spaces the converse holds.  Formal libraries define summability through a net over finite subsets, making rearrangement invariance explicit.  The Fabius development establishes summability before using topological infinite sums for Fourier modes, translate formulas, and periodic corrections.  A bound by a summable majorant is the standard route, often via geometric or Gamma-decaying estimates.
\end{glossaryentry}
\begin{glossaryentry}{supremum}{Supremum and maximum}{least upper bound; exact norm}
The \emph{supremum} of a set is its least upper bound; it need not belong to the set.  A maximum is an attained supremum.  Function norms such as $\|f\|_\infty=\sup_x|f(x)|$ are therefore first upper-bound statements and only later attainment statements.  The repository proves exact derivative suprema by giving both the universal bound and a point where equality holds.  Compactness plus continuity often guarantees attainment, but the signed global functions live on all of $\mathbb R$, so support, periodicity, or explicit witnesses are used instead.  Calling a constant ``sharp'' without an attainment or limiting argument overclaims what a one-sided estimate proves.
\end{glossaryentry}
\begin{glossaryentry}{symmetry-reflection}{Symmetry and reflection identity}{even symmetry; midpoint symmetry; complementarity}
A function is even when $f(-x)=f(x)$.  A CDF-style midpoint reflection has the complementary form $F(x)+F(1-x)=1$.  Folding the bounded Fabius function about $1/2$ converts the latter into evenness of Rvachev's up-function.  Symmetry halves many calculations: odd moments and odd Legendre coefficients vanish, endpoint asymptotics transfer from left to right, and monotonicity on one half determines the other.  Structural evenness of the folded function can hold for every bounded candidate before the Fabius differential equations are assumed, whereas exact support and positivity require the full specification.
\end{glossaryentry}
\clearpage
\hypertarget{letter:T}{}
\section*{\color{linkblue}T}
\addcontentsline{toc}{section}{T}
\markboth{T}{T}
\alphabetbar
\begin{glossaryentry}{two-adic-valuation}{$2$-adic valuation}{dyadic valuation; $v_2$; $\nu_2$}
For a nonzero rational number $r$, the \emph{$2$-adic valuation} $v_2(r)$ is the exponent of $2$ in its prime factorization: $r=2^{v_2(r)}u/v$ with $u,v$ odd.  It satisfies $v_2(rs)=v_2(r)+v_2(s)$ and $v_2(r+s)\ge\min(v_2(r),v_2(s))$, with equality when the two valuations differ.  The arithmetic paper proves the exact formula
\[
 v_2\bigl(F(2^{-n})\bigr)=-\binom n2-v_2(n!)-1
\]
for the relevant normalization.  Parity lemmas, Lucas's theorem, and oddness of moment numerators are used to show that no hidden factor of two cancels.  The valuation measures only the power of two; odd prime factors require separate denominator arguments.
\end{glossaryentry}
\begin{glossaryentry}{taylor-theorem}{Taylor's theorem}{Taylor polynomial; integral remainder}
\emph{Taylor's theorem} approximates a sufficiently differentiable function near $a$ by
\[
 f(x)=\sum_{k=0}^n\frac{f^{(k)}(a)}{k!}(x-a)^k+R_n(x).
\]
The integral remainder is $R_n(x)=\int_a^x f^{(n+1)}(t)(x-t)^n/n!\,dt$.  In the dyadic Fabius evaluator, Taylor expansion at a reciprocal power of two expresses $F(x)$ as a finite polynomial plus a reflected smaller argument.  At dyadic points sufficiently high derivatives vanish, so certain Taylor series terminate exactly.  At the endpoint all coefficients vanish but the function is nonzero nearby, demonstrating that Taylor's theorem with a remainder does not imply analyticity.
\end{glossaryentry}
\begin{glossaryentry}{telescoping-series}{Telescoping series}{finite telescope; cancellation of consecutive terms}
A \emph{telescoping series} has terms expressible as $b_n-b_{n+1}$, so a finite sum collapses to $b_0-b_N$.  If $b_N\to0$, the infinite sum equals $b_0$.  The global binary-reduction and $q$-binomial developments prove exact finite-remainder telescope identities before passing to infinite sums.  This is stronger and safer than guessing a limit from numerical partial sums: the only analytic task left is to bound the explicit remainder.  Reindexing can shift the initial boundary term, which is why the correct global outer index begins at zero.
\end{glossaryentry}
\begin{glossaryentry}{termwise-operations}{Termwise integration and differentiation}{interchange of series with integral or derivative}
A series may be integrated or differentiated \emph{termwise} only under hypotheses that control convergence.  Uniform convergence permits integration on compact intervals; dominated convergence permits integration under an integrable majorant; uniform convergence of the derivative series together with convergence at one point permits differentiation.  The Mellin identity integrates the logarithmic Bose exponential series term by term after proving absolute integrability of the norms.  Fourier derivatives of the periodic correction require sufficient weighted coefficient summability.  Formal power-series differentiation is always coefficientwise algebra, but it becomes an analytic derivative only inside a common convergence disk.
\end{glossaryentry}
\begin{glossaryentry}{thue-morse}{Thue--Morse sequence}{Prouhet--Thue--Morse sequence; binary-parity sequence}
The \emph{Thue--Morse sequence} is $t_n=s_2(n)\bmod2$, where $s_2(n)$ is the number of ones in the binary expansion of $n$.  Its sign form is $\varepsilon_n=(-1)^{s_2(n)}$.  It satisfies $\varepsilon_{2n}=\varepsilon_n$ and $\varepsilon_{2n+1}=-\varepsilon_n$, so each dyadic block is followed by its negation.  In the Fabius theory these signs build the signed global extension, express higher derivatives as translate sums, annihilate low-degree polynomials, and enter exact dyadic and $q$-binomial formulas.  The zero-one form and sign form encode the same sequence but behave differently inside powers and logarithms, so the formalization exposes both.
\end{glossaryentry}
\begin{glossaryentry}{toeplitz-summation}{Toeplitz summation theorem}{Toeplitz averaging; regular summability matrix}
A \emph{Toeplitz summation theorem} gives conditions under which weighted averages of a convergent sequence have the same limit.  For a matrix $(a_{n,k})$, typical conditions are uniformly bounded row $\ell^1$ norms, $a_{n,k}\to0$ for each fixed $k$, and row sums tending to one.  In the Fabius discrete-limit proof, reindexed finite rows become Toeplitz averages of centered spline approximants.  Bounded variation or explicit row-control estimates verify regularity, allowing convergence of the underlying approximants to pass to the weighted sums.  This proves limit independence of a complex shift even when individual finite rows depend on it.
\end{glossaryentry}
\begin{glossaryentry}{translation-invariance}{Translation invariance}{shift invariance; common-shift independence}
A quantity is \emph{translation invariant} when shifting every argument by the same amount leaves it unchanged.  Lebesgue measure and convolution have standard translation invariance.  More unusually, the finite Fabius $q$-binomial expressions are constant polynomials in a common translation parameter.  This follows because they extract a leading coefficient, which is unchanged by replacing $p(x)$ with $p(x+q)$.  The theorem extends the identity from rational shifts to arbitrary real or complex shifts.  A finite discrete-limit row need not share this invariance; in that setting only the limit is shift independent.
\end{glossaryentry}
\begin{glossaryentry}{triangular-cancellation}{Triangular cancellation}{vanishing of lower degrees; moment cancellation}
A \emph{triangular cancellation} is a family of identities in which an order-$n$ signed sum annihilates every polynomial of degree below $n$ and has a controlled value at degree $n$.  The resulting matrix from degrees to sums is triangular with nonzero diagonal, hence invertible.  Thue--Morse block sums provide one such cancellation, and geometric-node Gaussian-binomial weights provide another.  Their combination is the algebraic heart of the finite Fabius formulas: one mechanism removes low powers within blocks, while the other extracts the remaining leading coefficient across scales.  The word ``triangular'' refers to degree filtration, not geometric triangles.
\end{glossaryentry}
\clearpage
\hypertarget{letter:U}{}
\section*{\color{linkblue}U}
\addcontentsline{toc}{section}{U}
\markboth{U}{U}
\alphabetbar
\begin{glossaryentry}{uniform-convergence}{Uniform convergence}{convergence in the supremum norm}
A sequence $f_n$ converges \emph{uniformly} to $f$ on $E$ when $\sup_{x\in E}|f_n(x)-f(x)|\to0$.  Unlike pointwise convergence, uniform convergence preserves continuity and permits passage of limits through integrals on finite-measure domains.  Uniform convergence of the Fourier--Legendre expansion on $[-1,1]$ makes the identity valid at both endpoints.  Centered finite splines approximate the bounded Fabius function with an explicit uniform bound, which is crucial for computability.  Uniform convergence on every compact set is weaker than global uniform convergence but is the natural mode for infinite products and holomorphic series.
\end{glossaryentry}
\begin{glossaryentry}{uniform-distribution}{Uniform distribution}{continuous uniform law; rectangular density}
A random variable is \emph{uniform on $[a,b]$} when it has constant density $1/(b-a)$ there and zero outside.  Its characteristic function is an exponential phase times a sinc factor.  The Fabius random variable is an infinite weighted sum of independent uniform $[0,1]$ coordinates.  Finite sums have spline densities, and multiplication of their sinc characteristic functions yields the transform product.  ``Uniform distribution'' can also mean equidistribution of a deterministic sequence modulo one; the Fabius corpus primarily uses the probabilistic meaning, except where finite uniform-distribution CDFs appear as spline approximants.
\end{glossaryentry}
\begin{glossaryentry}{unimodality}{Unimodality}{single-peaked function; strict unimodality}
A nonnegative function is \emph{unimodal} if it increases up to a mode and decreases afterward; strict unimodality requires strict behavior away from the mode.  Rvachev's up-function is even and strictly unimodal with its unique maximum $1$ at zero.  This follows from the positivity of the appropriate translated Fabius term in its derivative equation.  The bounded Fabius function is instead sigmoid-shaped: it is strictly increasing, convex before the midpoint, concave after it, and has a unique inflection point at $1/2$.  Unimodality is a shape property and does not by itself imply log-concavity.
\end{glossaryentry}
\clearpage
\hypertarget{letter:V}{}
\section*{\color{linkblue}V}
\addcontentsline{toc}{section}{V}
\markboth{V}{V}
\alphabetbar
\begin{glossaryentry}{vertical-contour}{Vertical Bromwich contour}{inverse-Laplace contour; vertical line of integration}
A \emph{vertical Bromwich contour} is a line $z=c+it$ in the complex plane used to invert a Laplace transform:
\[
 f(x)=\frac1{2\pi i}\int_{c-i\infty}^{c+i\infty}e^{zx}F(z)\,dz.
\]
The real number $c$ lies to the right of the transform's singularities.  In the Fabius problem the transform is entire but the contour is moved to pass through a positive saddle chosen by $x$.  Bounds on the transform along this vertical line control the tails, while a local expansion in $t$ produces the Gaussian central contribution.  Orientation, the factor $1/(2\pi i)$, and whether the inversion targets a density or a cumulative distribution affect the amplitude.
\end{glossaryentry}
\clearpage
\hypertarget{letter:W}{}
\section*{\color{linkblue}W}
\addcontentsline{toc}{section}{W}
\markboth{W}{W}
\alphabetbar
\begin{glossaryentry}{weighted-sum}{Infinite weighted sum of random variables}{random series}
A \emph{random series} is a sum $\sum_{k\ge1}a_kX_k$ of random variables.  If $|X_k|\le1$ and $\sum|a_k|<\infty$, it converges absolutely for every outcome.  The Fabius variable uses $a_k=2^{-k}$ and independent uniforms $U_k\in[0,1]$, so the sum lies in $[0,1]$.  Splitting off the first coordinate gives a distributional fixed-point equation, while finite partial sums give convolution and spline approximants.  The geometric weights are responsible for exact dyadic self-similarity; replacing them by a general summable sequence would retain a smooth convolution law but usually lose the special differential and arithmetic identities.
\end{glossaryentry}
\begin{glossaryentry}{wavelet}{Wavelet and scaling function}{multiresolution function; refinement basis}
A \emph{wavelet} is a localized function whose dilates and translates analyze a signal or function at multiple scales.  A scaling function satisfies a refinement equation and generates nested approximation spaces; a wavelet spans the complementary detail spaces.  Rvachev's up-function and related atomic functions have compact support, smoothness, and exact refinement identities, making them natural ingredients in wavelet-like and approximation constructions.  The corpus cites applications to orthonormal systems and Meyer-wavelet constructions but does not identify the up-function itself with every technical notion of an orthonormal wavelet.  Orthogonality, normalization, and multiresolution completeness require additional theorems beyond a two-scale relation.
\end{glossaryentry}
\begin{glossaryentry}{weak-star-convergence}{Weak-* convergence}{weak-star convergence of measures; vague dual convergence}
\emph{Weak-* convergence} views measures as continuous linear functionals on a space of test functions and asks for pointwise convergence on that dual pairing.  For finite measures on a compact space it coincides with ordinary weak convergence against continuous functions.  The first Fabius source describes finite convolution measures converging weak-* to the measure with the sinc-product transform.  The star emphasizes that convergence is weak in a dual space rather than norm convergence of measures.  On noncompact spaces, the choice between bounded continuous, compactly supported continuous, or vanishing-at-infinity test functions affects the exact notion.
\end{glossaryentry}
\begin{glossaryentry}{weak-convergence}{Weak convergence of probability measures}{convergence in distribution; weak convergence of laws}
Probability measures $\mu_n$ converge \emph{weakly} to $\mu$ when $\int f\,d\mu_n\to\int f\,d\mu$ for every bounded continuous $f$.  Equivalently on $\mathbb R$, their cumulative distribution functions converge at every continuity point of the limit; characteristic functions may also identify the limit through L\'evy's theorem.  The finite convolution laws of weighted uniform partial sums converge weakly to the Fabius law.  Weak convergence does not imply convergence of densities pointwise, although additional spline estimates provide that stronger conclusion here.  It preserves total mass but not necessarily moments without uniform integrability.
\end{glossaryentry}
\begin{glossaryentry}{weierstrass-product}{Weierstrass product}{canonical product over zeros}
A \emph{Weierstrass product} represents an entire function by factors encoding its zeros, augmented when necessary by exponential convergence factors.  The identity
\[
 \frac{\sin\pi z}{\pi z}=\prod_{n=1}^\infty\left(1-\frac{z^2}{n^2}\right)
\]
is the standard example.  Substituting it into the dyadic sinc product expresses the up-transform as a product whose zero at an integer $m$ has multiplicity $1+v_2(m)$.  This zero structure explains partition identities and sign patterns of half-integer Fourier coefficients.  A zero product determines an entire function only up to an exponential factor unless normalization and growth conditions fix that ambiguity.
\end{glossaryentry}
\begin{glossaryentry}{well-founded-recursion}{Well-founded recursion}{terminating recursion; decreasing-measure recursion}
\emph{Well-founded recursion} defines a value by recursive calls on arguments smaller in a relation that admits no infinite descending chain.  It generalizes primitive recursion on natural numbers.  The exact dyadic Fabius evaluator removes the highest set bit of the numerator at each step, so binary weight or numerator size supplies a decreasing measure and proves termination.  Formalizing this measure separates mathematical termination from operational intuition.  Once termination is established, induction over the same well-founded relation proves correctness and complexity properties.  The resulting recursion is finite even when the original function is defined by infinite convolution or a fixed point.
\end{glossaryentry}
\clearpage
\hypertarget{letter:Z}{}
\section*{\color{linkblue}Z}
\addcontentsline{toc}{section}{Z}
\markboth{Z}{Z}
\alphabetbar
\begin{glossaryentry}{zero-run}{Zero run}{consecutive zeros; vanishing block}
A \emph{zero run} is a consecutive block of zero terms in a sequence.  Iterated prefix sums of Thue--Morse signs develop structured zero runs because complete dyadic blocks have vanishing low moments.  The $k$-fold Thue--Morse part of the corpus proves exact lengths and locations of these runs and uses them to analyze polygonal interpolation schemes.  Zero runs are discrete manifestations of cancellation, not evidence that the limiting smooth function vanishes on corresponding intervals.  Indexing and normalization are especially delicate at the beginning of a run, where an omitted initial term can shift every stated boundary.
\end{glossaryentry}
\clearpage
\hypertarget{alias-index}{}
\section*{Alias and synonym index}
\addcontentsline{toc}{section}{Alias and synonym index}
\markboth{Alias and synonym index}{Alias and synonym index}
Alternative names, common abbreviations, notation, and spelling variants point
to the consolidated headword.  Entries are sorted by the printed alias rather
than by their TeX source.
\begin{multicols}{2}
\raggedcolumns
\raggedright
\small

\aliasindexentry{$\gamma$}{euler-mascheroni}{Euler--Mascheroni constant}
\aliasindexentry{$\Gamma$}{gamma-function}{Gamma function}
\aliasindexentry{$\Psi$}{periodic-correction}{Periodic correction in a self-similar asymptotic}
\aliasindexentry{$\Up$}{rvachev-up-function}{Rvachev's up-function}
\aliasindexentry{$\log(1-e^{-x})$}{bose-kernel}{Logarithmic Bose kernel}
\aliasindexentry{$\nu_2$}{two-adic-valuation}{$2$-adic valuation}
\aliasindexentry{$2$-automatic sequence}{automatic-sequence}{Automatic sequence}
\aliasindexentry{$(2n+1)!!$}{double-factorial}{Double factorial}
\aliasindexentry{a.e.}{almost-everywhere}{Almost everywhere}
\aliasindexentry{$a_j$}{saddle-coefficients}{Saddle mass and logarithmic coefficients}
\aliasindexentry{$(a;q)_n$}{q-pochhammer}{$q$-Pochhammer symbol}
\aliasindexentry{absolute integrability}{integrability}{Integrability}
\aliasindexentry{absolutely continuous function}{absolute-continuity}{Absolute continuity}
\aliasindexentry{absolutely continuous measure}{absolute-continuity}{Absolute continuity}
\aliasindexentry{absolutely convergent series}{absolute-convergence}{Absolute convergence}
\aliasindexentry{abstract Fabius candidate}{is-fabius-specification}{The \texorpdfstring{\texttt{IsFabius}}{IsFabius} specification}
\aliasindexentry{affine change of variable}{dilation}{Dilation and translation}
\aliasindexentry{all-orders expansion}{asymptotic-expansion}{Asymptotic expansion}
\aliasindexentry{all-orders saddle coefficient}{saddle-coefficients}{Saddle mass and logarithmic coefficients}
\aliasindexentry{almost-everywhere equality}{almost-everywhere}{Almost everywhere}
\aliasindexentry{analytic at a point}{analytic-function}{Analytic function and analytic locus}
\aliasindexentry{\texttt{AnalyticAt}}{analytic-function}{Analytic function and analytic locus}
\aliasindexentry{approximation error}{numerical-error}{Absolute, relative, and truncation error}
\aliasindexentry{asymptotic equivalent}{asymptotic-equivalent}{Asymptotic equivalence}
\aliasindexentry{asymptotic filter}{filter-convergence}{Filter convergence}
\aliasindexentry{asymptotic order}{big-o-little-o}{Big-$O$ and little-$o$}
\aliasindexentry{average over one period}{mean-periodic}{Period mean and centered periodic function}
\aliasindexentry{averaging kernel}{approximate-identity}{Approximate identity}
\aliasindexentry{base-two logarithmic coordinate}{logarithmic-scale}{Logarithmic scale and log-periodicity}
\aliasindexentry{base-two representation}{binary-expansion}{Binary expansion, set bit, and highest set bit}
\aliasindexentry{basic factorial}{q-integer-factorial}{$q$-integer and $q$-factorial}
\aliasindexentry{basic hypergeometric calculus}{q-calculus}{$q$-calculus}
\aliasindexentry{Bernoulli recurrence}{bernoulli-numbers}{Bernoulli numbers}
\aliasindexentry{bilateral Laplace transform of a law}{moment-generating-function}{Moment-generating function}
\aliasindexentry{binary mesh}{dyadic-grid}{Dyadic grid}
\aliasindexentry{binary-parity sequence}{thue-morse}{Thue--Morse sequence}
\aliasindexentry{binary rational}{dyadic-rational}{Dyadic rational}
\aliasindexentry{binomial coefficients modulo a prime}{congruence-lucas}{Congruence and Lucas's theorem}
\aliasindexentry{binomial-parity power series}{parity-power-series}{Parity-power form}
\aliasindexentry{bit-removal expansion}{binary-reduction}{Binary-reduction formula}
\aliasindexentry{Borel measurable}{measurable-function}{Measurable function}
\aliasindexentry{Bose--Mellin kernel}{bose-kernel}{Logarithmic Bose kernel}
\aliasindexentry{boundary conventions}{empty-conventions}{Empty sum, empty product, and $0^0$ conventions}
\aliasindexentry{boundary-layer asymptotic}{endpoint-asymptotic}{Endpoint asymptotic}
\aliasindexentry{bounded derivatives}{derivative-bound}{Uniform derivative bound}
\aliasindexentry{bounded Fabius axioms}{is-fabius-specification}{The \texorpdfstring{\texttt{IsFabius}}{IsFabius} specification}
\aliasindexentry{bounded Fabius function}{fabius-function}{Bounded Fabius function}
\aliasindexentry{box spline}{b-spline}{B-spline}
\aliasindexentry{Bromwich contour}{bromwich-inversion}{Bromwich inversion formula}
\aliasindexentry{$C^\infty$}{regularity}{Regularity of a function}
\aliasindexentry{$C^k$}{regularity}{Regularity of a function}
\aliasindexentry{$c_n$}{moment}{Moment of a distribution}
\aliasindexentry{cancellation of consecutive terms}{telescoping-series}{Telescoping series}
\aliasindexentry{canonical Fabius CDF}{fabius-function}{Bounded Fabius function}
\aliasindexentry{canonical normalization}{normalization}{Normalization}
\aliasindexentry{canonical product over zeros}{weierstrass-product}{Weierstrass product}
\aliasindexentry{cardinal partition}{partition-of-unity}{Partition of unity}
\aliasindexentry{cardinal sine}{sinc-function}{Sinc function}
\aliasindexentry{CDF}{cumulative-distribution}{Cumulative distribution function}
\aliasindexentry{centered periodic correction}{centering}{Centering and mean-zero normalization}
\aliasindexentry{centered uniform spline}{centered-finite-spline}{Centered finite spline}
\aliasindexentry{central arc}{major-minor-arc}{Major arc and minor arc}
\aliasindexentry{characteristic function of a set}{indicator-function}{Indicator function}
\aliasindexentry{characteristic transform}{fourier-transform}{Fourier transform}
\aliasindexentry{choice of scale}{normalization}{Normalization}
\aliasindexentry{classical orthogonal polynomial}{legendre-polynomial}{Legendre polynomial}
\aliasindexentry{coefficientwise series}{formal-power-series}{Formal power series}
\aliasindexentry{combinatorial number}{binomial-coefficient}{Binomial coefficient}
\aliasindexentry{common-shift independence}{translation-invariance}{Translation invariance}
\aliasindexentry{common-shift invariance}{constant-polynomial}{Constant polynomial and translation independence}
\aliasindexentry{compact convergence}{compact-uniform-convergence}{Uniform convergence on compact sets}
\aliasindexentry{complementarity}{symmetry-reflection}{Symmetry and reflection identity}
\aliasindexentry{complex line integral}{contour-integral}{Contour integral}
\aliasindexentry{composition}{ordered-composition}{Ordered composition of an integer}
\aliasindexentry{computable modulus}{effective-uniform-continuity}{Effective uniform continuity}
\aliasindexentry{computable real function in the Grzegorczyk sense}{grzegorczyk-computability}{Grzegorczyk computability of a real function}
\aliasindexentry{consecutive zeros}{zero-run}{Zero run}
\aliasindexentry{constant term at a pole}{finite-part}{Finite part of a singular expression}
\aliasindexentry{constraint set}{admissible-class}{Admissible class}
\aliasindexentry{\texttt{ContDiff}}{smoothness}{$C^k$ and $C^\infty$ smoothness}
\aliasindexentry{continuity class}{regularity}{Regularity of a function}
\aliasindexentry{continuity modulus}{modulus-of-continuity}{Modulus of continuity}
\aliasindexentry{continuity theorem for characteristic functions}{levy-continuity}{L\'evy's continuity theorem}
\aliasindexentry{continuous uniform law}{uniform-distribution}{Uniform distribution}
\aliasindexentry{continuously differentiable}{smoothness}{$C^k$ and $C^\infty$ smoothness}
\aliasindexentry{continuum limit of finite sums}{discrete-limit}{Discrete limit}
\aliasindexentry{contraction constant}{contraction}{Contraction}
\aliasindexentry{contraction mapping theorem}{banach-fixed-point}{Banach fixed-point theorem}
\aliasindexentry{convergence at each point}{pointwise-convergence}{Pointwise convergence}
\aliasindexentry{convergence in distribution}{weak-convergence}{Weak convergence of probability measures}
\aliasindexentry{convergence in the supremum norm}{uniform-convergence}{Uniform convergence}
\aliasindexentry{convergence with probability one}{almost-sure-convergence}{Almost-sure convergence}
\aliasindexentry{convergent product}{infinite-product}{Infinite product}
\aliasindexentry{convolution measure}{convolution-measures}{Finite and infinite convolution of measures}
\aliasindexentry{convolution power}{convolution}{Convolution}
\aliasindexentry{convolution product}{convolution-measures}{Finite and infinite convolution of measures}
\aliasindexentry{convolution product of measures}{infinite-convolution}{Infinite convolution}
\aliasindexentry{cumulant coefficient}{cumulant}{Cumulant}
\aliasindexentry{cumulative sum}{prefix-sum}{Prefix sum and iterated prefix sum}
\aliasindexentry{$D_n$}{dyadic-denominator}{Dyadic common denominator}
\aliasindexentry{$d_n$}{half-moment}{Half moment}
\aliasindexentry{de Bruijn-type scale}{quadratic-logarithmic-law}{Quadratic logarithmic decay law}
\aliasindexentry{de Polignac formula}{legendre-factorial-formula}{Legendre's formula for factorial valuations}
\aliasindexentry{decreasing-measure recursion}{well-founded-recursion}{Well-founded recursion}
\aliasindexentry{delta measure}{dirac-measure}{Dirac measure}
\aliasindexentry{denominator sequence}{dyadic-denominator}{Dyadic common denominator}
\aliasindexentry{density function}{density}{Probability density}
\aliasindexentry{derivative definition of Legendre polynomials}{rodrigues-formula}{Rodrigues' formula}
\aliasindexentry{derivative supremum}{derivative-bound}{Uniform derivative bound}
\aliasindexentry{differentiability class}{regularity}{Regularity of a function}
\aliasindexentry{differential identity}{differential-equation}{Differential equation}
\aliasindexentry{dilation equation}{refinement-equation}{Refinement equation}
\aliasindexentry{discrete B-spline law}{b-spline}{B-spline}
\aliasindexentry{discrete derivative}{finite-difference}{Finite difference}
\aliasindexentry{discrete scale invariance}{self-similarity}{Self-similarity}
\aliasindexentry{distribution function}{cumulative-distribution}{Cumulative distribution function}
\aliasindexentry{distribution of a random variable}{probability-law}{Probability law}
\aliasindexentry{division-free numerator sequences}{moment-numerators}{Moment and half-moment numerators}
\aliasindexentry{dyadic point}{dyadic-rational}{Dyadic rational}
\aliasindexentry{dyadic refinement invariance}{refinement-invariance}{Refinement and representation invariance}
\aliasindexentry{dyadic self-similarity}{self-similarity}{Self-similarity}
\aliasindexentry{dyadic valuation}{two-adic-valuation}{$2$-adic valuation}
\aliasindexentry{EGF}{exponential-generating-function}{Exponential generating function}
\aliasindexentry{endpoint-corrected indicator}{half-endpoint-indicator}{Half-weighted interval indicator}
\aliasindexentry{\texttt{ENNReal}}{extended-nonnegative-reals}{Extended nonnegative reals}
\aliasindexentry{entire complex function}{entire-function}{Entire function}
\aliasindexentry{entire function of exponential type}{exponential-type}{Exponential type}
\aliasindexentry{entire moment-generating series}{complex-egf}{Complex exponential generating function}
\aliasindexentry{equal-power partition}{prouhet-thue-morse}{Prouhet--Thue--Morse cancellation}
\aliasindexentry{equality near a point}{germ}{Germ of a function}
\aliasindexentry{Euler gamma function}{gamma-function}{Gamma function}
\aliasindexentry{Euler's constant}{euler-mascheroni}{Euler--Mascheroni constant}
\aliasindexentry{evaluator error}{numerical-error}{Absolute, relative, and truncation error}
\aliasindexentry{even integer}{parity}{Parity, evenness, and oddness}
\aliasindexentry{even symmetry}{symmetry-reflection}{Symmetry and reflection identity}
\aliasindexentry{exact arithmetic recurrence}{rational-recurrence}{Rational recurrence}
\aliasindexentry{exact norm}{supremum}{Supremum and maximum}
\aliasindexentry{exact reproduction}{polynomial-reproduction}{Polynomial reproduction}
\aliasindexentry{exact supremum}{attained-bound}{Attained bound and sharp constant}
\aliasindexentry{exponential formula}{bell-polynomial}{Bell-polynomial and exponential-coefficient recurrence}
\aliasindexentry{extended Fabius function}{signed-global-fabius}{Signed global Fabius extension}
\aliasindexentry{$\mathcal F$}{signed-global-fabius}{Signed global Fabius extension}
\aliasindexentry{$F(2^{-n})$}{inverse-dyadic-value}{Inverse-dyadic value}
\aliasindexentry{$f\sim g$}{asymptotic-equivalent}{Asymptotic equivalence}
\aliasindexentry{$f^{(n)}$}{iterated-derivative}{Iterated derivative}
\aliasindexentry{$F_n$ and $G_n$ sequences}{moment-numerators}{Moment and half-moment numerators}
\aliasindexentry{\texttt{fabiusReal}}{fabius-function}{Bounded Fabius function}
\aliasindexentry{fast Cauchy name}{computable-real-sequence}{Computable real number and computable real sequence}
\aliasindexentry{fast dyadic name}{computable-real-sequence}{Computable real number and computable real sequence}
\aliasindexentry{finite convolution approximant}{polynomial-step-approximant}{Polynomial step approximant}
\aliasindexentry{finite function}{atomic-function}{Atomic function}
\aliasindexentry{finite function}{rvachev-up-function}{Rvachev's up-function}
\aliasindexentry{finite jet}{derivative-jet}{Derivative jet}
\aliasindexentry{finite-spline approximant}{centered-finite-spline}{Centered finite spline}
\aliasindexentry{finite telescope}{telescoping-series}{Telescoping series}
\aliasindexentry{first-kind Stirling numbers}{signed-stirling}{Signed Stirling numbers of the first kind}
\aliasindexentry{first saddle correction}{saddle-correction}{Saddle correction}
\aliasindexentry{fixed-point equation}{fixed-point}{Fixed point}
\aliasindexentry{formal exponential recurrence}{bell-polynomial}{Bell-polynomial and exponential-coefficient recurrence}
\aliasindexentry{forward difference}{finite-difference}{Finite difference}
\aliasindexentry{Fourier mode}{fourier-coefficient}{Fourier coefficient}
\aliasindexentry{Fourier support--growth theorem}{paley-wiener}{Paley--Wiener theorem}
\aliasindexentry{Fourier transform of a probability law}{characteristic-function}{Characteristic function}
\aliasindexentry{FPS}{formal-power-series}{Formal power series}
\aliasindexentry{frequency transform}{fourier-transform}{Fourier transform}
\aliasindexentry{FTC}{fundamental-theorem-calculus}{Fundamental theorem of calculus}
\aliasindexentry{functional--differential equation}{delay-equation}{Delay and dilation equation}
\aliasindexentry{Gamma--zeta mode}{gamma-zeta-coefficient}{Gamma--zeta Fourier coefficient}
\aliasindexentry{Gaussian binomial theorem}{q-binomial-theorem}{Finite $q$-binomial theorem}
\aliasindexentry{Gaussian coefficient}{gaussian-binomial}{Gaussian binomial coefficient}
\aliasindexentry{Gaussian expectation operator}{gaussian-contraction}{Gaussian contraction}
\aliasindexentry{generalized Euler constants}{stieltjes-constant}{Stieltjes constants}
\aliasindexentry{generalized function}{distribution-generalized}{Distribution in the sense of generalized functions}
\aliasindexentry{geometric progression of nodes}{geometric-nodes}{Geometric interpolation nodes}
\aliasindexentry{global binary recursion}{binary-reduction}{Binary-reduction formula}
\aliasindexentry{global Fabius function}{signed-global-fabius}{Signed global Fabius extension}
\aliasindexentry{greatest-integer function}{floor-function}{Floor and ceiling functions}
\aliasindexentry{grid of order $n$}{dyadic-grid}{Dyadic grid}
\aliasindexentry{Grzegorczyk computability}{computable-real-function}{Computable real function}
\aliasindexentry{$H_n$}{harmonic-number}{Harmonic number}
\aliasindexentry{Hadamard finite part}{finite-part}{Finite part of a singular expression}
\aliasindexentry{half-endpoint convention}{half-endpoint-indicator}{Half-weighted interval indicator}
\aliasindexentry{Hamming weight}{binary-weight}{Binary weight}
\aliasindexentry{harmonic coefficient}{fourier-coefficient}{Fourier coefficient}
\aliasindexentry{higher derivative}{iterated-derivative}{Iterated derivative}
\aliasindexentry{higher-order steepest-descent term}{saddle-correction}{Saddle correction}
\aliasindexentry{holomorphic continuation}{analytic-continuation}{Analytic continuation}
\aliasindexentry{homothetic copy}{dilation}{Dilation and translation}
\aliasindexentry{$Icc$, $Ioo$, $Ioc$, $Ico$}{intervals}{Closed, open, and half-open intervals}
\aliasindexentry{image measure}{pushforward-measure}{Pushforward measure}
\aliasindexentry{increasing function}{monotonicity}{Monotonicity}
\aliasindexentry{independent coordinates}{product-measure-independence}{Product measure and independence}
\aliasindexentry{infinite sum law}{infinite-convolution}{Infinite convolution}
\aliasindexentry{infinitely differentiable}{smoothness}{$C^k$ and $C^\infty$ smoothness}
\aliasindexentry{infinitely flat}{flat-function}{Flat function}
\aliasindexentry{integral--derivative theorem}{fundamental-theorem-calculus}{Fundamental theorem of calculus}
\aliasindexentry{integral over an unbounded or singular interval}{improper-integral}{Improper integral}
\aliasindexentry{integral remainder}{taylor-theorem}{Taylor's theorem}
\aliasindexentry{interchange of series with integral or derivative}{termwise-operations}{Termwise integration and differentiation}
\aliasindexentry{interpolation polynomial}{interpolation}{Polynomial interpolation}
\aliasindexentry{invariant class}{admissible-class}{Admissible class}
\aliasindexentry{invariant point}{fixed-point}{Fixed point}
\aliasindexentry{inverse-Laplace contour}{vertical-contour}{Vertical Bromwich contour}
\aliasindexentry{inverse Laplace integral}{bromwich-inversion}{Bromwich inversion formula}
\aliasindexentry{inversion formula}{fourier-inversion}{Fourier inversion}
\aliasindexentry{iterative solution of the saddle equation}{newton-iteration}{Fixed-point and Newton iteration for the Lambert phase}
\aliasindexentry{jet}{derivative-jet}{Derivative jet}
\aliasindexentry{$L^2$ approximation}{least-squares}{Least-squares approximation}
\aliasindexentry{Lagrange interpolation}{interpolation}{Polynomial interpolation}
\aliasindexentry{Lagrange mean value theorem}{mean-value-theorem}{Mean value theorem}
\aliasindexentry{Landau notation}{big-o-little-o}{Big-$O$ and little-$o$}
\aliasindexentry{LCM}{least-common-multiple}{Least common multiple}
\aliasindexentry{least constant}{attained-bound}{Attained bound and sharp constant}
\aliasindexentry{least-integer function}{floor-function}{Floor and ceiling functions}
\aliasindexentry{least Lipschitz constant}{lipschitz-continuity}{Lipschitz continuity}
\aliasindexentry{least upper bound}{supremum}{Supremum and maximum}
\aliasindexentry{Lebesgue dominated convergence}{dominated-convergence}{Dominated convergence theorem}
\aliasindexentry{Lebesgue integrability}{integrability}{Integrability}
\aliasindexentry{Legendre expansion}{fourier-legendre-series}{Fourier--Legendre series}
\aliasindexentry{Lipschitz bound}{lipschitz-continuity}{Lipschitz continuity}
\aliasindexentry{local germ}{germ}{Germ of a function}
\aliasindexentry{locally finite family}{local-finiteness}{Local finiteness}
\aliasindexentry{locally uniform convergence}{compact-uniform-convergence}{Uniform convergence on compact sets}
\aliasindexentry{log moment-generating function}{cumulant-generating}{Cumulant-generating function}
\aliasindexentry{log-periodic fluctuation}{periodic-correction}{Periodic correction in a self-similar asymptotic}
\aliasindexentry{log-squared asymptotic}{quadratic-logarithmic-law}{Quadratic logarithmic decay law}
\aliasindexentry{log transform}{cumulant-generating}{Cumulant-generating function}
\aliasindexentry{logarithmic product}{negative-laplace-product}{Negative-Laplace product}
\aliasindexentry{map of a measure}{pushforward-measure}{Pushforward measure}
\aliasindexentry{mean of a random variable}{expectation-variance}{Expectation and variance}
\aliasindexentry{Mellin--Fourier coefficient}{gamma-zeta-coefficient}{Gamma--zeta Fourier coefficient}
\aliasindexentry{meromorphic expansion}{laurent-series}{Laurent series and pole}
\aliasindexentry{MGF}{moment-generating-function}{Moment-generating function}
\aliasindexentry{midpoint symmetry}{symmetry-reflection}{Symmetry and reflection identity}
\aliasindexentry{moment cancellation}{triangular-cancellation}{Triangular cancellation}
\aliasindexentry{moment-generating transform}{laplace-transform}{Laplace transform}
\aliasindexentry{moment sequence}{moment}{Moment of a distribution}
\aliasindexentry{monotone map}{monotonicity}{Monotonicity}
\aliasindexentry{most significant bit}{binary-expansion}{Binary expansion, set bit, and highest set bit}
\aliasindexentry{multiplicative Fourier transform}{mellin-transform}{Mellin transform}
\aliasindexentry{multiplicative periodicity}{logarithmic-scale}{Logarithmic scale and log-periodicity}
\aliasindexentry{multiresolution function}{wavelet}{Wavelet and scaling function}
\aliasindexentry{$n!$}{factorial}{Factorial}
\aliasindexentry{$\gamma_n$}{stieltjes-constant}{Stieltjes constants}
\aliasindexentry{$N(0,1)$}{standard-gaussian}{Standard Gaussian distribution}
\aliasindexentry{neighborhood filter}{filter-convergence}{Filter convergence}
\aliasindexentry{nested polynomial evaluation}{horner-method}{Horner's method}
\aliasindexentry{neutral-element convention}{empty-conventions}{Empty sum, empty product, and $0^0$ conventions}
\aliasindexentry{nonanalytic at every point}{nowhere-analytic}{Nowhere analytic on a set}
\aliasindexentry{normal moment}{gaussian-moment}{Gaussian moment}
\aliasindexentry{numerical approximation}{approximation-theory}{Approximation theory}
\aliasindexentry{odd factorial}{double-factorial}{Double factorial}
\aliasindexentry{odd integer}{parity}{Parity, evenness, and oddness}
\aliasindexentry{one-periodic function}{periodic-function}{Periodic function}
\aliasindexentry{one-sided Laplace transform}{laplace-transform}{Laplace transform}
\aliasindexentry{one-sided moment}{half-moment}{Half moment}
\aliasindexentry{one-to-one correspondence}{bijection-inverse}{Bijection and inverse function}
\aliasindexentry{optimal bound}{attained-bound}{Attained bound and sharp constant}
\aliasindexentry{ordinary differential equation}{differential-equation}{Differential equation}
\aliasindexentry{ordinary generating function}{generating-function}{Generating function}
\aliasindexentry{orthogonal functions}{orthogonality}{Orthogonality}
\aliasindexentry{orthogonal-polynomial expansion}{fourier-legendre-series}{Fourier--Legendre series}
\aliasindexentry{orthogonal projection}{least-squares}{Least-squares approximation}
\aliasindexentry{$P_n$}{legendre-polynomial}{Legendre polynomial}
\aliasindexentry{pantograph-type equation}{delay-equation}{Delay and dilation equation}
\aliasindexentry{parity argument}{parity}{Parity, evenness, and oddness}
\aliasindexentry{parity-power summand}{parity-power-series}{Parity-power form}
\aliasindexentry{partial integration}{integration-by-parts}{Integration by parts}
\aliasindexentry{Pascal coefficient}{binomial-coefficient}{Binomial coefficient}
\aliasindexentry{period}{periodic-function}{Periodic function}
\aliasindexentry{periodization--sampling identity}{poisson-summation}{Poisson summation formula}
\aliasindexentry{permutation-cycle coefficients}{signed-stirling}{Signed Stirling numbers of the first kind}
\aliasindexentry{piecewise-polynomial approximant}{polynomial-step-approximant}{Polynomial step approximant}
\aliasindexentry{Poincar\'e expansion}{asymptotic-expansion}{Asymptotic expansion}
\aliasindexentry{point mass}{dirac-measure}{Dirac measure}
\aliasindexentry{pointwise finite sum}{local-finiteness}{Local finiteness}
\aliasindexentry{polynomial approximation}{approximation-theory}{Approximation theory}
\aliasindexentry{primitive-recursive function}{primitive-recursion}{Primitive recursion}
\aliasindexentry{principal part}{laurent-series}{Laurent series and pole}
\aliasindexentry{product logarithm}{lambert-w}{Lambert $W$ function}
\aliasindexentry{product probability space}{product-measure-independence}{Product measure and independence}
\aliasindexentry{Prouhet identity}{prouhet-thue-morse}{Prouhet--Thue--Morse cancellation}
\aliasindexentry{Prouhet--Thue--Morse sequence}{thue-morse}{Thue--Morse sequence}
\aliasindexentry{pushforward probability measure}{probability-law}{Probability law}
\aliasindexentry{$q$-binomial coefficient}{gaussian-binomial}{Gaussian binomial coefficient}
\aliasindexentry{$\mathbb Q(i)$}{gaussian-rational}{Gaussian rational}
\aliasindexentry{quantitative infinite-order vanishing}{effective-flatness}{Effective flatness}
\aliasindexentry{quantum calculus}{q-calculus}{$q$-calculus}
\aliasindexentry{quantum integer}{q-integer-factorial}{$q$-integer and $q$-factorial}
\aliasindexentry{$\mathbb R_{\ge0}\cup\{\infty\}$}{extended-nonnegative-reals}{Extended nonnegative reals}
\aliasindexentry{$R_n$}{reshetnikov-numbers}{Reshetnikov numbers}
\aliasindexentry{Radon--Nikodym density}{density}{Probability density}
\aliasindexentry{Radon probability measure}{borel-radon-measure}{Borel and Radon measure}
\aliasindexentry{random series}{weighted-sum}{Infinite weighted sum of random variables}
\aliasindexentry{rapidly decreasing function}{rapid-decrease}{Rapid decrease and Schwartz space}
\aliasindexentry{rational complex number}{gaussian-rational}{Gaussian rational}
\aliasindexentry{raw moment}{moment}{Moment of a distribution}
\aliasindexentry{real analytic}{analytic-function}{Analytic function and analytic locus}
\aliasindexentry{reciprocal-power-of-two value}{inverse-dyadic-value}{Inverse-dyadic value}
\aliasindexentry{rectangular density}{uniform-distribution}{Uniform distribution}
\aliasindexentry{recursive modulus of continuity}{effective-uniform-continuity}{Effective uniform continuity}
\aliasindexentry{recursive sequence equation}{recurrence-relation}{Recurrence relation}
\aliasindexentry{refinement basis}{wavelet}{Wavelet and scaling function}
\aliasindexentry{regular summability matrix}{toeplitz-summation}{Toeplitz summation theorem}
\aliasindexentry{regularized value}{finite-part}{Finite part of a singular expression}
\aliasindexentry{Rvachev atomic function}{atomic-function}{Atomic function}
\aliasindexentry{Rvachev function}{rvachev-up-function}{Rvachev's up-function}
\aliasindexentry{\texttt{rvachevUp}}{rvachev-up-function}{Rvachev's up-function}
\aliasindexentry{$\zeta(s)$}{riemann-zeta}{Riemann zeta function}
\aliasindexentry{saddle jet}{derivative-jet}{Derivative jet}
\aliasindexentry{saddle-point method}{laplace-method}{Laplace's method}
\aliasindexentry{scaling}{dilation}{Dilation and translation}
\aliasindexentry{Schwartz distribution}{distribution-generalized}{Distribution in the sense of generalized functions}
\aliasindexentry{Schwartz function}{rapid-decrease}{Rapid decrease and Schwartz space}
\aliasindexentry{second central moment}{expectation-variance}{Expectation and variance}
\aliasindexentry{self-adjoint second-order eigenvalue problem}{sturm-liouville}{Sturm--Liouville equation}
\aliasindexentry{self-convolution}{convolution}{Convolution}
\aliasindexentry{sequence transform}{generating-function}{Generating function}
\aliasindexentry{sequential computability}{computable-real-function}{Computable real function}
\aliasindexentry{set indicator}{indicator-function}{Indicator function}
\aliasindexentry{shift invariance}{translation-invariance}{Translation invariance}
\aliasindexentry{shifted $q$-factorial}{q-pochhammer}{$q$-Pochhammer symbol}
\aliasindexentry{$\operatorname{sinc}$}{sinc-function}{Sinc function}
\aliasindexentry{single-peaked function}{unimodality}{Unimodality}
\aliasindexentry{small-argument asymptotic}{endpoint-asymptotic}{Endpoint asymptotic}
\aliasindexentry{smooth compactly supported function}{bump-function}{Bump function}
\aliasindexentry{smoothing kernel}{approximate-identity}{Approximate identity}
\aliasindexentry{standard normal law}{standard-gaussian}{Standard Gaussian distribution}
\aliasindexentry{stationary point of an exponent}{saddle-point}{Saddle point}
\aliasindexentry{steepest descent}{laplace-method}{Laplace's method}
\aliasindexentry{steepest-descent point}{saddle-point}{Saddle point}
\aliasindexentry{Stirling approximation}{stirling-asymptotic}{Stirling's asymptotic formula}
\aliasindexentry{Strang--Fix-type identity}{polynomial-reproduction}{Polynomial reproduction}
\aliasindexentry{strict contraction}{contraction}{Contraction}
\aliasindexentry{strict unimodality}{unimodality}{Unimodality}
\aliasindexentry{strictly concave}{convexity-concavity}{Convexity, concavity, and inflection point}
\aliasindexentry{strictly convex}{convexity-concavity}{Convexity, concavity, and inflection point}
\aliasindexentry{strictly increasing}{monotonicity}{Monotonicity}
\aliasindexentry{strongly measurable}{measurable-function}{Measurable function}
\aliasindexentry{sum--integral expansion}{euler-maclaurin}{Euler--Maclaurin summation formula}
\aliasindexentry{sum of binary digits}{binary-weight}{Binary weight}
\aliasindexentry{sum of independent uniforms}{irwin-hall-distribution}{Irwin--Hall distribution}
\aliasindexentry{summability witness}{has-sum-tsum}{\texorpdfstring{\texttt{HasSum}}{HasSum} and infinite sum}
\aliasindexentry{summable series}{summability}{Summability of a family}
\aliasindexentry{summatory sequence}{prefix-sum}{Prefix sum and iterated prefix sum}
\aliasindexentry{support}{compact-support}{Compact support and topological support}
\aliasindexentry{tail}{partial-sum}{Partial sum and remainder}
\aliasindexentry{tail of a saddle contour}{major-minor-arc}{Major arc and minor arc}
\aliasindexentry{Taylor polynomial}{taylor-theorem}{Taylor's theorem}
\aliasindexentry{Taylor series about the origin}{power-series}{Power series and radius of convergence}
\aliasindexentry{tendsto}{filter-convergence}{Filter convergence}
\aliasindexentry{terminating recursion}{well-founded-recursion}{Well-founded recursion}
\aliasindexentry{test function}{rapid-decrease}{Rapid decrease and Schwartz space}
\aliasindexentry{Toeplitz averaging}{toeplitz-summation}{Toeplitz summation theorem}
\aliasindexentry{top-coefficient extractor}{leading-coefficient-functional}{Leading-coefficient functional}
\aliasindexentry{topological sum}{has-sum-tsum}{\texorpdfstring{\texttt{HasSum}}{HasSum} and infinite sum}
\aliasindexentry{total variation}{bounded-variation}{Bounded variation}
\aliasindexentry{totalized inverse}{bijection-inverse}{Bijection and inverse function}
\aliasindexentry{translate partition}{partition-of-unity}{Partition of unity}
\aliasindexentry{triangular recurrence}{recurrence-relation}{Recurrence relation}
\aliasindexentry{trigonometric expansion}{fourier-series}{Fourier series}
\aliasindexentry{truncation}{partial-sum}{Partial sum and remainder}
\aliasindexentry{\texttt{tsum}}{has-sum-tsum}{\texorpdfstring{\texttt{HasSum}}{HasSum} and infinite sum}
\aliasindexentry{\texttt{tsupport}}{compact-support}{Compact support and topological support}
\aliasindexentry{two-scale relation}{refinement-equation}{Refinement equation}
\aliasindexentry{two-sided Laplace product}{negative-laplace-product}{Negative-Laplace product}
\aliasindexentry{two-sided Laplace transform}{bilateral-laplace-transform}{Bilateral Laplace transform}
\aliasindexentry{unconditional convergence}{summability}{Summability of a family}
\aliasindexentry{uniform-sum spline}{irwin-hall-distribution}{Irwin--Hall distribution}
\aliasindexentry{$v_2$}{two-adic-valuation}{$2$-adic valuation}
\aliasindexentry{$v_p(n!)$}{legendre-factorial-formula}{Legendre's formula for factorial valuations}
\aliasindexentry{vague dual convergence}{weak-star-convergence}{Weak-* convergence}
\aliasindexentry{vanishing block}{zero-run}{Zero run}
\aliasindexentry{vanishing of lower degrees}{triangular-cancellation}{Triangular cancellation}
\aliasindexentry{vanishing to infinite order}{flat-function}{Flat function}
\aliasindexentry{vertical line of integration}{vertical-contour}{Vertical Bromwich contour}
\aliasindexentry{$W_0$}{lambert-w}{Lambert $W$ function}
\aliasindexentry{$W_{-1}$}{lambert-w}{Lambert $W$ function}
\aliasindexentry{weak convergence of laws}{weak-convergence}{Weak convergence of probability measures}
\aliasindexentry{weak-star convergence of measures}{weak-star-convergence}{Weak-* convergence}
\aliasindexentry{weighted composition}{ordered-composition}{Ordered composition of an integer}
\aliasindexentry{well-definedness on equivalence classes}{representation-independence}{Representation independence}
\aliasindexentry{Wick contraction}{gaussian-contraction}{Gaussian contraction}
\aliasindexentry{Wick moment}{gaussian-moment}{Gaussian moment}
\aliasindexentry{Wolfram \texttt{DiscreteLimit}}{discrete-limit}{Discrete limit}
\aliasindexentry{$\operatorname{wt}_2(n)$}{binary-weight}{Binary weight}
\aliasindexentry{$[z^n]f(z)$}{coefficient-extraction}{Coefficient extraction}
\aliasindexentry{zero inner product}{orthogonality}{Orthogonality}
\aliasindexentry{zero-mean function}{centering}{Centering and mean-zero normalization}
\aliasindexentry{zero-mean normalization}{mean-periodic}{Period mean and centered periodic function}
\aliasindexentry{zeta-type series}{dirichlet-series}{Dirichlet series}
\end{multicols}

\clearpage
\section*{Corpus note}
\addcontentsline{toc}{section}{Corpus note}
\markboth{Corpus note}{Corpus note}
The inventory was assembled recursively from the mathematical prose in the
project README and documentation audits; the two archived Arias de Reyna TeX
sources; the human-readable Fabius/up synthesis; the small-argument asymptotic
article; the Lean walkthrough; draft articles and submission notes; and the
module comments and docstrings throughout \texttt{Lean/FabiusFunction} and its
umbrella import.  Definitions are written to stand on their own, while the
second half of most entries records the term's exact role or normalization in
the repository.

\begin{boxedremark}
This is a glossary, not a replacement for the theorem statements.  When a
boundary condition, domain, coercion, normalization, or convergence mode is
material, the Lean declaration remains the authoritative specification.
\end{boxedremark}

\vfill
\begin{center}
\textcolor{rulegray}{\rule{0.72\textwidth}{0.6pt}}\\[0.7em]
{\small Generated for \texttt{Analysis/FabiusFunction}; repository snapshot
inspected 25 August 2026.}
\end{center}

\end{document}
